\documentclass[a4paper,12pt]{report}

% ─── Codifica e lingua ───────────────────────────────────────────────────────
\usepackage[utf8]{inputenc}
\usepackage[T1]{fontenc}
\usepackage{lmodern}
\usepackage[italian]{babel}
\usepackage{float}

% ─── Geometria e layout ──────────────────────────────────────────────────────
\usepackage{geometry}
\geometry{margin=2.5cm}
\usepackage{graphicx}
\usepackage{multicol}
\usepackage{booktabs}
\usepackage{longtable}
\usepackage{array}
\usepackage{tabularx}
\usepackage{multirow}

% ─── Colori ──────────────────────────────────────────────────────────────────
\usepackage{xcolor}

% Palette codice Java (come J-Pou)
\definecolor{javared}{rgb}{0.6,0,0}
\definecolor{javagreen}{rgb}{0.25,0.5,0.35}
\definecolor{javapurple}{rgb}{0.5,0,0.35}
\definecolor{javadocblue}{rgb}{0.25,0.35,0.75}

% Palette codice TypeScript / YAML
\definecolor{tsblue}{rgb}{0.0,0.2,0.6}
\definecolor{tsorange}{rgb}{0.8,0.4,0.0}
\definecolor{tsgray}{rgb}{0.5,0.5,0.5}

% Palette ambienti sicurezza
\definecolor{vulnred}{RGB}{220,50,47}
\definecolor{fixgreen}{RGB}{42,161,152}
\definecolor{warnyellow}{RGB}{181,137,0}
\definecolor{infoblue}{RGB}{38,139,210}
\definecolor{lightred}{RGB}{255,235,235}
\definecolor{lightgreen}{RGB}{235,255,235}
\definecolor{lightyellow}{RGB}{255,252,220}

% ─── Listings ────────────────────────────────────────────────────────────────
\usepackage{listings}
\usepackage{listingsutf8}

% Stile Java (identico a J-Pou)
\lstdefinestyle{java}{
  language=Java,
  basicstyle=\ttfamily\footnotesize,
  keywordstyle=\color{javapurple}\bfseries,
  stringstyle=\color{javared},
  commentstyle=\color{javagreen},
  morecomment=[s][\color{javadocblue}]{/**}{*/},
  tabsize=4,
  showspaces=false,
  showstringspaces=false,
  frame=single,
  breaklines=true,
  breakatwhitespace=false,
  columns=flexible,
  postbreak=\mbox{\textcolor{red}{$\hookrightarrow$}\space},
}

% Stile TypeScript / JavaScript
\lstdefinestyle{typescript}{
  language=JavaScript,
  basicstyle=\ttfamily\footnotesize,
  keywordstyle=\color{tsblue}\bfseries,
  stringstyle=\color{javared},
  commentstyle=\color{javagreen},
  tabsize=2,
  showspaces=false,
  showstringspaces=false,
  frame=single,
  breaklines=true,
  breakatwhitespace=false,
  columns=flexible,
  postbreak=\mbox{\textcolor{red}{$\hookrightarrow$}\space},
}

% Stile YAML / Properties
\lstdefinestyle{yaml}{
  basicstyle=\ttfamily\footnotesize,
  keywordstyle=\color{tsblue}\bfseries,
  stringstyle=\color{javared},
  commentstyle=\color{javagreen},
  tabsize=2,
  showspaces=false,
  showstringspaces=false,
  frame=single,
  breaklines=true,
  breakatwhitespace=false,
  columns=flexible,
  postbreak=\mbox{\textcolor{red}{$\hookrightarrow$}\space},
}

% Stile SQL
\lstdefinestyle{sql}{
  language=SQL,
  basicstyle=\ttfamily\footnotesize,
  keywordstyle=\color{tsblue}\bfseries,
  stringstyle=\color{javared},
  commentstyle=\color{javagreen},
  tabsize=2,
  showspaces=false,
  showstringspaces=false,
  frame=single,
  breaklines=true,
  breakatwhitespace=false,
  columns=flexible,
  postbreak=\mbox{\textcolor{red}{$\hookrightarrow$}\space},
}

% Stile Bash / Shell
\lstdefinestyle{bash}{
  language=bash,
  basicstyle=\ttfamily\footnotesize,
  keywordstyle=\color{tsblue}\bfseries,
  stringstyle=\color{javared},
  commentstyle=\color{javagreen},
  tabsize=2,
  showspaces=false,
  showstringspaces=false,
  frame=single,
  breaklines=true,
  breakatwhitespace=false,
  columns=flexible,
  postbreak=\mbox{\textcolor{red}{$\hookrightarrow$}\space},
}

% Definizione linguaggio JavaScript/TypeScript (non incluso in tutte le versioni di listings)
\lstdefinelanguage{JavaScript}{
  keywords={break,case,catch,continue,debugger,default,delete,do,else,finally,
    for,function,if,in,instanceof,new,return,switch,this,throw,try,typeof,
    var,void,while,with,let,const,class,export,import,extends,super,async,
    await,from,of,static,yield,null,true,false,undefined},
  sensitive=true,
  comment=[l]{//},
  morecomment=[s]{/*}{*/},
  morestring=[b]',
  morestring=[b]",
  morestring=[b]`
}

% Stile predefinito + UTF-8 + mapping caratteri non-ASCII nei listing
\lstset{
  style=java,
  inputencoding=utf8/latin1,
  extendedchars=true,
  literate={—}{{---}}1
           {→}{{$\to$}}1
           {←}{{$\leftarrow$}}1
           {↵}{{$\hookleftarrow$}}1
           {à}{{\`a}}1 {è}{{\`e}}1 {é}{{\'e}}1
           {ì}{{\`i}}1 {ò}{{\`o}}1 {ù}{{\`u}}1
           {À}{{\`A}}1 {È}{{\`E}}1
           {ó}{{\'o}}1 {ú}{{\'u}}1
           {─}{{\textemdash}}1
}

% ─── Ambienti personalizzati ─────────────────────────────────────────────────
\usepackage{tcolorbox}
\tcbuselibrary{skins, breakable, listings}

% Box codice INSICURO (sfondo rosso chiaro)
\newtcolorbox{insecurebox}[1][]{
  colback=lightred,
  colframe=vulnred,
  fonttitle=\bfseries\small,
  title={\textcolor{vulnred}{Codice Vulnerabile} #1},
  breakable,
  #1
}

% Box codice SICURO (sfondo verde chiaro)
\newtcolorbox{securebox}[1][]{
  colback=lightgreen,
  colframe=fixgreen,
  fonttitle=\bfseries\small,
  title={\textcolor{fixgreen}{Codice Sicuro} #1},
  breakable,
  #1
}

% Box nota/avviso
\newtcolorbox{notebox}[1][]{
  colback=lightyellow,
  colframe=warnyellow,
  fonttitle=\bfseries\small,
  title={\textcolor{warnyellow}{Nota} #1},
  breakable,
}

% ─── Overflow prevention ─────────────────────────────────────────────────────
\usepackage[protrusion=true,expansion=false]{microtype}
\setlength{\emergencystretch}{3em}
\usepackage[hyphens]{url}
\usepackage{seqsplit}

% ─── Layout e tipografia avanzata ────────────────────────────────────────────
\usepackage{csquotes}
\usepackage{enumitem}
\usepackage{titlesec}
\usepackage{fancyhdr}
\usepackage{caption}

% ─── Hyperref ────────────────────────────────────────────────────────────────
\usepackage{hyperref}
\usepackage[all]{hypcap}
\usepackage[italian]{cleveref}

\hypersetup{
    colorlinks=true,
    linkcolor=black,
    filecolor=magenta,
    urlcolor=blue,
    pdftitle={Progetto Secure Coding — Hotel PMS},
    pdfauthor={Diego Andruccioli},
}

% ─── Comandi personalizzati ──────────────────────────────────────────────────

% Intestazione vulnerabilità: \vuln{Nome Vuln}{OWASP A0X}{Threat ID}
\newcommand{\vuln}[3]{%
  \subsubsection{#1 \hfill \textcolor{vulnred}{\small [#2 — #3]}}%
}

% Badge OWASP colorato
\newcommand{\owasp}[1]{\textcolor{white}{\colorbox{vulnred}{\textbf{\small #1}}}}

% Riferimento commit Git
\newcommand{\commitref}[2]{\texttt{\small \href{https://github.com/diegoandruccioli/hotel-pms/commit/#1}{#2}}}

% Threat ID inline
\newcommand{\threat}[1]{\texttt{\small\textcolor{vulnred}{[#1]}}}

% ─── Fix: \subtitle e \subsubtitle non definiti dalla classe report ──────────
\makeatletter
\newcommand{\subtitle}[1]{\gdef\@subtitle{#1}}
\newcommand{\subsubtitle}[1]{\gdef\@subsubtitle{#1}}
\gdef\@subtitle{}
\gdef\@subsubtitle{}
\renewcommand{\maketitle}{%
  \begin{titlepage}
    \centering
    \vspace*{3.5cm}
    {\LARGE\bfseries\@title}\par\bigskip
    \ifx\@subtitle\empty\else{\large\@subtitle}\par\medskip\fi
    \ifx\@subsubtitle\empty\else{\normalsize\@subsubtitle}\par\fi
    \vspace{2cm}
    {\large\@author\par}
    \vspace{0.5cm}
    {\normalsize\@date\par}
    \vfill
  \end{titlepage}
  \setcounter{page}{1}
}
\makeatother

% ─── Metadati documento ──────────────────────────────────────────────────────
\title{\textbf{Progetto di Secure Coding}}
\subtitle{Analisi delle vulnerabilità e hardening di \\ Hotel PMS}
\subsubtitle{Property Management System per hotel}
\author{Diego Andruccioli}
\date{A.A. 2025/26}

% ─────────────────────────────────────────────────────────────────────────────
\begin{document}

\maketitle

% ─── Sommario ────────────────────────────────────────────────────────────────
\chapter*{Sommario}

Questa relazione documenta il processo di analisi della sicurezza e di \textit{hardening} applicato
a \textbf{Hotel PMS}, un sistema gestionale per la gestione alberghiera sviluppato come progetto
personale e utilizzato come base per questo elaborato del corso di \textit{Laboratorio di Sicurezza dei Sistemi e Privacy}.

\par\medskip

Il sistema è un'applicazione enterprise a microservizi composta da 9 servizi backend (Java 21,
Spring Boot 3.5, Spring Cloud 2025.0) e un frontend React 19, accessibili tramite un API Gateway centrale.
La superficie di attacco è ampia e significativa: il sistema gestisce credenziali utente, sessioni
JWT, dati personali di ospiti (PII), prenotazioni, fatture e report obbligatori per la Polizia di Stato.

\par\medskip

Il documento segue l'approccio richiesto dalla traccia d'esame: partendo dalla baseline del codice
originale (branch \texttt{pre-secure-coding}), identifica le vulnerabilità tramite analisi statica e threat
modeling STRIDE, le mappa rispetto all'OWASP Top 10 2025, e descrive per ciascuna l'intervento
di mitigazione applicato sul branch \texttt{main}, mostrando i diff
rilevanti del codice prima e dopo l'hardening.

\par\medskip

Ogni intervento è accompagnato dal riferimento al commit Git corrispondente e da un frammento
della tabella di tracciamento del \texttt{THREAT\_MODEL.md}, che evolve iterativamente
con lo sviluppo.

\tableofcontents

% ═════════════════════════════════════════════════════════════════════════════
\chapter{Introduzione}
% ═════════════════════════════════════════════════════════════════════════════

\section{Descrizione del dominio applicativo}

\textbf{Hotel PMS} (Property Management System) è un gestionale enterprise per strutture
ricettive che supporta l'intero ciclo operativo di un hotel: dalla registrazione degli ospiti
alla gestione delle prenotazioni, dal check-in/check-out alla fatturazione, fino al servizio
di Food \& Beverage del ristorante interno.

Il sistema è progettato in ottica multi-hotel: un singolo deployment gestisce più strutture,
garantendo l'isolamento dei dati tramite il campo \texttt{hotel\_id} propagato su tutti i servizi.

\section{Architettura del sistema}

Il backend è composto da 9 microservizi indipendenti, ciascuno con il proprio database PostgreSQL.
Tutte le richieste client transitano attraverso l'\textbf{API Gateway} (:8080), che valida i token JWT
e inietta gli header di autenticazione (\texttt{X-Auth-User}, \texttt{X-Auth-Role},
\texttt{X-Internal-Signature}) prima di instradare verso il servizio competente.

\begin{table}[H]
\centering
\begin{tabularx}{\textwidth}{l l X}
\toprule
\textbf{Servizio} & \textbf{Porta} & \textbf{Funzione} \\
\midrule
api-gateway & 8080 & Routing, JWT validation, rate limiting, CORS \\
config-service & 8888 & Configurazione centralizzata (segreti) \\
auth-service & 8087 & Autenticazione, JWT, gestione utenti \\
guest-service & 8083 & Anagrafica ospiti (PII) \\
inventory-service & 8081 & Gestione camere e disponibilità \\
reservation-service & 8082 & Prenotazioni \\
stay-service & 8084 & Check-in/check-out, report Alloggiati \\
billing-service & 8085 & Fatture e pagamenti \\
fb-service & 8086 & POS Food \& Beverage \\
\bottomrule
\end{tabularx}
\caption{Microservizi del sistema Hotel PMS}
\label{tab:services}
\end{table}

\section{Requisiti funzionali rilevanti per la sicurezza}

L'applicazione soddisfa tutti i requisiti minimi obbligatori indicati dalla traccia d'esame:

\begin{enumerate}
  \item \textbf{Sistema di login e registrazione} — \texttt{auth-service}: form di autenticazione
        con JWT in cookie httpOnly, gestione ruoli (ADMIN, OWNER, RECEPTIONIST).
  \item \textbf{Interazione persistente con database} — PostgreSQL per ogni microservizio,
        ORM JPA/Hibernate con query parametrizzate tramite Spring Data JPA.
  \item \textbf{Gestione dei ruoli (RBAC)} — tre ruoli distinti con endpoint amministrativi
        riservati ad ADMIN, verificati lato server tramite Spring Security.
  \item \textbf{Inserimento e visualizzazione dati utente} — campi testo libero su ospiti,
        note prenotazione, descrizioni ordini F\&B.
  \item \textbf{Risorse protette per proprietà} — ogni prenotazione/fattura/stay è legata a
        un ospite e un hotel; accesso verificato server-side tramite \texttt{hotel\_id}.
  \item \textbf{Gestione sessioni/stato} — JWT in cookie httpOnly con flag \texttt{Secure}
        e \texttt{SameSite=Strict}.
\end{enumerate}

\textbf{Funzionalità avanzate (bonus)} presenti:
\begin{itemize}
  \item Integrazione API esterna: report Alloggiati via Web Service Polizia di Stato
        (gestione sicura delle credenziali PS).
  \item Servizi esterni mockati: config-service come secret manager centralizzato.
\end{itemize}

\section{Metodologia di analisi}

L'analisi delle vulnerabilità segue la metodologia \textbf{STRIDE} di Microsoft, con mappatura
rispetto all'\textbf{OWASP Top 10 2025}. Le vulnerabilità identificate sono documentate nel file
\texttt{THREAT\_MODEL.md} del repository e tracciate nella tabella di mitigazione che evolve
con ogni commit di hardening.

Il branch \texttt{pre-secure-coding} rappresenta la \textbf{baseline} (codice originale), mentre il branch
\texttt{main} contiene tutti gli interventi di messa in sicurezza.

% ═════════════════════════════════════════════════════════════════════════════
\chapter{Analisi della Baseline Insicura}
% ═════════════════════════════════════════════════════════════════════════════

Questo capitolo documenta le vulnerabilità rilevate nel codice del branch \texttt{pre\-secure\-coding}
(baseline), organizzate per area tematica. Per ciascuna categoria viene presentata
un'analisi del pattern insicuro con i relativi snippet. I dettagli completi di ogni
intervento di mitigazione --- con codice pre/post e riferimento al commit ---
sono nel Capitolo~3 (Implementazione del Secure Coding).

\section{Panoramica delle vulnerabilità identificate}

\begin{table}[H]
\centering
\begin{tabularx}{\textwidth}{l l X l}
\toprule
\textbf{ID} & \textbf{OWASP} & \textbf{Descrizione} & \textbf{Impatto} \\
\midrule
T-AUTH-01 & A07:2025 & User Enumeration al login & ALTO \\
T-AUTH-02 & A07:2025 & Brute Force senza protezione & CRITICO \\
T-AUTH-03 & A04:2025 & Gestione password (hashing) & CRITICO \\
T-AUTH-04 & A07:2025 & Refresh Token senza revoca & ALTO \\
T-GW-01   & A01:2025 & Header X-Auth iniettabili & CRITICO \\
T-GW-03   & A02:2025 & Assenza security headers HTTP & MEDIO \\
T-GW-04   & A02:2025 & CORS misconfiguration: allowedHeaders wildcard & ALTO \\
T-GW-05   & A01:2025 & Assenza protezione CSRF & ALTO \\
T-GST-01  & A01:2025 & IDOR su ospiti (hotel\_id) & CRITICO \\
T-GST-03  & A01:2025 & Multi-tenant leak PII ospiti & CRITICO \\
T-RES-01  & A06:2025 & Double booking (no date lock) & CRITICO \\
T-FB-02   & A06:2025 & Prezzi non ricalcolati server-side & CRITICO \\
T-CFG-01  & A02:2025 & Actuator esposti & CRITICO \\
T-CFG-02  & A04:2025 & Segreti in chiaro & CRITICO \\
T-FE-01   & A05:2025 & Assenza guard statico XSS (dangerouslySetInnerHTML) & ALTO \\
T-FE-03   & A01:2025 & Route guard solo lato client & MEDIO \\
T-FE-04   & A02:2025 & Assenza CSP & ALTO \\
\bottomrule
\end{tabularx}
\caption{Riepilogo vulnerabilità identificate nella baseline}
\label{tab:vulns}
\end{table}

% ─────────────────────────────────────────────────────────────────────────────
\section{Autenticazione e Gestione Password}
% ─────────────────────────────────────────────────────────────────────────────

L'analisi della baseline ha rilevato tre vulnerabilità nell'area dell'autenticazione,
classificabili OWASP A04:2025 e A07:2025, per un impatto complessivo da \textbf{alto} a
\textbf{critico}.

\paragraph{T-AUTH-02 — Assenza di protezione al Brute Force}
Il metodo \texttt{login()} in \\ \texttt{AuthServiceImpl} non manteneva alcun contatore
di tentativi falliti né implementava alcun meccanismo di lockout applicativo.
Un attaccante poteva eseguire un attacco a dizionario sull'endpoint \texttt{POST /api/v1/auth/login}
senza limitazioni di sorta: il rate limiting del gateway (5~req/s per IP) costituiva l'unica
difesa, facilmente aggirata distribuendo le richieste su più indirizzi IP (botnet, Tor).

\begin{insecurebox}[title={AuthServiceImpl.java --- assenza di account lockout (baseline)}]
\begin{lstlisting}[style=java]
@Transactional(readOnly = true)
public AuthResponse login(LoginRequest request) {
    UserAccount user = userRepository.findByUsername(request.username())
        .orElseThrow(() -> new BadCredentialsException("INVALID_CREDENTIALS"));

    if (!passwordEncoder.matches(request.password(), user.getPasswordHash())) {
        throw new BadCredentialsException("INVALID_CREDENTIALS");
        // Nessun contatore, nessun lockout: tentativi illimitati
    }
    return buildAuthResponse(user);
}
\end{lstlisting}
\end{insecurebox}

\paragraph{T-AUTH-03 — Hashing password con cost factor insufficiente}
La baseline istanziava \texttt{BCryptPasswordEncoder} senza argomenti, usando il valore
di default \textbf{cost~10} (circa 100\,ms su hardware del~2013). Su GPU moderne questo
valore permette di verificare centinaia di migliaia di hash al secondo, rendendo un attacco
offline praticabile entro ore nel caso di furto del database.

Un secondo problema correlato: anche alzando il parametro in seguito, gli hash esistenti
sarebbero rimasti a \texttt{cost~10} per tutti gli utenti che non avessero effettuato il
logout e un nuovo login --- senza un meccanismo di rehash progressivo, la protezione era
retroattiva solo per i nuovi account.

\begin{insecurebox}[title={SecurityConfig.java --- cost factor 10 implicito (baseline)}]
\begin{lstlisting}[style=java]
@Bean
public PasswordEncoder passwordEncoder() {
    return new BCryptPasswordEncoder(); // strength=10 implicito, sotto soglia OWASP
}
\end{lstlisting}
\end{insecurebox}

\paragraph{T-AUTH-04 — Assenza di revoca del token (Refresh Token)}
La baseline emetteva un unico access token JWT con TTL configurabile ma senza alcun
meccanismo di revoca anticipata. Non esisteva il concetto di \emph{refresh token}: una volta
emesso, il JWT rimaneva valido fino alla scadenza naturale anche dopo il logout esplicito,
dopo una compromissione rilevata o in seguito a un cambio di credenziali.
L'unica azione di ``logout'' era cancellare il cookie lato client, ma il token continuava
a essere accettato dal gateway per tutta la sua durata residua.
Classificato OWASP A07 --- Authentication Failures.

% ─────────────────────────────────────────────────────────────────────────────
\section{Validazione degli Input}
% ─────────────────────────────────────────────────────────────────────────────

La baseline presentava quattro aree di mancata validazione server-side, classificabili
OWASP A05:2025 e A06:2025, che aprivano la via a injection, manipolazione dei prezzi e
corruzione dei dati di prenotazione.

\paragraph{T-GST-02 — Assenza di vincoli di formato sui dati ospite}
I DTO di richiesta del \texttt{guest-service} (\texttt{GuestRequest},
\texttt{IdentityDocumentRequestDTO}) erano protetti solo con \texttt{@NotBlank} e
\texttt{@Size}, ma privi di vincoli di formato sui campi testo libero.
Un attaccante poteva inviare tag HTML o payload SQL nei campi \texttt{firstName},
\texttt{phone}, \texttt{address}, oppure una data di nascita nel futuro.

\begin{insecurebox}[title={GuestRequest.java --- nessun \texttt{@Pattern} (baseline)}]
\begin{lstlisting}[style=java]
public record GuestRequest(
    @NotBlank @Size(max = 100) String firstName,  // nessun @Pattern
    @NotBlank @Size(max = 100) String lastName,   // nessun @Pattern
    @NotBlank @Email @Size(max = 150) String email,
    @Size(max = 20) String phone,                 // nessun @Pattern
    @Size(max = 255) String address,              // nessun @Pattern
    @Size(max = 50) String city,                  // nessun @Pattern
    @Size(max = 50) String country,               // nessun @Pattern
    LocalDate dateOfBirth) {                      // nessun @Past
}
// Payload valido inviabile: firstName = "<script>alert(1)</script>"
\end{lstlisting}
\end{insecurebox}

\paragraph{T-GST-04 — Mass Assignment nel mapper ospite}
\texttt{GuestMapper.toEntity()} non dichiarava esplicitamente i campi da ignorare durante
il mapping da DTO a entità JPA. In assenza di \texttt{@Mapping(ignore=true)} su
\texttt{id}, \texttt{hotelId}, \texttt{active}, \texttt{createdAt}, \texttt{updatedAt},
MapStruct avrebbe tentato di copiare qualsiasi campo omonimo presente nel DTO, aprendo
la strada a una sovrascrittura dei metadati di sistema da parte del client (mass assignment).

\paragraph{T-RES-03 — Mancanza di validazione cross-field sulle date di prenotazione}
\texttt{ReservationRequest} non verificava che \texttt{checkOutDate} fosse successiva a
\texttt{checkInDate}. Un client poteva inviare una prenotazione con date invertite o con
\texttt{checkIn == checkOut} (zero-night stay), producendo una query di overlap detection
semanticamente errata che poteva ignorare conflitti reali.

\paragraph{T-FB-02 — Prezzi degli ordini F\&B controllati dal client}
\texttt{OrderItemRequest} accettava i campi \texttt{unitPrice} e \texttt{itemName}
direttamente dalla richiesta HTTP. Il servizio si fidava del prezzo dichiarato dal client
anziché consultare il catalogo server-side, consentendo a un utente di creare ordini con
importi arbitrari (incluso zero o negativi) e di alterare il nome degli articoli.
Classificato OWASP A06:2025 --- Insecure Design.

\paragraph{T-BILL-02 — Assenza di validazione degli importi pagamento}
\texttt{PaymentRequest} non vincolava il campo \texttt{amount} con \texttt{@Digits} né
\texttt{@Positive}, e \texttt{PaymentServiceImpl} non normalizzava l'importo prima di
aggiungere il pagamento alla fattura. Importi negativi o con più di due decimali erano
accettati, corrompendo il saldo e la logica di chiusura fattura.

% ─────────────────────────────────────────────────────────────────────────────
\section{Autorizzazione e Controllo degli Accessi}
% ─────────────────────────────────────────────────────────────────────────────

Questa categoria rappresenta l'area di rischio più estesa della baseline: sette
vulnerabilità IDOR distribuite su cinque microservizi, più una rottura strutturale
del meccanismo di isolamento multi-tenant a livello di gateway.
Tutte classificate OWASP A01 --- Broken Access Control.

\paragraph{T-GST-01 / T-GST-03 / T-RES-02 / T-BILL-01 / T-FB-01 — IDOR su cinque servizi}
I repository JPA di \texttt{guest-service}, \texttt{reservation-service},
\texttt{billing-service} e \texttt{fb-service} operavano senza scoping per
\texttt{hotel\_id}. Le query \texttt{findById()} e \texttt{findAll()} restituivano
risorse di qualsiasi hotel, permettendo a un operatore dell'Hotel~A, conoscendo
o indovinando un UUID, di leggere e modificare ospiti, prenotazioni, fatture e
ordini F\&B dell'Hotel~B.

\begin{insecurebox}[title={Pattern IDOR --- GuestServiceImpl.java (baseline)}]
\begin{lstlisting}[style=java]
// INSICURO: lookup per solo UUID, nessun check ownership hotel
private Guest resolveGuest(UUID id) {
    return guestRepository.findById(id)   // trova chiunque, qualsiasi hotel
        .orElseThrow(() -> new NotFoundException("GUEST_NOT_FOUND"));
}

// INSICURO: lista tutti gli ospiti di tutti gli hotel
public Page<GuestResponse> getAllGuests(Pageable pageable) {
    return guestRepository.findAll(pageable)  // nessun filtro hotel_id
        .map(guestMapper::toResponse);
}
\end{lstlisting}
\end{insecurebox}

Lo stesso pattern era replicato in \texttt{ReservationServiceImpl},
\texttt{InvoiceServiceImpl} e \texttt{RestaurantOrderServiceImpl}: nessuno dei quattro
servizi aveva \texttt{hotel\_id} nella propria entità JPA né nelle query del repository.

\paragraph{T-GW-06 — Propagazione X-Auth-Hotel mancante end-to-end}
Anche con i fix IDOR applicati ai singoli servizi, l'isolamento multi-tenant avrebbe
comunque fallito in produzione per una ragione strutturale più profonda: l'entità
\texttt{UserAccount} in \texttt{auth-service} non possedeva il campo \texttt{hotel\_id},
quindi \texttt{JwtService.generateToken()} non includeva il claim \texttt{hotelId}
nel JWT emesso. Di conseguenza \texttt{AuthenticationFilter} nell'API Gateway non
poteva iniettare l'header \texttt{X-Auth-Hotel} --- perché il dato non era disponibile
nel token --- e i servizi interni ricevevano sempre l'header vuoto, restituendo~401
a ogni richiesta autenticata via gateway.

\begin{insecurebox}[title={JwtService.java (auth-service) --- claim hotelId assente (baseline)}]
\begin{lstlisting}[style=java]
// INSICURO: il JWT non contiene il claim hotelId
public String generateToken(UserAccount user) {
    return Jwts.builder()
        .setSubject(user.getUsername())
        .claim("role", user.getRole().name())
        // hotelId mai aggiunto: l'isolamento multi-tenant non funziona
        .setIssuedAt(new Date())
        .setExpiration(new Date(System.currentTimeMillis() + expiration))
        .signWith(key, SignatureAlgorithm.HS256)
        .compact();
}
\end{lstlisting}
\end{insecurebox}

\paragraph{T-FE-03 — Route guard solo lato client}
Il frontend React proteggeva le route amministrative (\texttt{/owner-dashboard})
nascondendo i link nella navigazione in base al ruolo memorizzato nello store Zustand.
Tuttavia, \texttt{ProtectedRoute} non verificava il ruolo: un utente con ruolo
\texttt{RECEPTIONIST} poteva accedere direttamente digitando l'URL nel browser,
ricevendo dati o form riservati ad \texttt{ADMIN}/\texttt{OWNER}.
Classificato OWASP A01 --- Broken Access Control (frontend enforcement only).

\paragraph{T-RES-01 — Double Booking (assenza di lock concorrente)}\mbox{}\\
\texttt{\seqsplit{ReservationServiceImpl.createReservation()}} eseguiva un controllo di overlap
sulle date di prenotazione, ma senza \emph{pessimistic locking} sulla query.
In condizioni di concorrenza (due richieste simultanee per la stessa camera e le
stesse date) entrambe le transazioni potevano superare il controllo e inserire due
prenotazioni in conflitto, violando il requisito di unicità operativa.
Classificato OWASP A06:2025 --- Insecure Design.

% ─────────────────────────────────────────────────────────────────────────────
\section{Gestione delle Sessioni}
% ─────────────────────────────────────────────────────────────────────────────

La baseline gestiva le sessioni tramite JWT in cookie \texttt{httpOnly} --- scelta
corretta per prevenire l'accesso JavaScript al token --- ma presentava due lacune
rilevanti nella gestione del ciclo di vita del token e nella protezione dalle
richieste cross-site.

\paragraph{T-AUTH-04 — Refresh token senza rotazione né blacklist Redis}
Il sistema non implementava il meccanismo di \emph{refresh token}: il solo access token
JWT, una volta emesso, era accettato dal gateway per tutta la sua durata (configurabile,
default 15~min) anche dopo il logout. Non esisteva alcun identificatore univoco
(\texttt{jti}) per tracciare il token, né una blacklist su Redis per invalidarlo
anticipatamente. Di conseguenza:

\begin{itemize}
  \item Un logout non revocava effettivamente la sessione: il token rubato prima del
        logout restava valido.
  \item Un cambio di password o ruolo non invalidava i token emessi in precedenza.
  \item Uno scenario di \emph{token theft} non aveva rimedio applicativo immediato.
\end{itemize}

\paragraph{T-GW-05 — Assenza di protezione CSRF}
L'API Gateway non implementava alcun meccanismo esplicito di protezione CSRF.
Il flag \texttt{SameSite=Strict} sul cookie JWT offre una prima difesa, ma non è
sufficiente: browser datati, proxy aziendali o subdomain takeover possono violarne
il rispetto. In assenza di un token CSRF verificato server-side, un sito malevolo
poteva indurre il browser dell'utente autenticato a inviare richieste mutanti
(\texttt{POST}, \texttt{PUT}, \texttt{DELETE}) verso il sistema, sfruttando
l'invio automatico del cookie da parte del browser.

\begin{insecurebox}[title={api-gateway (baseline) --- nessun filtro CSRF}]
\begin{lstlisting}[style=java]
// INSICURO: il gateway accetta qualsiasi richiesta con cookie JWT valido,
// indipendentemente dall'origine della richiesta.
// Non esiste alcun GlobalFilter che validi un CSRF token.
// Un sito malevolo puo' scatenare:
//   POST https://hotel.example.com/api/v1/reservations
//   Cookie: jwt=<token valido> (inviato automaticamente dal browser)
// La richiesta viene eseguita senza che il server verifichi l'origine.
\end{lstlisting}
\end{insecurebox}

% ─────────────────────────────────────────────────────────────────────────────
\section{Configurazioni di Sicurezza}
% ─────────────────────────────────────────────────────────────────────────────

La baseline presentava sei categorie di misconfiguration, alcune con impatto critico
(segreti in chiaro, actuator esposti), altre con impatto medio-alto (rate limiting
parziale, assenza di header HTTP, CORS permissivo, mancanza di CSP).
Classificate prevalentemente OWASP A02:2025 --- Security Misconfiguration.

\paragraph{T-CFG-01 — Endpoint \texttt{/actuator} esposti sulla porta applicativa}
Tutti e 9 i microservizi pubblicavano gli endpoint Spring Boot Actuator
(\texttt{/actuator/health}, \texttt{/actuator/prometheus}, \texttt{/actuator/info})
sulla stessa porta dell'API (\texttt{8080}--\texttt{8087}, \texttt{8888}),
a sua volta esposta nel \texttt{docker-compose.yml}.
\texttt{/actuator/prometheus} rivelava metriche interne del processo JVM (thread pool,
connection pool, circuit breaker state) preziose per un attaccante in fase di
\emph{reconnaissance}. \texttt{/actuator/health} con \texttt{show-details: always}
esponeva lo stato dei sottosistemi e i dettagli di connessione a database e Redis.

\paragraph{T-CFG-02 — Segreti in chiaro nei file versionati}
Due categorie di segreti erano presenti in chiaro nel repository:
\begin{enumerate}
  \item Il fallback hardcoded nella proprietà \texttt{jwt.secret} in
        \texttt{auth-service.yml} e in \texttt{@Value} di \texttt{api-gateway/JwtUtil.java}:
        entrambi contenevano la stringa Base64 nota \\
        \texttt{\seqsplit{bXktMzItYnl0ZS1zZWNyZXQta2V5LWZvci10ZXN0LWxvY2FsLWRldi0xMjM0NQ==}}.
        Se la variabile d'ambiente \texttt{JWT\_SECRET} non veniva impostata, il servizio
        avviava silenziosamente con questo segreto debole, permettendo a chiunque
        conosca il repository di forgiare JWT validi.
  \item \texttt{POSTGRES\_PASSWORD} hardcoded a \texttt{password} nel
        \texttt{docker-compose.yml}: password identica per tutti gli ambienti, senza
        possibilità di rotazione senza modifica al codice sorgente.
\end{enumerate}

\begin{insecurebox}[title={auth-service.yml + JwtUtil.java --- fallback hardcoded (baseline)}]
\begin{lstlisting}[style=yaml]
# auth-service.yml: fallback noto, committato nel repository
jwt:
  secret: "${JWT_SECRET:bXktMzItYnl0ZS1zZWNyZXQta2V5LWZvci10ZXN0LWxvY2FsLWRldi0xMjM0NQ==}"
\end{lstlisting}
\begin{lstlisting}[style=java]
// JwtUtil.java (api-gateway): stesso fallback anche nel gateway
@Value("${jwt.secret:bXktMzItYnl0ZS1zZWNyZXQta2V5LWZvci10ZXN0LWxvY2FsLWRldi0xMjM0NQ==}")
private String secret;
\end{lstlisting}
\end{insecurebox}

\paragraph{T-CFG-03 — Config Server senza autenticazione}
Il \texttt{config-service} (Spring Cloud Config) distribuiva la configurazione centralizzata
--- inclusi segreti, credenziali e chiavi --- senza alcuna protezione HTTP Basic.
Qualsiasi processo raggiungibile sulla porta \texttt{8888} della rete Docker poteva
ottenere la configurazione completa di ogni microservizio con una semplice richiesta GET.

\paragraph{T-GW-02 — Rate limiting applicato solo alla route \texttt{/auth}}
Nella baseline il \\ \texttt{RequestRateLimiter} era configurato esclusivamente sulla
route di autenticazione (\texttt{/api/v1/auth/**}). Le sei route di business
(\texttt{/guests}, \texttt{/reservations}, \texttt{/stays}, \texttt{/invoices},
\texttt{/fb}, \texttt{/rooms}) non avevano alcun limite di frequenza,
esponendo i servizi interni ad attacchi di enumerazione e a scenari di denial-of-service
applicativo. Inoltre, il \texttt{KeyResolver} leggeva l'indirizzo TCP remoto
(\texttt{RemoteAddress}), risultando inutile per deployment dietro proxy/load balancer
che presentano un unico IP al gateway.

\paragraph{T-GW-03 / T-GW-04 — Assenza di security headers HTTP e CORS permissivo}
Nessuno dei microservizi né il gateway aggiungeva gli header di sicurezza HTTP fondamentali:
\texttt{Strict-Transport-Security}, \texttt{X-Content-Type-Options},
\texttt{X-Frame-Options}, \texttt{Referrer-Policy} e \texttt{Permissions-Policy}
erano assenti da tutte le risposte. La configurazione CORS usava
\texttt{allowedHeaders: "*"} (wildcard), trasmettendo qualsiasi header client al backend
senza whitelist esplicita.

\paragraph{T-FE-04 — Assenza di Content Security Policy}
La SPA React servita da Nginx non definiva alcuna \texttt{Content-Security-Policy}.
In assenza di CSP, un XSS riuscito avrebbe potuto caricare script arbitrari da domini
esterni, estrarre dati e contattare server C\&C senza restrizioni. Ugualmente assente
era qualsiasi meccanismo di prevenzione statica del \texttt{dangerouslySetInnerHTML}
lato frontend.

\paragraph{T-AUTH-05 / T-STAY-02 / T-BILL-03 — Assenza di audit logging}
Nella baseline nessun servizio emetteva log strutturati per gli eventi sensibili:
tentativi di login (riusciti e falliti), check-in/check-out e transazioni di pagamento
avvenivano senza traccia applicativa. In uno scenario di incidente o frode, l'assenza
di audit trail avrebbe reso impossibile ricostruire la sequenza degli eventi.
Classificato OWASP A09:2025 --- Security Logging \& Alerting Failures.

% ═════════════════════════════════════════════════════════════════════════════
\chapter{Implementazione del Secure Coding}
\label{ch:storico}
% ═════════════════════════════════════════════════════════════════════════════

% ─────────────────────────────────────────────────────────────────────────────
\section{Autenticazione Sicura}
% ─────────────────────────────────────────────────────────────────────────────

% TEMPLATE PER OGNI VULNERABILITÀ CORRETTA:
%
% \vuln{Nome Vulnerabilità}{OWASP AXX}{T-ID}
%
% \paragraph{Descrizione}
% Testo della vulnerabilità originale.
%
% \begin{insecurebox}[title={Codice Vulnerabile — NomeFile.java}]
% \begin{lstlisting}[style=java]
% // codice insicuro
% \end{lstlisting}
% \end{insecurebox}
%
% \paragraph{Mitigazione applicata}
% Testo della mitigazione.
%
% \begin{securebox}[title={Codice Sicuro — NomeFile.java}]
% \begin{lstlisting}[style=java]
% // codice sicuro
% \end{lstlisting}
% \end{securebox}
%
% \paragraph{Riferimento} Commit: \commitref{SHA}{messaggio commit}

% ── Nota sui threat già mitigati nella baseline ──────────────────────────────
\begin{notebox}[title={Mitigazioni presenti nella baseline (branch \texttt{pre-secure-coding})}]
Le vulnerabilità \textbf{T-AUTH-01}, \textbf{T-GW-01} e \textbf{T-FE-02} risultano già
mitigate nel codice del branch \texttt{pre-secure-coding} (baseline). Tali protezioni sono state introdotte
consapevolmente durante la fase di sviluppo iniziale, non come hardening successivo.
Le sezioni seguenti documentano il pattern insicuro che è stato evitato, a dimostrazione
della consapevolezza del rischio durante la progettazione.
\end{notebox}

\subsection{User Enumeration — Risposta Uniforme al Login \threat{T-AUTH-01}}

\paragraph{Descrizione}
Un endpoint di login che restituisce messaggi di errore differenziati (es.\
\texttt{"USER\_NOT\_FOUND"} vs \texttt{"INVALID\_PASSWORD"}) consente a un attaccante
di enumerare gli username validi tramite una scansione automatizzata.
Classificato OWASP A07 --- Authentication Failures.

\begin{insecurebox}[title={Pattern Vulnerabile Evitato --- AuthServiceImpl.java}]
\begin{lstlisting}[style=java]
// INSICURO: messaggi diversi rivelano l'esistenza dell'account
public AuthResponse login(LoginRequest request) {
    UserAccount user = userRepository.findByUsername(request.username())
        .orElseThrow(
            () -> new BadCredentialsException("USER_NOT_FOUND")); // leak!

    if (!passwordEncoder.matches(request.password(), user.getPasswordHash())) {
        throw new BadCredentialsException("INVALID_PASSWORD"); // diverso -> enum
    }
    ...
}
\end{lstlisting}
\end{insecurebox}

\paragraph{Mitigazione applicata (baseline)}
\texttt{AuthServiceImpl.login()} utilizza lo stesso messaggio \texttt{"INVALID\_CREDENTIALS"}
sia quando l'utente non esiste sia quando la password è errata, rendendo impossibile
distinguere i due casi dall'esterno.

\begin{securebox}[title={Codice Sicuro --- AuthServiceImpl.java (branch \texttt{main})}]
\begin{lstlisting}[style=java]
// SICURO: stesso messaggio in entrambi i casi
public AuthResponse login(LoginRequest request) {
    UserAccount user = userRepository.findByUsername(request.username())
        .orElseThrow(
            () -> new BadCredentialsException("INVALID_CREDENTIALS"));

    if (!passwordEncoder.matches(request.password(), user.getPasswordHash())) {
        throw new BadCredentialsException("INVALID_CREDENTIALS"); // identico
    }
    ...
}
\end{lstlisting}
\end{securebox}

\paragraph{Riferimento} Presente nella baseline --- \texttt{auth-service/AuthServiceImpl.java}.

% ─────────────────────────────────────────────────────────────────────────────
\subsection{Iniezione Header X-Auth via Client \threat{T-GW-01}}

\paragraph{Descrizione}
In un'architettura a microservizi, i servizi interni ricevono dal gateway gli header
\texttt{X-Auth-User} e \texttt{X-Auth-Role} per identificare il chiamante senza
ri-validare il JWT. Un servizio che accetti tali header senza verificarne l'origine
è vulnerabile a impersonation: un client che raggiunga direttamente la porta interna
del servizio può iniettare \texttt{X-Auth-Role: ADMIN} e ottenere privilegi massimi.
Classificato OWASP A01 --- Broken Access Control.

\begin{insecurebox}[title={Pattern Vulnerabile Evitato --- filtro servizi interni}]
\begin{lstlisting}[style=java]
// INSICURO: il servizio trust gli header senza verifica di origine
String username = request.getHeader("X-Auth-User");
String role    = request.getHeader("X-Auth-Role");
// Chiunque raggiunga direttamente il servizio puo' iniettare ROLE_ADMIN
SecurityContextHolder.getContext().setAuthentication(
    new UsernamePasswordAuthenticationToken(
        username, "", List.of(new SimpleGrantedAuthority("ROLE_" + role))));
\end{lstlisting}
\end{insecurebox}

\paragraph{Mitigazione applicata (baseline)}
Il gateway firma ogni richiesta con \\ \texttt{HMAC-SHA256(secret, "username:role")} e
inietta la firma in \texttt{X-Internal-Signature}. L'\texttt{InternalAuthFilter}
presente in ogni microservizio ricalcola la firma e la confronta con
\texttt{MessageDigest.isEqual} (confronto constant-time, immune a timing attack)
prima di popolare il \texttt{SecurityContext}. Senza il segreto HMAC condiviso è
impossibile forgiare una firma valida.

\begin{securebox}[title={Codice Sicuro --- InternalAuthFilter.java (tutti i servizi)}]
\begin{lstlisting}[style=java]
// SICURO: verifica HMAC-SHA256 prima di fidarsi degli header
private boolean isSignatureValid(String username, String role, String sig) {
    String expected = computeHmac(username, role);
    // Confronto constant-time: previene timing attack
    return MessageDigest.isEqual(
        expected.getBytes(StandardCharsets.UTF_8),
        sig.getBytes(StandardCharsets.UTF_8));
}

// In doFilterInternal():
if (!isSignatureValid(username, role, signature)) {
    rejectRequest(response, "Invalid internal request signature"); // 401
    return;
}
\end{lstlisting}
\end{securebox}

\paragraph{Riferimento} Presente nella baseline --- \path{guest-service/.../InternalAuthFilter.java}
(replicato in tutti i microservizi).

% ─────────────────────────────────────────────────────────────────────────────
\subsection{Token JWT in \texttt{localStorage} --- Esposizione a XSS \threat{T-FE-02}}

\paragraph{Descrizione}
Memorizzare il token di sessione in \texttt{localStorage} lo rende leggibile da qualsiasi
script JavaScript eseguito nella pagina, incluso codice iniettato via XSS. Un attaccante
può estrarlo con \texttt{localStorage.getItem('jwt')} ed eseguire richieste autenticate
su server arbitrari. Classificato OWASP A04:2025 --- Cryptographic Failures / A07.

\begin{insecurebox}[title={Pattern Vulnerabile Evitato --- gestione token frontend}]
\begin{lstlisting}[style=typescript]
// INSICURO: token in localStorage, leggibile da XSS
const response = await axios.post('/api/v1/auth/login', credentials);
localStorage.setItem('jwt', response.data.token);  // esposto a JS
axios.defaults.headers.common['Authorization'] =
    `Bearer ${localStorage.getItem('jwt')}`;
\end{lstlisting}
\end{insecurebox}

\paragraph{Mitigazione applicata (baseline)}
Il server imposta il JWT in un cookie \texttt{httpOnly; Secure; SameSite=Strict}.
Il flag \texttt{httpOnly} lo rende inaccessibile a JavaScript; \texttt{Secure}
ne impedisce la trasmissione in chiaro; \texttt{SameSite=Strict} blocca l'invio
cross-site. Il frontend React non gestisce mai il valore del token.

\begin{securebox}[title={Codice Sicuro --- AuthController.java (branch \texttt{main})}]
\begin{lstlisting}[style=java]
// SICURO: cookie httpOnly, invisibile a JavaScript
private ResponseCookie createCookie(String value, int maxAge) {
    return ResponseCookie.from("jwt", value)
        .httpOnly(true)       // inaccessibile a JavaScript
        .secure(true)         // solo HTTPS
        .sameSite("Strict")   // blocco invio cross-site
        .path("/")
        .maxAge(maxAge)
        .build();
}
// Il frontend non tocca mai il token: il browser lo gestisce in automatico
\end{lstlisting}
\end{securebox}

\paragraph{Riferimento} Presente nella baseline --- \texttt{auth-service/AuthController.java}.

% ─────────────────────────────────────────────────────────────────────────────
\subsection{Brute Force sul Login --- Assenza di Account Lockout \threat{T-AUTH-02}}

\paragraph{Descrizione}
L'assenza di un meccanismo di account lockout consente attacchi a dizionario o brute
force sull'endpoint di login senza limitazioni applicative. Il rate limiting del gateway
(5 req/s per IP) è insufficiente contro attacchi distribuiti da IP multipli. Classificato
OWASP A07 --- Authentication Failures.

La baseline non implementava alcun contatore di tentativi falliti:

\begin{insecurebox}[title={Codice Vulnerabile --- AuthServiceImpl.java (baseline, branch \texttt{pre-secure-coding})}]
\begin{lstlisting}[style=java]
@Transactional(readOnly = true)
public AuthResponse login(LoginRequest request) {
    UserAccount user = userRepository.findByUsername(request.username())
        .orElseThrow(() -> new BadCredentialsException("INVALID_CREDENTIALS"));

    if (!passwordEncoder.matches(request.password(), user.getPasswordHash())) {
        throw new BadCredentialsException("INVALID_CREDENTIALS");
        // Nessun contatore, nessun lockout:
        // un attaccante puo' provare password illimitatamente
    }

    return new AuthResponse(
        jwtService.generateToken(user.getUsername(), user.getRole()));
}
\end{lstlisting}
\end{insecurebox}

\paragraph{Mitigazione applicata}
Il metodo \texttt{login()} è stato esteso con una politica di lockout progressivo:
\begin{itemize}
  \item Dopo \textbf{5 tentativi falliti consecutivi} il campo \texttt{locked\_until}
        viene impostato a \texttt{now + 15 minuti}.
  \item Se l'account è bloccato, viene lanciata \texttt{AccountLockedException}
        (HTTP 429 Too Many Requests) \emph{prima} del controllo password,
        evitando timing leak sulla validità della password.
  \item Un login riuscito resetta il contatore e cancella il lock.
  \item I valori (soglia e durata) sono costanti named, senza magic numbers.
\end{itemize}

La Flyway migration \texttt{V2\_\_add\_login\_lockout.sql} aggiunge le colonne
\\ \texttt{failed\_attempts} e \texttt{locked\_until} alla tabella \texttt{user\_account}.

\begin{securebox}[title={Codice Sicuro --- AuthServiceImpl.java (main)}]
\begin{lstlisting}[style=java]
private static final int MAX_FAILED_ATTEMPTS = 5;
private static final Duration LOCKOUT_DURATION = Duration.ofMinutes(15);

@Transactional
public AuthResponse login(LoginRequest request) {
    UserAccount user = userRepository.findByUsername(request.username())
        .orElseThrow(() -> new BadCredentialsException("INVALID_CREDENTIALS"));

    // 1. Verifica lockout PRIMA del controllo password (no timing leak)
    if (user.getLockedUntil() != null
            && Instant.now().isBefore(user.getLockedUntil())) {
        throw new AccountLockedException("ACCOUNT_TEMPORARILY_LOCKED"); // 429
    }

    // 2. Verifica password
    if (!passwordEncoder.matches(request.password(), user.getPasswordHash())) {
        int newAttempts = user.getFailedAttempts() + 1;
        user.setFailedAttempts(newAttempts);
        if (newAttempts >= MAX_FAILED_ATTEMPTS) {
            user.setLockedUntil(Instant.now().plus(LOCKOUT_DURATION));
        }
        userRepository.save(user);
        throw new BadCredentialsException("INVALID_CREDENTIALS");
    }

    // 3. Login riuscito: reset contatore e lock
    if (user.getFailedAttempts() > 0) {
        user.setFailedAttempts(0);
        user.setLockedUntil(null);
        userRepository.save(user);
    }
    return new AuthResponse(
        jwtService.generateToken(user.getUsername(), user.getRole()));
}
\end{lstlisting}
\end{securebox}

\begin{securebox}[title={V2\_\_add\_login\_lockout.sql --- Flyway migration}]
\begin{lstlisting}[style=yaml]
ALTER TABLE user_account
    ADD COLUMN IF NOT EXISTS failed_attempts INTEGER     NOT NULL DEFAULT 0,
    ADD COLUMN IF NOT EXISTS locked_until   TIMESTAMPTZ;
\end{lstlisting}
\end{securebox}

\paragraph{Test} Aggiunti 4 test in \texttt{AuthServiceImplTest}:
\texttt{loginThrowsWhenAccountIsLocked},
\texttt{loginIncrementsFailedAttemptsOnWrongPassword}, \\
\texttt{loginLocksAccountAfterMaxFailedAttempts},
\texttt{loginResetsCounterOnSuccess}.

\paragraph{Riferimento} Commit: \commitref{a90e462}{security(auth): add account lockout for brute-force \\ protection (T-AUTH-02)}

\subsection{Sicurezza Hashing Password \threat{T-AUTH-03}}

\paragraph{Descrizione}
L'algoritmo BCrypt adotta un \textit{cost factor} (o \textit{strength}) che determina il numero
di iterazioni dell'algoritmo: un valore troppo basso rende il calcolo dell'hash banalmente veloce
su hardware moderno, abbattendo il tempo necessario per un attacco brute-force su hash rubati.
La configurazione originale usava \texttt{new BCryptPasswordEncoder()} senza argomenti,
che corrisponde al valore di default \textbf{10} (circa 100\,ms su hardware del 2013).
L'OWASP Authentication Cheat Sheet raccomanda un cost factor minimo di \textbf{10} e
preferibilmente \textbf{12} per hardware contemporaneo (\(\approx\)400\,ms per hash).
Classificato OWASP A04:2025 --- Cryptographic Failures.

Un secondo problema correlato è l'\textbf{assenza di rehash progressivo}: anche dopo
aver alzato il cost factor, gli hash esistenti nel database rimangono con il vecchio valore
finché non vengono aggiornati esplicitamente. Senza un meccanismo di aggiornamento automatico,
la protezione è retroattiva solo per i nuovi utenti.

\begin{insecurebox}[title={Codice Vulnerabile --- SecurityConfig.java (cost factor 10 implicito)}]
\begin{lstlisting}[style=java]
// INSICURO: BCryptPasswordEncoder() senza argomento usa strength=10,
// valore di default del 2013, insufficiente su hardware moderno.
@Bean
public PasswordEncoder passwordEncoder() {
    return new BCryptPasswordEncoder(); // strength=10 implicito
}
\end{lstlisting}
\end{insecurebox}

\begin{insecurebox}[title={Codice Vulnerabile --- AuthServiceImpl.java (nessun rehash progressivo)}]
\begin{lstlisting}[style=java]
// INSICURO: anche alzando il cost factor, gli hash esistenti
// restano a strength=10 finche' l'utente non cambia password.
// Nessun aggiornamento automatico al login.
if (user.getFailedAttempts() > 0) {
    user.setFailedAttempts(0);
    user.setLockedUntil(null);
    userRepository.save(user);
}
// Hash con cost=10 rimane in DB invariato
\end{lstlisting}
\end{insecurebox}

\paragraph{Mitigazione applicata}
Sono stati applicati due interventi distinti e complementari:

\begin{enumerate}
  \item \textbf{Upgrade cost factor a 12} in \texttt{SecurityConfig}: tutti i nuovi hash
        (registrazione + rehash) vengono calcolati con strength 12, che corrisponde a circa
        400\,ms su hardware contemporaneo, allineandosi alla raccomandazione OWASP.
  \item \textbf{Lazy rehash on login}: ad ogni login riuscito, il metodo \\
        \texttt{passwordEncoder.upgradeEncoding(storedHash)} verifica se l'hash corrente
        è stato calcolato con un cost factor inferiore a quello configurato.
        Se sì, la password in chiaro (disponibile solo in questo momento) viene ricalcolata
        e persistita, migrando silenziosamente l'intero parco utenti senza downtime
        né password reset forzati.
\end{enumerate}

L'evento di rehash è tracciato con log strutturato \texttt{[AUTH] PASSWORD\_REHASHED}
per audit e monitoraggio della progressione della migrazione.

\begin{securebox}[title={Codice Sicuro --- SecurityConfig.java (cost factor 12)}]
\begin{lstlisting}[style=java]
// SICURO: cost factor 12 (OWASP-raccomandato per hardware moderno,
// ~400 ms per hash; valore di default era 10 = ~100 ms).
@Bean
public PasswordEncoder passwordEncoder() {
    return new BCryptPasswordEncoder(12);
}
\end{lstlisting}
\end{securebox}

\begin{securebox}[title={Codice Sicuro --- AuthServiceImpl.java (lazy rehash on login)}]
\begin{lstlisting}[style=java]
// SICURO: dopo un login riuscito, verifica se l'hash in DB usa
// un cost factor inferiore a quello corrente (12). Se necessario,
// ricalcola e persiste l'hash aggiornato (lazy migration).
final boolean needsRehash =
    passwordEncoder.upgradeEncoding(user.getPasswordHash());
final boolean needsReset = user.getFailedAttempts() > 0;
if (needsRehash) {
    user.setPasswordHash(passwordEncoder.encode(request.password()));
    log.info("[AUTH] PASSWORD_REHASHED | user={}", user.getUsername());
}
if (needsReset) {
    user.setFailedAttempts(0);
    user.setLockedUntil(null);
}
if (needsRehash || needsReset) {
    userRepository.save(user); // unica save anche se entrambi true
}
\end{lstlisting}
\end{securebox}

\paragraph{Test} Aggiunti 2 test in \texttt{AuthServiceImplTest}: \\
\texttt{loginRehashesPasswordWhenEncodingNeedsUpgrade} (verifica che il nuovo hash venga
persistito quando \texttt{upgradeEncoding} restituisce \texttt{true}) e\\
\texttt{\seqsplit{loginSkipsRehashWhenEncodingIsUpToDate}} (verifica che nessuna \texttt{save}
venga eseguita quando l'hash è già aggiornato e i tentativi falliti sono zero).

\paragraph{Riferimento SEI CERT} La mitigazione applica la regola \textbf{MSC62-J} del SEI CERT Java Coding Standard: \emph{``Store passwords using a hash function specifically designed for the purpose''}. BCrypt con cost~12 è la scelta corretta per il 2025/26 secondo OWASP; MD5 e SHA-1 (funzioni general-purpose) non sono accettabili.

\paragraph{Riferimento} Commit: \commitref{4e1f1e9}{fix(security): T-AUTH-03 — BCrypt cost 12 + lazy rehash on login on auth-service}

% ─── Upgrade T-AUTH-03: Argon2id + password complexity (A04:2025 + A07:2021) ───────────────
\paragraph{Upgrade T-AUTH-03 — Argon2id e Password Complexity Policy (esercizi OWASP)}

\subparagraph{Argon2id — KDF memory-hard}
L'analisi degli esercizi OWASP Top10 ha evidenziato un ulteriore gap: BCrypt, pur con
cost~12, è un algoritmo \textit{memory-weak} — il calcolo di un singolo hash usa al
massimo pochi kB di RAM, rendendo gli attacchi con GPU/ASIC efficienti.
\textbf{Argon2id} (vincitore del Password Hashing Competition 2015, raccomandato da
NIST SP 800-63B) è \textit{memory-hard}: il parametro \texttt{memory} forza l'attaccante
a allocare RAM proporzionale, abbattendo drasticamente il throughput su hardware
parallelo. Classificato OWASP A04:2025 --- Cryptographic Failures.

\begin{enumerate}
  \item \textbf{DelegatingPasswordEncoder} (default: \texttt{argon2}): nuove password
        codificate come \texttt{\{argon2\}\$argon2id\$v=19\$m=19456,t=2,p=1\$...\$...}
        Parametri: memoria 19\,MiB, 2 iterazioni, parallelismo~1, coerenti con
        l'esercizio A04:2025.
  \item \textbf{Legacy BCrypt}: hash già presenti nel DB mantengono il prefisso
        \texttt{\{bcrypt\}\$2a\$12\$...}. Il meccanismo di lazy rehash esistente
        (\texttt{upgradeEncoding}) li migra ad Argon2id al primo login successivo,
        senza downtime né password reset forzati.
  \item \textbf{Dipendenza}: \texttt{org.bouncycastle:bcprov-jdk18on:1.84}
        (già presente nel BOM per CVE-2026-5598; aggiunta come
        \texttt{implementation} dependency).
\end{enumerate}

\begin{securebox}[title={Codice Sicuro --- SecurityConfig.java (Argon2id + DelegatingPasswordEncoder)}]
\begin{lstlisting}[style=java]
// SICURO: Argon2id memory-hard (19 MiB, 2 iter, parallelismo 1)
// DelegatingPasswordEncoder gestisce anche gli hash BCrypt legacy.
@Bean
public PasswordEncoder passwordEncoder() {
    final Map<String, PasswordEncoder> encoders = new HashMap<>();
    encoders.put("argon2", new Argon2PasswordEncoder(
            16, 32, 1, 19 * 1024, 2));   // salt=16B, hash=32B
    encoders.put("bcrypt", new BCryptPasswordEncoder(12));
    return new DelegatingPasswordEncoder("argon2", encoders);
}
\end{lstlisting}
\end{securebox}

\subparagraph{Password Complexity Policy — A07:2021}
L'esercizio A07:2021 richiedeva la validazione della password in fase di registrazione:
almeno \textbf{16 caratteri}, \textbf{2 maiuscole}, \textbf{2 cifre}, \textbf{2 caratteri speciali}.
La baseline accettava password di soli 8 caratteri senza requisiti di complessità
(\texttt{@Size(min~=~8)}), classificato OWASP A07:2021 --- Identification and Authentication Failures.

La mitigazione applica un'annotazione \texttt{@Pattern} con lookahead regex su tutti
i DTO che accettano password in input: \texttt{RegisterRequest}, \texttt{CreateUserRequest},
\texttt{ChangePasswordRequest} (solo \texttt{newPassword}).
\texttt{LoginRequest} non è soggetto alla policy (l'utente non può modificare la password
storica a retroattività).

\begin{securebox}[title={Codice Sicuro --- RegisterRequest.java (password complexity)}]
\begin{lstlisting}[style=java]
// SICURO: >=16 caratteri, 2 maiuscole, 2 cifre, 2 speciali.
// LoginRequest escluso: la policy non e' retroattiva sul login.
@NotBlank
@Pattern(
    regexp = "^(?=.*[A-Z].*[A-Z])(?=.*[0-9].*[0-9])"
           + "(?=.*[^A-Za-z0-9].*[^A-Za-z0-9]).{16,}$",
    message = "PASSWORD_TOO_WEAK")
String password
\end{lstlisting}
\end{securebox}

\paragraph{Riferimento} Commit: \commitref{446ad8d}{feat(security): T-AUTH-03 upgrade — Argon2id KDF + password complexity policy (A04+A07:2025)}

\subsection{Refresh Token Rotation e Blacklist \threat{T-AUTH-04}}

\paragraph{Descrizione}
La baseline emetteva un unico token JWT con scadenza di \textbf{24 ore} privo di un meccanismo di
revoca. Un token intercettato (es.\ via XSS o man-in-the-middle su HTTP) rimaneva valido per
l'intera finestra temporale, senza alcuna possibilità di invalidarlo lato server. L'assenza di
\textit{token rotation} aggravava il problema: un refresh token mai revocato poteva essere riutilizzato
a tempo indeterminato, anche dopo il logout esplicito dell'utente.
Classificato OWASP A07 --- Authentication Failures.

\begin{insecurebox}[title={Codice Vulnerabile --- JwtService.java + AuthController.java (baseline)}]
\begin{lstlisting}[style=java]
// Un solo token JWT con TTL 24 ore, nessun refresh token separato
@Value("${jwt.expiration}")
private long jwtExpiration; // 86_400_000 ms = 24 ore

// AuthController: un unico cookie "jwt", nessun cookie "refresh_token"
private static final int COOKIE_MAX_AGE = 86_400; // 24 ore

// Login: emette un solo token
return new AuthResponse(
    jwtService.generateToken(user.getUsername(), user.getRole()));

// Logout: azzera solo il cookie "jwt"
// Il token non viene blacklistato: se qualcuno lo ha catturato,
// rimane valido fino alla sua scadenza naturale (24 ore)
ResponseEntity.ok()
    .header(HttpHeaders.SET_COOKIE, createCookie("", 0).toString())
    .build();
\end{lstlisting}
\end{insecurebox}

\paragraph{Mitigazione applicata}
Sono stati applicati quattro interventi coordinati:

\begin{enumerate}
  \item \textbf{Separazione access/refresh token}: il login emette ora due JWT distinti:
        un \textit{access token} con TTL \textbf{15 minuti} e un \textit{refresh token} con
        TTL \textbf{7 giorni}. Il refresh token porta una claim \texttt{jti} (UUID univoco) e
        una claim \texttt{typ=refresh} che lo distingue dall'access token.
  \item \textbf{Blacklist Redis}: al momento del refresh o del logout, il JTI del token consumato
        viene scritto in Redis con chiave \texttt{rt:blacklist:<jti>} e TTL pari alla vita residua
        del token. L'entry si auto-cancella alla scadenza naturale, senza job di pulizia.
  \item \textbf{Token rotation}: \texttt{POST /api/v1/auth/refresh} valida il token ricevuto nel
        cookie \texttt{refresh\_token}, ne blacklista il JTI, e risponde con un \textbf{nuovo}
        token pair completo. Ogni refresh produce una coppia fresca --- un token intercettato
        prima della rotation risulta immediatamente blacklistato.
  \item \textbf{Frontend interceptor}: l'interceptor Axios intercetta i 401, tenta silenziosamente
        il refresh e riprova la richiesta originale. Nessun logout visibile per sessioni attive
        entro i 7 giorni; logout automatico se il refresh token è scaduto o revocato.
\end{enumerate}

\begin{securebox}[title={Codice Sicuro --- JwtService.java (generazione refresh token con jti)}]
\begin{lstlisting}[style=java]
// Refresh token: TTL 7 giorni, jti UUID, claim typ=refresh
public String generateRefreshToken(String username, Role role) {
    Map<String, Object> claims = new HashMap<>();
    claims.put("role", role.name());
    claims.put("typ", "refresh");  // distingue da access token
    return buildRefreshToken(
        claims, username, refreshExpiration,
        UUID.randomUUID().toString()); // jti univoco
}

// Metodi di supporto per la blacklist
public String extractJti(String token) {
    return extractClaim(token, Claims::getId);
}
public Instant extractExpirationInstant(String token) {
    return extractClaim(token, c -> c.getExpiration().toInstant());
}
public boolean isRefreshToken(String token) {
    return "refresh".equals(
        extractClaim(token, c -> c.get("typ", String.class)));
}
\end{lstlisting}
\end{securebox}

\begin{securebox}[title={Codice Sicuro --- RefreshTokenServiceImpl.java (blacklist Redis)}]
\begin{lstlisting}[style=java]
// Redis key: "rt:blacklist:<jti>" con TTL = vita residua del token
@Override
public void blacklist(String jti, Instant expiresAt) {
    long ttlSeconds = Duration.between(Instant.now(), expiresAt).getSeconds();
    if (ttlSeconds > 0) {
        // Scadenza automatica: nessun job di pulizia necessario
        redisTemplate.opsForValue()
            .set("rt:blacklist:" + jti, "1", Duration.ofSeconds(ttlSeconds));
    }
}

@Override
public boolean isBlacklisted(String jti) {
    return Boolean.TRUE.equals(redisTemplate.hasKey("rt:blacklist:" + jti));
}
\end{lstlisting}
\end{securebox}

\begin{securebox}[title={Codice Sicuro --- AuthServiceImpl.java (refresh con rotation)}]
\begin{lstlisting}[style=java]
@Override
@Transactional(readOnly = true)
public AuthResponse refresh(String refreshToken) {
    // 1. Verifica che sia un refresh token
    if (!jwtService.isRefreshToken(refreshToken)) {
        log.warn("[AUTH] REFRESH_REJECTED | reason=NOT_REFRESH_TOKEN");
        throw new BadCredentialsException("INVALID_REFRESH_TOKEN");
    }
    // 2. Verifica che il JTI non sia in blacklist
    String jti = jwtService.extractJti(refreshToken);
    if (jti == null || refreshTokenService.isBlacklisted(jti)) {
        log.warn("[AUTH] REFRESH_REJECTED | reason=TOKEN_BLACKLISTED");
        throw new BadCredentialsException("INVALID_REFRESH_TOKEN");
    }
    // 3. Verifica che l'utente esista ancora
    String username = jwtService.extractUsername(refreshToken);
    UserAccount user = userRepository.findByUsername(username)
        .orElseThrow(() -> new BadCredentialsException("INVALID_REFRESH_TOKEN"));
    // 4. Blacklista il token consumato (rotation)
    refreshTokenService.blacklist(jti, jwtService.extractExpirationInstant(refreshToken));
    // 5. Emette il nuovo token pair
    log.info("[AUTH] REFRESH_SUCCESS | user={}", username);
    return new AuthResponse(
        jwtService.generateToken(username, user.getRole()),
        jwtService.generateRefreshToken(username, user.getRole()));
}
\end{lstlisting}
\end{securebox}

\begin{securebox}[title={Codice Sicuro --- api.ts (Axios interceptor con silent refresh)}]
\begin{lstlisting}[style=typescript]
// Interceptor: su 401, tenta refresh silenzioso prima di fare logout
api.interceptors.response.use(
  (response) => response,
  async (error) => {
    const url = error.config?.url ?? '';
    if (error.response?.status === 401) {
      if (url.includes('/refresh')) {
        performLogout(); return Promise.reject(error); // refresh scaduto
      }
      if (!url.includes('/login') && !url.includes('/me')) {
        if (isRefreshing) {
          // Accoda la richiesta — verra' ritentata dopo il refresh
          return new Promise((resolve, reject) => {
            pendingQueue.push({ resolve, reject });
          }).then(() => api(originalConfig));
        }
        isRefreshing = true;
        try {
          await api.post('/api/v1/auth/refresh'); // cookie gestito dal browser
          drainQueue(true, null);
          return api(originalConfig);            // riprova la richiesta originale
        } catch (err) {
          drainQueue(false, err); performLogout();
        } finally { isRefreshing = false; }
      }
    }
    return Promise.reject(error);
  }
);
\end{lstlisting}
\end{securebox}

\paragraph{Test} Aggiunti 6 test in \texttt{AuthServiceImplTest}:
\texttt{refreshSuccessIssuesNewTokenPair} (verifica rotation e blacklist del JTI consumato), \\
\texttt{refreshThrowsWhenTokenIsNotRefreshType},
\texttt{refreshThrowsWhenJtiIsBlacklisted},\\
\texttt{refreshThrowsWhenUserNotFound},
\texttt{invalidateRefreshTokenBlacklistsValidToken},\\
\texttt{invalidateRefreshTokenSilentlyHandlesInvalidToken}.
Aggiunti 5 test in \\ \texttt{RefreshTokenServiceImplTest} che verificano il comportamento
della blacklist Redis (store con TTL, skip su token già scaduto, lookup true/false/null).

\paragraph{Effetto}
La finestra di esposizione di un access token intercettato è ora di \textbf{15 minuti}
invece di 24 ore. Il refresh token, con TTL 7 giorni, viene ruotato a ogni uso:
intercettarlo è inutile poiché è già blacklistato non appena viene usato dal legittimo
titolare. Il logout revoca immediatamente il refresh token, chiudendo la sessione
server-side senza dipendere dalla scadenza naturale del token.

\paragraph{Riferimento} Commit: \commitref{5ce3010}{fix(security): T-AUTH-04 --- refresh token rotation + Redis blacklist}

% ─────────────────────────────────────────────────────────────────────────────
\subsection{Cambio password non revoca le sessioni attive \threat{T-AUTH-04 residuo}}

\paragraph{Descrizione}
La rotation introdotta con \texttt{5ce3010} copre il caso in cui un token viene
intercettato durante la trasmissione, ma lasciava aperta una finestra di attacco
distinta: un utente che \textbf{cambia la propria password} non revocava i refresh
token già emessi. Un attaccante in possesso di un refresh token (es.\ estratto da
un dispositivo rubato o da un cookie leak) poteva continuare a ruotarlo per
\textbf{7 giorni}, ottenendo nuovi access token validi anche dopo che la vittima
aveva aggiornato la password. Il sistema non possedeva alcun meccanismo per
correlare un token in volo con la ``generazione'' di credenziali corrente.
Classificato OWASP A07 --- Authentication Failures.

\begin{insecurebox}[title={Codice Vulnerabile --- AuthServiceImpl.java (refresh senza controllo versione)}]
\begin{lstlisting}[style=java]
// Prima del fix: refresh accetta qualsiasi token valido non blacklistato,
// indipendentemente da quando e' stato emesso rispetto all'ultimo cambio password.
@Override
@Transactional(readOnly = true)
public AuthResponse refresh(final String refreshToken) {
    // ...
    String username = jwtService.extractUsername(refreshToken);
    UserAccount user = userRepository.findByUsername(username)
        .orElseThrow(() -> new BadCredentialsException("INVALID_REFRESH_TOKEN"));
    // MANCANZA: nessun confronto con la "generazione" corrente delle credenziali.
    // Un token emesso 6 giorni fa (prima del cambio password) supera questo check.
    refreshTokenService.blacklist(jti, jwtService.extractExpirationInstant(refreshToken));
    return new AuthResponse(
        jwtService.generateToken(username, user.getRole(), user.getHotelId()),
        jwtService.generateRefreshToken(username, user.getRole(), user.getHotelId()));
}

// JwtService.generateToken / generateRefreshToken: nessun claim tv
public String generateToken(String username, Role role, UUID hotelId) {
    Map<String, Object> claims = new HashMap<>();
    claims.put("role", role.name());
    claims.put("hotelId", hotelId.toString());
    // claim "tv" assente: il token non porta informazioni sulla generazione
    return buildToken(claims, username, jwtExpiration);
}
\end{lstlisting}
\end{insecurebox}

\paragraph{Mitigazione applicata}
L'approccio scelto è il \textbf{token versioning}: un contatore monotòno \texttt{tokenVersion}
in \texttt{user\_account} viene embedded come claim \texttt{tv} in ogni JWT (access
e refresh). Al cambio password il contatore viene incrementato e il nuovo valore
sovrascrive la chiave Redis \texttt{user:tv:<username>}. Il metodo \texttt{refresh()}
confronta il claim \texttt{tv} del token in ingresso con il valore in Redis: se
divergono, il token appartiene a una generazione precedente e viene rifiutato con
HTTP~401.

\noindent Questo approccio è superiore alla blacklist JTI per il caso del cambio
password perché:
\begin{itemize}
  \item Opera come \emph{blanket revocation}: un singolo write Redis revoca
        \emph{tutti} i token precedenti, senza dover enumerare i JTI.
  \item Non cresce con il numero di sessioni: la chiave è una sola per utente.
  \item I token access (TTL 15 min) scadono naturalmente in pochi minuti;
        la protezione immediata copre i refresh token (TTL 7 giorni).
\end{itemize}

\noindent\textbf{Migrazione graceful:} se la chiave Redis è assente
(\texttt{storedTv = -1}), il check viene saltato. Questo permette agli utenti
che avevano sessioni attive prima del deploy di continuare a usarle.
La chiave viene scritta al primo login/register successivo al deploy.

\begin{securebox}[title={Codice Sicuro --- Flyway V4 (migrazione schema)}]
\begin{lstlisting}[style=sql]
-- V4__add_token_version.sql
ALTER TABLE user_account
    ADD COLUMN IF NOT EXISTS token_version INT NOT NULL DEFAULT 0;
COMMENT ON COLUMN user_account.token_version IS
    'Monotonically incrementing counter embedded in JWT "tv" claim. '
    'Incrementing it (on password change) invalidates all previously issued tokens.';
\end{lstlisting}
\end{securebox}

\begin{securebox}[title={Codice Sicuro --- JwtService.java (claim tv nei JWT)}]
\begin{lstlisting}[style=java]
private static final String CLAIM_TOKEN_VERSION = "tv";

// generateToken e generateRefreshToken accettano ora tokenVersion come 4 argomento
public String generateToken(String username, Role role, UUID hotelId, int tokenVersion) {
    Map<String, Object> claims = new HashMap<>();
    claims.put("role", role.name());
    claims.put("hotelId", hotelId.toString());
    claims.put(CLAIM_TOKEN_VERSION, tokenVersion); // nuovo claim
    return buildToken(claims, username, jwtExpiration);
}

// Estrae il claim tv: -1 se assente (token pre-feature → check saltato)
public int extractTokenVersion(String token) {
    Integer tv = extractClaim(token, c -> c.get(CLAIM_TOKEN_VERSION, Integer.class));
    return tv != null ? tv : -1;
}
\end{lstlisting}
\end{securebox}

\begin{securebox}[title={Codice Sicuro --- RefreshTokenServiceImpl.java (Redis token version)}]
\begin{lstlisting}[style=java]
private static final String TV_PREFIX = "user:tv:";

// Scrive "user:tv:<username>" = versione corrente con TTL = lifetime refresh token
@Override
public void storeTokenVersion(String username, int version, Duration ttl) {
    redisTemplate.opsForValue()
        .set(TV_PREFIX + username, String.valueOf(version), ttl);
}

// Restituisce -1 se la chiave non esiste (migrazione graceful)
@Override
public int getTokenVersion(String username) {
    String raw = redisTemplate.opsForValue().get(TV_PREFIX + username);
    if (raw == null) return -1;
    try { return Integer.parseInt(raw); }
    catch (NumberFormatException e) { return -1; }
}
\end{lstlisting}
\end{securebox}

\begin{securebox}[title={Codice Sicuro --- AuthServiceImpl.java (check tv in refresh + changePassword)}]
\begin{lstlisting}[style=java]
// In refresh(): check versione prima di ruotare il token
final int storedTv = refreshTokenService.getTokenVersion(username);
final int jwtTv = jwtService.extractTokenVersion(refreshToken);
// storedTv >= 0: chiave presente -> check attivo
// jwtTv = -1 (claim assente) != storedTv positivo -> rifiuto corretto
if (storedTv >= 0 && jwtTv != storedTv) {
    log.warn("[AUTH] REFRESH_REJECTED | reason=TOKEN_VERSION_MISMATCH"
             + " | user={} | jwtTv={} | storedTv={}", username, jwtTv, storedTv);
    throw new BadCredentialsException("INVALID_REFRESH_TOKEN");
}

// POST /api/v1/auth/change-password
@Override
@Transactional
public AuthResponse changePassword(String username, ChangePasswordRequest request) {
    UserAccount user = userRepository.findByUsername(username)
        .orElseThrow(() -> new BadCredentialsException("INVALID_CREDENTIALS"));
    // Secondo fattore: verifica password corrente
    if (!passwordEncoder.matches(request.currentPassword(), user.getPasswordHash())) {
        log.warn("[AUTH] CHANGE_PASSWORD_REJECTED | reason=BAD_CURRENT_PASSWORD");
        throw new BadCredentialsException("INVALID_CURRENT_PASSWORD");
    }
    user.setPasswordHash(passwordEncoder.encode(request.newPassword()));
    user.setTokenVersion(user.getTokenVersion() + 1); // invalida tutte le sessioni
    userRepository.save(user);
    // Sovrascrive Redis: tutti i token con tv=vecchio verranno rifiutati in refresh()
    refreshTokenService.storeTokenVersion(username, user.getTokenVersion(), REFRESH_TOKEN_TTL);
    log.info("[AUTH] PASSWORD_CHANGED | user={} | newTokenVersion={}", username,
             user.getTokenVersion());
    // Emette un nuovo token pair per mantenere attiva la sessione corrente
    return new AuthResponse(
        jwtService.generateToken(username, user.getRole(), user.getHotelId(),
                                  user.getTokenVersion()),
        jwtService.generateRefreshToken(username, user.getRole(), user.getHotelId(),
                                         user.getTokenVersion()));
}
\end{lstlisting}
\end{securebox}

\paragraph{Test}
Aggiunti 6 test tra \texttt{AuthServiceImplTest} e \texttt{RefreshTokenServiceImplTest}:
\begin{itemize}
  \item \texttt{refreshThrowsWhenTokenVersionMismatch}: verifica che il refresh
        venga rifiutato quando il claim \texttt{tv} del JWT diverge dal valore Redis
        (simula un token emesso prima del cambio password).
  \item \texttt{refreshSkipsVersionCheckWhenRedisKeyAbsent}: verifica la migrazione
        graceful (\texttt{storedTv = -1} $\Rightarrow$ check saltato).
  \item \texttt{changePasswordSuccessIncreasesVersionAndIssuesNewTokenPair}: verifica
        l'incremento di \texttt{tokenVersion}, il write Redis e l'emissione del nuovo
        token pair.
  \item \texttt{changePasswordThrowsWhenCurrentPasswordIsWrong}: verifica che la
        password corrente errata blocchi il cambio.
  \item \texttt{changePasswordThrowsWhenUserNotFound}: verifica la gestione di utenti
        inesistenti.
  \item \texttt{storeTokenVersionWritesCorrectKeyAndValue} e \\
        \texttt{getTokenVersionReturnsStoredValue} in \texttt{RefreshTokenServiceImplTest}.
\end{itemize}

\paragraph{Effetto}
Dopo il cambio password, qualsiasi tentativo di ruotare un refresh token
``pre-cambio'' viene respinto con HTTP~401 non appena il token giunge al \\
\texttt{POST /api/v1/auth/refresh}. La finestra di esposizione residua si riduce
ai soli access token ancora in volo (massimo \textbf{15 minuti}), che scadono
naturalmente. Il cambio password reimposta anche i cookie del browser della
sessione corrente con il nuovo token pair, garantendo continuità di servizio
al legittimo proprietario.

\paragraph{Riferimento} Commit: \commitref{b473265}{fix(security): GAP-1 --- password change revokes existing sessions via tokenVersion}

% ─────────────────────────────────────────────────────────────────────────────
\section{Audit Log e Security Monitoring \textnormal{\small\textcolor{warnyellow}{[Appendice A]}}}
% ─────────────────────────────────────────────────────────────────────────────

\begin{notebox}[title={Nota metodologica --- Implementazioni AI-assisted}]
Le implementazioni di questa sezione (T-AUTH-05, T-STAY-01, T-STAY-02, T-STAY-03, T-BILL-03) 
sono state sviluppate con supporto di Claude AI. Il codice è funzionante e integrato nel sistema.
La critica critica all'output dell'AI è documentata nell'Appendice~B.
\end{notebox}

\section{Validazione, Sanificazione ed Encoding}
% ─────────────────────────────────────────────────────────────────────────────

\subsection{Input Validation su dati ospite \threat{T-GST-02}}
% ─────────────────────────────────────────────────────────────────────────────

\paragraph{Descrizione}
Il \texttt{guest-service} espone endpoint REST per la creazione e l'aggiornamento
di profili ospite (\texttt{POST /api/v1/guests}, \texttt{PUT /api/v1/guests/\{id\}}).
I DTO di richiesta --- \texttt{GuestRequest} e \texttt{IdentityDocumentRequestDTO} ---
erano protetti solo con \texttt{@NotBlank} e \texttt{@Size}, ma \textbf{privi di
vincoli di formato} sui campi testo libero. Un attaccante poteva inviare payload
contenenti tag HTML o script nei campi nome, telefono, città e numero documento,
oppure una data di nascita nel futuro. Classificato OWASP A05:2025 --- Injection.

\begin{insecurebox}[title={Codice Vulnerabile --- GuestRequest.java (baseline)}]
\begin{lstlisting}[style=java]
public record GuestRequest(
    @NotBlank @Size(max = 100) String firstName,  // nessun @Pattern
    @NotBlank @Size(max = 100) String lastName,   // nessun @Pattern
    @NotBlank @Email @Size(max = 150) String email,
    @Size(max = 20) String phone,                 // nessun @Pattern
    @Size(max = 255) String address,              // nessun @Pattern
    @Size(max = 50) String city,                  // nessun @Pattern
    @Size(max = 50) String country,               // nessun @Pattern
    LocalDate dateOfBirth) {                      // nessun @Past
}
// Payload valido inviabile: firstName = "<script>alert(1)</script>"
// oppure phone = "'; DROP TABLE guests;--"
\end{lstlisting}
\end{insecurebox}

\begin{securebox}[title={Codice Sicuro --- GuestRequest.java + ValidationConstants.java}]
\begin{lstlisting}[style=java]
// ValidationConstants.java
public static final String NAME_PATTERN     = "^[\\p{L} '\\-]+$";
public static final String PHONE_PATTERN    = "^[+]?[0-9 ()\\-]+$";
public static final String TEXT_SAFE_PATTERN = "^[^<>&\"]+$";
public static final String LOCATION_PATTERN = "^[\\p{L}0-9 '\\-.]+$";
public static final String DOCUMENT_NUMBER_PATTERN = "^[A-Za-z0-9\\-]+$";

// GuestRequest.java
public record GuestRequest(
    @NotBlank @Size(max = MAX_FIRST_NAME_LENGTH)
    @Pattern(regexp = NAME_PATTERN) String firstName,

    @NotBlank @Size(max = MAX_LAST_NAME_LENGTH)
    @Pattern(regexp = NAME_PATTERN) String lastName,

    @NotBlank @Email @Size(max = MAX_EMAIL_LENGTH) String email,

    @Size(max = MAX_PHONE_LENGTH)
    @Pattern(regexp = PHONE_PATTERN) String phone,

    @Size(max = MAX_ADDRESS_LENGTH)
    @Pattern(regexp = TEXT_SAFE_PATTERN) String address,

    @Size(max = MAX_LOCATION_LENGTH)
    @Pattern(regexp = LOCATION_PATTERN) String city,

    @Size(max = MAX_LOCATION_LENGTH)
    @Pattern(regexp = LOCATION_PATTERN) String country,

    @Past LocalDate dateOfBirth) {
}

// IdentityDocumentRequestDTO.java
public record IdentityDocumentRequestDTO(
    @NotNull DocumentType documentType,
    @NotBlank @Size(max = MAX_DOCUMENT_NUMBER_LENGTH)
    @Pattern(regexp = DOCUMENT_NUMBER_PATTERN) String documentNumber,
    @NotNull @Past LocalDate issueDate,
    @NotNull @FutureOrPresent LocalDate expiryDate,
    @Size(max = MAX_COUNTRY_LENGTH)
    @Pattern(regexp = LOCATION_PATTERN) String issuingCountry) {
}
\end{lstlisting}
\end{securebox}

\paragraph{Effetto}
Jakarta Bean Validation rifiuta automaticamente al layer controller qualsiasi
richiesta che non rispetti i pattern dichiarati: la risposta è
\texttt{400 Bad Request} prima che il dato raggiunga il service o il database.
I pattern scelti coprono nomi internazionali (Unicode \texttt{\textbackslash p\{L\}}),
numeri di telefono in formato internazionale, e bloccano esplicitamente i caratteri
usati negli attacchi XSS e SQL Injection (\texttt{<}, \texttt{>}, \texttt{\&},
\texttt{"}, \texttt{;}).

\paragraph{Riferimento SEI CERT} Regola \textbf{IDS00-J}: \emph{``Sanitize untrusted data passed across a trust boundary''}. Principio chiave: \textbf{la normalizzazione deve precedere la validazione}; i pattern allow-list bloccano esplicitamente i caratteri usati in attacchi XSS e SQL Injection (\texttt{<}, \texttt{>}, \texttt{;}, \texttt{\&}).

\paragraph{Riferimento} Commit: \commitref{fb6fdff}{fix(security): T-GST-02 + T-GST-04 --- input validation \\
patterns and explicit mass-assignment guards on guest-service}

% ─────────────────────────────────────────────────────────────────────────────
\subsection{Difesa contro Mass Assignment \threat{T-GST-04}}
% ─────────────────────────────────────────────────────────────────────────────

\paragraph{Descrizione}
Il \emph{mass assignment} è una vulnerabilità per cui un client può impostare
campi sensibili dell'entità (chiave primaria, \texttt{hotel\_id}, flag
\texttt{active}, timestamp di auditing) includendoli nel payload JSON della
richiesta. In Java con Spring e MapStruct il rischio è mitigato dal tipo record
(immutabile) e dalla policy \texttt{ReportingPolicy.IGNORE}, ma il metodo
\texttt{GuestMapper\#toEntity()} non dichiarava esplicitamente quali target
fossero esclusi dal mapping, rendendo la protezione implicita e non verificabile
a ispezione statica del codice. Classificato OWASP A06:2025 --- Insecure Design.

\begin{insecurebox}[title={Codice Vulnerabile --- GuestMapper.java (baseline)}]
\begin{lstlisting}[style=java]
@Mapper(componentModel = "spring",
        unmappedTargetPolicy = ReportingPolicy.IGNORE,  // protezione implicita
        uses = IdentityDocumentMapper.class)
public interface GuestMapper {

    // Nessun @Mapping(ignore=true): i campi id, hotelId, active,
    // createdAt, updatedAt sono esclusi solo per comportamento di default
    // di ReportingPolicy.IGNORE, non per dichiarazione esplicita.
    Guest toEntity(GuestRequest request);

    @Mapping(target = "id",                ignore = true)
    @Mapping(target = "hotelId",           ignore = true)
    @Mapping(target = "identityDocuments", ignore = true)
    @Mapping(target = "active",            ignore = true)
    @Mapping(target = "createdAt",         ignore = true)
    @Mapping(target = "updatedAt",         ignore = true)
    void updateEntityFromRequest(GuestRequest request, @MappingTarget Guest target);
}
\end{lstlisting}
\end{insecurebox}

\begin{securebox}[title={Codice Sicuro --- GuestMapper.java}]
\begin{lstlisting}[style=java]
@Mapper(componentModel = "spring",
        unmappedTargetPolicy = ReportingPolicy.IGNORE,
        uses = IdentityDocumentMapper.class)
public interface GuestMapper {

    // Ignores espliciti: nessun dato client puo' influenzare questi campi,
    // indipendentemente dalla policy del mapper o da future refactoring.
    // gdprConsentDate aggiunto successivamente nell'ambito della fix T-GST-05
    // (GDPR data retention): campo di sistema impostato dal service, mai dal client.
    @Mapping(target = "id",                ignore = true)
    @Mapping(target = "hotelId",           ignore = true)
    @Mapping(target = "identityDocuments", ignore = true)
    @Mapping(target = "active",            ignore = true)
    @Mapping(target = "createdAt",         ignore = true)
    @Mapping(target = "updatedAt",         ignore = true)
    @Mapping(target = "gdprConsentDate",   ignore = true)
    Guest toEntity(GuestRequest request);

    @Mapping(target = "id",                ignore = true)
    @Mapping(target = "hotelId",           ignore = true)
    @Mapping(target = "identityDocuments", ignore = true)
    @Mapping(target = "active",            ignore = true)
    @Mapping(target = "createdAt",         ignore = true)
    @Mapping(target = "updatedAt",         ignore = true)
    @Mapping(target = "gdprConsentDate",   ignore = true)
    void updateEntityFromRequest(GuestRequest request, @MappingTarget Guest target);
}
\end{lstlisting}
\end{securebox}

\paragraph{Effetto}
La dichiarazione esplicita degli ignore garantisce che qualunque refactoring
futuro (rimozione di \texttt{ReportingPolicy.IGNORE}, cambio di libreria di
mapping) non reintroduca accidentalmente la possibilità di impostare campi
sensibili da input client. La protezione diventa \emph{fail-safe by design}
anziché dipendente da un comportamento di default.

\paragraph{Riferimento SEI CERT} Regola \textbf{OBJ13-J}: \emph{``Ensure that references to mutable objects are not exposed''}. La separazione DTO/entità con \texttt{@Mapping(target=..., ignore=true)} dichiarativo è la concretizzazione di questo principio: nessun riferimento al modello interno è raggiungibile dall'input client.

\paragraph{Riferimento} Commit: \commitref{fb6fdff}{fix(security): T-GST-02 + T-GST-04 --- input validation \\
patterns and explicit mass-assignment guards on guest-service}

\subsection{Server-side Price Validation \threat{T-FB-02}}
% ─────────────────────────────────────────────────────────────────────────────

\paragraph{Descrizione}
Il \texttt{fb-service} gestisce gli ordini del Food \& Beverage (bar, ristorante)
addebitati al conto camera. Prima della mitigazione, il DTO \texttt{OrderItemRequest}
accettava il campo \texttt{unitPrice} inviato direttamente dal client. Il service
usava quel valore senza alcuna verifica per calcolare il \texttt{totalAmount}
dell'ordine. Un utente malintenzionato poteva manipolare il payload JSON
(e.g.\ con un proxy interceptor come Burp Suite) e impostare prezzi arbitrari,
inclusi valori prossimi allo zero, ottenendo beni e servizi a costo irrisorio
senza essere rilevato. Classificato OWASP A06:2025 --- Insecure Design.

\begin{insecurebox}[title={Codice Vulnerabile --- OrderItemRequest.java + service (baseline)}]
\begin{lstlisting}[style=java]
// INSICURO: il prezzo arriva dal client senza validazione
public record OrderItemRequest(
    @NotBlank String itemName,
    @NotNull @Positive Integer quantity,
    @NotNull @Positive BigDecimal unitPrice  // attaccante imposta 0.01
) {}

// INSICURO: totalAmount calcolato con il prezzo del client
BigDecimal totalAmount = order.getItems().stream()
    .map(item -> item.getUnitPrice()         // valore iniettato dal client
                     .multiply(BigDecimal.valueOf(item.getQuantity())))
    .reduce(BigDecimal.ZERO, BigDecimal::add);
\end{lstlisting}
\end{insecurebox}

\paragraph{Mitigazione applicata}
L'intervento introduce un catalogo di voci di menu server-side e rimuove
il prezzo dal perimetro di controllo del client:

\begin{enumerate}
    \item \textbf{Entita' MenuItem} --- nuova tabella \texttt{menu\_items}
          (Flyway V2) con campi \texttt{name} e \texttt{price}. I prezzi
          sono scritti solo da operatori autenticati con accesso DB;
          il client non puo' mai influenzarli.
    \item \textbf{Seed data} --- 10 voci tipiche di un hotel (espresso,
          cappuccino, club sandwich, bistecca\ldots) con prezzi realistici.
    \item \textbf{OrderItemRequest ridisegnato} --- il campo \texttt{unitPrice}
          e \texttt{itemName} sono rimossi; il client invia solo
          \texttt{menuItemId} (UUID) e \texttt{quantity}.
    \item \textbf{buildItemsFromCatalog()} --- per ogni item richiesto, il service
          cerca il \texttt{MenuItem} per UUID. Se l'ID non esiste lancia
          \texttt{OrderValidationException} (HTTP 400). Il prezzo e' sempre
          quello del catalogo.
    \item \textbf{totalAmount ricalcolato server-side} --- dopo il lookup, il
          totale e' la somma di \texttt{menuItem.getPrice() * quantity}.
          Nessun valore del client partecipa al calcolo.
\end{enumerate}

\begin{securebox}[title={Codice Sicuro --- OrderItemRequest.java (fix)}]
\begin{lstlisting}[style=java]
// SICURO: nessun campo prezzo esposto al client
public record OrderItemRequest(
    @NotNull UUID menuItemId,   // riferimento al catalogo server-side
    @NotNull @Positive Integer quantity
) {}
\end{lstlisting}
\end{securebox}

\begin{securebox}[title={Codice Sicuro --- buildItemsFromCatalog() nel service}]
\begin{lstlisting}[style=java]
// SICURO: prezzo sempre letto dal catalogo DB, mai dal client
private List<OrderItem> buildItemsFromCatalog(
        List<OrderItemRequest> itemRequests, RestaurantOrder order) {

    List<OrderItem> items = new ArrayList<>();
    for (OrderItemRequest req : itemRequests) {
        MenuItem menuItem = menuItemRepository.findById(req.menuItemId())
            .orElseThrow(() -> new OrderValidationException(
                "MENU_ITEM_NOT_FOUND: " + req.menuItemId()));

        OrderItem item = OrderItem.builder()
            .restaurantOrder(order)
            .itemName(menuItem.getName())
            .unitPrice(menuItem.getPrice())  // server-side — mai dal client
            .quantity(req.quantity())
            .build();
        items.add(item);
    }
    return items;
}

// SICURO: totalAmount calcolato sui prezzi del catalogo
BigDecimal totalAmount = items.stream()
    .map(i -> i.getUnitPrice().multiply(BigDecimal.valueOf(i.getQuantity())))
    .reduce(BigDecimal.ZERO, BigDecimal::add);
order.setTotalAmount(totalAmount);
\end{lstlisting}
\end{securebox}

\begin{securebox}[title={V2\_\_add\_menu\_items.sql --- catalogo prezzi server-side}]
\begin{lstlisting}[style=sql]
CREATE TABLE IF NOT EXISTS menu_items (
    id     UUID           NOT NULL DEFAULT gen_random_uuid(),
    name   VARCHAR(255)   NOT NULL,
    price  NUMERIC(10,2)  NOT NULL,
    active BOOLEAN        NOT NULL DEFAULT TRUE,
    ...
    CONSTRAINT chk_menu_items_price CHECK (price >= 0)
);

-- Seed: prezzi autoritativi, non modificabili dal client
INSERT INTO menu_items (id, name, price) VALUES
    ('a1000000-...01', 'Espresso',          2.50),
    ('a1000000-...06', 'Club Sandwich',    12.00),
    ('a1000000-...09', 'Bistecca ai Ferri', 22.00);
\end{lstlisting}
\end{securebox}

\paragraph{Riferimento} Commit: \commitref{d7af61c}{fix(security): T-FB-02 --- server-side price catalog \\
on fb-service}

% ─────────────────────────────────────────────────────────────────────────────
\subsection{Validazione importi pagamento \threat{T-BILL-02}}
% ─────────────────────────────────────────────────────────────────────────────

\paragraph{Descrizione}
Il \texttt{billing-service} accettava importi di pagamento senza vincoli di
precisione decimale. Il DTO \texttt{PaymentRequest} era vincolato solo con
\texttt{@Positive} (no negativi, no zero), ma non limitava il numero di cifre
decimali. Un attaccante poteva inviare importi con precisione arbitraria
(es.\ \texttt{amount: 0.00000000001}) che: (1) bypassano \texttt{@Positive}
superando il controllo di validazione, (2) inquinano i calcoli di bilancio
per arrotondamento cumulativo, (3) possono causare errori di conversione
quando il valore viene persistito in colonne \texttt{NUMERIC(10,2)} di PostgreSQL.
Parallelamente, \texttt{transactionReference} non aveva un limite di lunghezza,
aprendo a payload di decine di KB con caratteri arbitrari.
Classificato OWASP A06:2025 --- Insecure Design.

\begin{insecurebox}[title={Codice Vulnerabile --- PaymentRequest.java (baseline)}]
\begin{lstlisting}[style=java]
// INSICURO: nessun vincolo di precisione decimale
public record PaymentRequest(
    @NotNull @Positive BigDecimal amount,   // accetta 0.00000000001
    @NotNull PaymentMethod paymentMethod,
    @NotBlank String transactionReference)  // nessun @Size: payload illimitato
{}

// Il service usa request.amount() raw per tutte le comparazioni:
final BigDecimal balanceDue = invoice.getTotalAmount().subtract(currentTotalPaid);
if (request.amount().compareTo(balanceDue) > 0) { ... }
final BigDecimal newTotalPaid = currentTotalPaid.add(request.amount());
\end{lstlisting}
\end{insecurebox}

\paragraph{Mitigazione applicata}
L'intervento opera su tre livelli ortogonali:

\begin{enumerate}
    \item \textbf{Precisione decimale} --- aggiunto
          \texttt{@Digits(integer = 10, fraction = 2)} su
          \texttt{PaymentRequest.amount} e su \texttt{InvoiceRequest.totalAmount}.
          Jakarta Bean Validation rifiuta con HTTP 400 qualsiasi importo con pi\`u
          di 2 cifre decimali o pi\`u di 10 cifre intere, allineato allo standard
          ISO 4217 per valute a centesimo.
    \item \textbf{Lunghezza e formato del riferimento transazione} ---
          aggiunto \texttt{@Size(max = 100)} e\\
          \texttt{@Pattern(regexp = "\^{}[A-Za-z0-9\textbackslash-\textbackslash\_/\textbackslash.]+\$")}
          su \texttt{transactionReference}: limita la lunghezza a 100 caratteri
          e blocca caratteri di iniezione (\texttt{<}, \texttt{>},
          \texttt{;}, \texttt{\&}, spazi).
    \item \textbf{Normalizzazione nel service} --- difesa in profondit\`a: \\
          \texttt{PaymentServiceImpl.addPayment()} normalizza l'importo a
          2 decimali con \texttt{setScale(2, RoundingMode.HALF\_UP)} prima
          di qualsiasi aritmetica e persiste il valore normalizzato
          tramite \texttt{payment.setAmount(paymentAmount)}.
\end{enumerate}

\begin{securebox}[title={Codice Sicuro --- PaymentRequest.java (fix)}]
\begin{lstlisting}[style=java]
// SICURO: precisione monetaria imposta a livello DTO
public record PaymentRequest(
    @NotNull(message = "Amount is required")
    @Positive(message = "Amount must be positive")
    @Digits(integer = 10, fraction = 2,
            message = "Amount must have at most 10 integer digits and 2 decimal places")
    BigDecimal amount,

    @NotNull(message = "Payment method is required")
    PaymentMethod paymentMethod,

    @NotBlank(message = "Transaction reference is required")
    @Size(max = 100, message = "Transaction reference must not exceed 100 characters")
    @Pattern(regexp = "^[A-Za-z0-9\\-_/\\.]+$",
            message = "Transaction reference contains invalid characters")
    String transactionReference) {}
\end{lstlisting}
\end{securebox}

\begin{securebox}[title={Codice Sicuro --- PaymentServiceImpl.java (normalizzazione defense-in-depth)}]
\begin{lstlisting}[style=java]
// SICURO: normalizzazione a 2 decimali prima di ogni aritmetica
final BigDecimal paymentAmount = request.amount()
        .setScale(2, RoundingMode.HALF_UP);   // defense-in-depth

if (paymentAmount.compareTo(balanceDue) > 0) {
    log.warn("[BILLING] PAYMENT_REJECTED | reason=PAYMENT_EXCEEDS_BALANCE"
            + " | amount={} | balanceDue={}", paymentAmount, balanceDue);
    throw new BillingValidationException("PAYMENT_EXCEEDS_BALANCE");
}

final Payment payment = paymentMapper.toEntity(request);
payment.setAmount(paymentAmount);             // valore normalizzato persistito
payment.setPaymentDate(LocalDateTime.now());
// ...
final BigDecimal newTotalPaid = currentTotalPaid.add(paymentAmount);
\end{lstlisting}
\end{securebox}

\paragraph{Effetto}
Un importo come \texttt{0.00000000001} viene ora rifiutato con HTTP 400 al
livello controller prima di raggiungere il service; un eventuale bypass
verrebbe comunque normalizzato a \texttt{0.00} dalla
\texttt{setScale} difensiva. Il riferimento transazione non pu\`o superare
100 caratteri n\'e contenere caratteri non whitelisted.
Due nuovi test coprono il ramo \texttt{INVOICE\_ALREADY\_PAID} e la
normalizzazione decimale (\texttt{100.005} $\rightarrow$ \texttt{100.01}).

\paragraph{Riferimento} Commit: \commitref{d346fe8}{fix(security): T-BILL-02 --- server-side payment amount \\
validation on billing-service}

% ─────────────────────────────────────────────────────────────────────────────
\subsection{Validazione date di prenotazione (range cross-field) \threat{T-RES-03}}
% ─────────────────────────────────────────────────────────────────────────────

\paragraph{Descrizione}
Il \texttt{ReservationRequest} DTO del \texttt{reservation-service} validava le date
di check-in e check-out singolarmente tramite \texttt{@FutureOrPresent}, ma mancava
di una validazione \emph{cross-field} che imponesse \texttt{checkOutDate > checkInDate}.

Questa lacuna consentiva due classi di input anomali:

\begin{enumerate}
  \item \textbf{Zero-night stay}: \texttt{checkIn == checkOut} --- la camera appare
        prenotata, ma la sovrapposizione non viene rilevata perché la condizione
        \texttt{r.checkOutDate > :checkIn} nella query non produce mai un risultato
        corretto.
  \item \textbf{Date invertite}: \texttt{checkOut < checkIn} --- la query di rilevamento overlap \\
        (\texttt{r.checkInDate < :checkOut AND r.checkOutDate > :checkIn})
        fallisce silenziosamente, rendendo la stanza erroneamente disponibile.
\end{enumerate}

Entrambi i casi passano la validazione standard, raggiungono il layer persistenza e
corrompono la logica di double-booking prevention introdotta con T-RES-01.
La mappatura OWASP è \textbf{A03 --- Injection} (input non sanificato che altera il
comportamento di una query parametrizzata).

\paragraph{Mitigazione applicata}

\textbf{Livello 1 --- Bean Validation (controller boundary)}

È stata introdotta l'annotazione \texttt{@ValidDateRange}, un vincolo di classe \\
(\texttt{@Target(ElementType.TYPE)}) implementato da \texttt{DateRangeValidator}:

\begin{insecurebox}[title={Codice Vulnerabile --- ReservationRequest.java (prima)}]
\begin{lstlisting}[style=java]
// Nessuna validazione cross-field: entrambe le date passano
// @FutureOrPresent anche se checkOut <= checkIn
public record ReservationRequest(
    @NotNull @FutureOrPresent LocalDate checkInDate,
    @NotNull @FutureOrPresent LocalDate checkOutDate,
    ...
) { }
\end{lstlisting}
\end{insecurebox}

\begin{securebox}[title={Codice Sicuro --- ReservationRequest.java (dopo)}]
\begin{lstlisting}[style=java]
@ValidDateRange   // class-level: checkOutDate > checkInDate
public record ReservationRequest(
    @NotNull @FutureOrPresent LocalDate checkInDate,
    @NotNull @FutureOrPresent LocalDate checkOutDate,
    ...
) { }
\end{lstlisting}
\end{securebox}

\begin{securebox}[title={Codice Sicuro --- DateRangeValidator.java}]
\begin{lstlisting}[style=java]
public class DateRangeValidator
    implements ConstraintValidator<ValidDateRange, ReservationRequest> {

  @Override
  public boolean isValid(ReservationRequest request,
                         ConstraintValidatorContext ctx) {
    if (request == null
        || request.checkInDate() == null
        || request.checkOutDate() == null) {
      return true;  // null-check delegato a @NotNull
    }
    return request.checkOutDate().isAfter(request.checkInDate());
  }
}
\end{lstlisting}
\end{securebox}

Una violazione produce HTTP 400 con corpo RFC 7807 tramite il
\texttt{\seqsplit{MethodArgumentNotValidException}} handler già presente nel
\texttt{GlobalExceptionHandler}.

\textbf{Livello 2 --- Guard nel service (difesa in profondità)}

Per proteggere i chiamanti programmatici che bypassano la validazione del
controller, è stato aggiunto un guard statico \texttt{verifyDateRange()} invocato
come prima operazione in \texttt{createReservation()} e \texttt{updateReservation()},
prima di qualsiasi accesso a servizi esterni o al DB:

\begin{securebox}[title={Codice Sicuro --- ReservationServiceImpl.java}]
\begin{lstlisting}[style=java]
private static void verifyDateRange(ReservationRequest request) {
    if (request.checkInDate() != null
        && request.checkOutDate() != null
        && !request.checkOutDate().isAfter(request.checkInDate())) {
        throw new BadRequestException("CHECKOUT_MUST_BE_AFTER_CHECKIN");
    }
}
\end{lstlisting}
\end{securebox}

\paragraph{Test introdotti}
Sei nuovi test coprono la mitigazione:
\begin{itemize}
  \item \texttt{shouldRejectCreateWhenCheckOutSameDayAsCheckIn} --- zero-night stay su create
  \item \texttt{shouldRejectCreateWhenCheckOutBeforeCheckIn} --- date invertite su create
  \item \texttt{shouldRejectUpdateWhenCheckOutSameDayAsCheckIn} --- zero-night stay su update
  \item \texttt{shouldPassWhenCheckOutIsAfterCheckIn} --- caso valido nel validator
  \item \texttt{shouldFailWhenCheckOutSameDayAsCheckIn} --- validator diretto
  \item \texttt{shouldFailWhenCheckOutBeforeCheckIn} --- validator diretto con range invertito
\end{itemize}

\paragraph{Riferimento} Commit: \commitref{43be320}{fix(security): T-RES-03 --- cross-field date range validation on reservation-service}

% ─────────────────────────────────────────────────────────────────────────────
\section{Autorizzazione e Controllo degli Accessi}
% ─────────────────────────────────────────────────────────────────────────────

\subsection{Strip Header X-Auth dal client \threat{T-GW-01}}
\textit{Mitigato nella baseline tramite HMAC-SHA256 (InternalAuthFilter). Documentato nella
sezione \textbf{Autenticazione Sicura} sopra (T-AUTH-01/T-GW-01/T-FE-02).}

\subsection{IDOR su risorse ospite/hotel \threat{T-GST-01, T-RES-02, T-BILL-01}}
Le vulnerabilità IDOR che consentivano l'accesso cross-tenant a ospiti (T-GST-01),
prenotazioni (T-RES-02) e fatture (T-BILL-01) sono state risolte con lo stesso pattern
architetturale applicato in cascata a tutti e tre i servizi.
Il dettaglio di ciascun intervento è documentato nell'appendice storica
(Capitolo~\ref{ch:storico}) alle sezioni dedicate ai rispettivi servizi.

% ─────────────────────────────────────────────────────────────────────────────
\subsection{IDOR su ordini F\&B --- Accesso Cross-Hotel \threat{T-FB-01}}
% ─────────────────────────────────────────────────────────────────────────────

\paragraph{Descrizione}
Il \texttt{fb-service} gestisce gli ordini del ristorante addebitati al conto camera.
Nella baseline, l'entità \texttt{RestaurantOrder} non possedeva il campo \texttt{hotel\_id}
e tutti i metodi del repository operavano senza scoping multi-tenant:
\texttt{findAll(pageable)} restituiva \emph{tutti} gli ordini di \emph{tutti} gli hotel,
e \texttt{findByStayId(stayId)} restituiva gli ordini di qualsiasi hotel purché il
\texttt{stayId} fosse noto all'attaccante.
Un operatore autenticato su Hotel~A poteva quindi enumerare e leggere gli ordini F\&B
dell'Hotel~B con una semplice richiesta \texttt{GET /api/v1/fb/orders}.
Classificato \textbf{OWASP A01 --- Broken Access Control / IDOR}.

\begin{insecurebox}[title={Codice Vulnerabile --- RestaurantOrderRepository.java (baseline)}]
\begin{lstlisting}[style=java]
// INSICURO: nessun filtro hotel_id — tutti gli ordini di tutti gli hotel
@Repository
public interface RestaurantOrderRepository
        extends JpaRepository<RestaurantOrder, UUID> {

    // Restituisce gli ordini di QUALSIASI hotel che abbiano questo stayId
    List<RestaurantOrder> findByStayId(UUID stayId);

    // findAll(pageable) ereditato da JpaRepository: restituisce TUTTI gli ordini
}
\end{lstlisting}
\end{insecurebox}

\begin{insecurebox}[title={Codice Vulnerabile --- RestaurantOrderServiceImpl.java (baseline)}]
\begin{lstlisting}[style=java]
// getAllOrders: nessun filtro hotel
public Page<RestaurantOrderResponse> getAllOrders(Pageable pageable) {
    return orderRepository.findAll(pageable)   // cross-hotel leak
            .map(orderMapper::toResponse);
}

// getOrdersByStayId: nessun check ownership
public List<RestaurantOrderResponse> getOrdersByStayId(UUID stayId) {
    return orderRepository.findByStayId(stayId)  // cross-hotel IDOR
            .stream().map(orderMapper::toResponse).toList();
}

// createOrder: hotel_id non impostato sull'entita'
public RestaurantOrderResponse createOrder(RestaurantOrderRequest request) {
    RestaurantOrder order = orderMapper.toEntity(request);
    // hotelId mai valorizzato: ordine non scoped a nessun hotel
    order.setStatus(OrderStatus.BILLED_TO_ROOM);
    ...
}
\end{lstlisting}
\end{insecurebox}

\paragraph{Mitigazione applicata}
L'intervento segue il medesimo pattern già applicato a \\
\texttt{guest-service},\texttt{reservation-service} e \texttt{billing-service}:

\begin{enumerate}
  \item \textbf{Flyway V3}: aggiunta colonna \texttt{hotel\_id UUID NOT NULL} alla tabella \\
        \texttt{restaurant\_orders} con indici compositi \texttt{(hotel\_id)} e
        \texttt{(hotel\_id, stay\_id)} per le query più frequenti.
        La stessa migration corregge anche il CHECK constraint sullo stato \\
        (\texttt{PENDING/PREPARED/DELIVERED/BILLED\_TO\_ROOM}), che nella baseline
        elencava valori obsoleti \\
        (\texttt{OPEN/IN\_PROGRESS/SERVED/CANCELLED}).
  \item \textbf{InternalAuthFilter}: lettura header \texttt{X-Auth-Hotel} iniettato
        dall'API Gateway; rifiuto con \texttt{401} se assente; valore memorizzato in
        \texttt{auth.setDetails(hotelId)} e disponibile in tutto il thread corrente.
  \item \textbf{FeignHeaderConfig}: propagazione di \texttt{X-Auth-Hotel} su tutte le
        chiamate Feign in uscita (es.\ verso \texttt{stay-service}).
  \item \textbf{RestaurantOrderServiceImpl}: metodo \texttt{resolveHotelId()} estrae
        \texttt{hotel\_id} dal \texttt{SecurityContext} (mai dal client); tutte le
        operazioni usano query hotel-scoped; \texttt{ORDER\_CREATED} loggato con
        \texttt{orderId+stayId+hotelId+total} (prefisso \texttt{[FB]}).
  \item \textbf{RestaurantOrderRepository}: nuovi metodi \\
        \texttt{findByStayIdAndHotelId} e \texttt{findAllByHotelId(Pageable)}
        al posto delle query non scoped.
\end{enumerate}

\begin{securebox}[title={Codice Sicuro --- InternalAuthFilter.java (lettura X-Auth-Hotel)}]
\begin{lstlisting}[style=java]
// T-FB-01: legge X-Auth-Hotel e lo rifiuta se assente
private static final String HEADER_HOTEL = "X-Auth-Hotel";

final String hotelId = request.getHeader(HEADER_HOTEL);
if (!StringUtils.hasText(hotelId)) {
    LOG.warn("[InternalAuthFilter] Missing X-Auth-Hotel header for user={}",
             username);
    rejectRequest(response, "Missing X-Auth-Hotel header"); // 401
    return;
}

UsernamePasswordAuthenticationToken auth =
    new UsernamePasswordAuthenticationToken(
        username, "", List.of(new SimpleGrantedAuthority("ROLE_" + role)));
auth.setDetails(hotelId);   // disponibile via auth.getDetails() nel service
SecurityContextHolder.getContext().setAuthentication(auth);
\end{lstlisting}
\end{securebox}

\begin{securebox}[title={Codice Sicuro --- RestaurantOrderServiceImpl.java (hotel-scoped)}]
\begin{lstlisting}[style=java]
// Estrae hotel_id dal SecurityContext; mai accettato dal client (T-FB-01)
private UUID resolveHotelId() {
    final Authentication auth =
        SecurityContextHolder.getContext().getAuthentication();
    final Object details = auth.getDetails();
    if (!(details instanceof String hotelIdStr)
            || !StringUtils.hasText(hotelIdStr)) {
        throw new IllegalStateException("X-Auth-Hotel_MISSING");
    }
    return UUID.fromString(hotelIdStr);
}

// createOrder: hotel_id impostato server-side
public RestaurantOrderResponse createOrder(RestaurantOrderRequest request) {
    final UUID hotelId = resolveHotelId();
    RestaurantOrder order = orderMapper.toEntity(request);
    order.setHotelId(hotelId);   // mai dal client
    ...
    log.info("[FB] ORDER_CREATED | orderId={} | stayId={} | hotelId={} | total={}",
             savedOrder.getId(), savedOrder.getStayId(), hotelId,
             savedOrder.getTotalAmount());
}

// getAllOrders: solo ordini dell'hotel autenticato
public Page<RestaurantOrderResponse> getAllOrders(Pageable pageable) {
    final UUID hotelId = resolveHotelId();
    return orderRepository.findAllByHotelId(hotelId, pageable)
                          .map(orderMapper::toResponse);
}

// getOrdersByStayId: IDOR-safe — lista vuota se stayId di altro hotel
public List<RestaurantOrderResponse> getOrdersByStayId(UUID stayId) {
    final UUID hotelId = resolveHotelId();
    return orderRepository.findByStayIdAndHotelId(stayId, hotelId)
                          .stream().map(orderMapper::toResponse).toList();
}
\end{lstlisting}
\end{securebox}

\paragraph{Test} Aggiunti/aggiornati 8 test in \texttt{RestaurantOrderServiceImplTest}:
mock del \texttt{SecurityContext} con \texttt{hotelId} in \texttt{@BeforeEach},
\texttt{clearContext()} in \texttt{@AfterEach};
\texttt{\seqsplit{shouldSetHotelIdFromSecurityContextOnCreate}} (verifica che \texttt{order.getHotelId()} == hotelId da SecurityContext),
\texttt{\seqsplit{shouldReturnEmptyListForWrongHotelOnStayLookup}} (IDOR scenario --- lista vuota per stayId di altro hotel),
\texttt{\seqsplit{shouldReturnOnlyHotelOrdersOnPaginatedListing}} (getAllOrders usa \texttt{findAllByHotelId}).

\paragraph{Riferimento} Commit: \commitref{18a3768}{fix(security): T-FB-01 --- IDOR fix, hotel\_id scope \\
on fb-service orders}

\subsection{Double Booking Prevention \threat{T-RES-01}}

\paragraph{Descrizione}
In un sistema di prenotazione alberghiera, il \emph{double booking} si verifica quando due
richieste concorrenti riescono a prenotare la stessa camera per date sovrapposte.
La condizione di gara (race condition) è del tipo \emph{Time-Of-Check/Time-Of-Use} (TOCTOU):
entrambe le transazioni superano il controllo di disponibilità prima che l'altra abbia
effettuato la scrittura. Classificato OWASP A06:2025 --- Insecure Design.

La baseline non implementava alcun meccanismo di locking:

\begin{insecurebox}[title={Codice Vulnerabile --- Reservation.java (baseline)}]
\begin{lstlisting}[style=java]
@Entity
@Table(name = "reservations")
public class Reservation {
    // Nessun @Version: aggiornamenti concorrenti alla stessa
    // prenotazione sovrascrivono silenziosamente i dati altrui.
    @Id
    @GeneratedValue(strategy = GenerationType.UUID)
    private UUID id;
    // ...
}
\end{lstlisting}
\end{insecurebox}

\begin{insecurebox}[title={Codice Vulnerabile --- ReservationRepository.java (baseline)}]
\begin{lstlisting}[style=java]
// Nessun @Lock: il SELECT sul controllo overlap non acquisisce
// alcun lock a livello DB. Due transazioni simultanee possono
// entrambe passare il check e inserire due prenotazioni
// per la stessa camera nelle stesse date.
List<Reservation> findOverlappingReservationsForNew(
    List<UUID> roomIds,
    LocalDate checkIn,
    LocalDate checkOut
);
\end{lstlisting}
\end{insecurebox}

\paragraph{Mitigazione applicata}
La mitigazione agisce su due livelli complementari:

\begin{itemize}
  \item \textbf{Optimistic locking JPA} (\texttt{@Version}): il campo \texttt{version} è
        incrementato ad ogni \texttt{UPDATE}. Se due transazioni leggono la stessa versione
        e tentano di aggiornare la stessa prenotazione, la seconda riceve
        \texttt{ObjectOptimisticLockingFailureException} che viene tradotta in HTTP 409
        Conflict dal \texttt{GlobalExceptionHandler}.

  \item \textbf{Pessimistic write lock sul controllo overlap} (\texttt{@Lock(PESSIMISTIC\_WRITE)}):
        la query \texttt{findOverlappingReservationsForNew} esegue un
        \texttt{SELECT \ldots\ FOR UPDATE} sulle prenotazioni esistenti che si sovrappongono
        alle date richieste. Questo serializza le richieste concorrenti per la stessa camera:
        la seconda transazione viene messa in attesa finché la prima non fa commit, eliminando
        la finestra TOCTOU.

  \item \textbf{Formula date-overlap} già corretta:\\
        \texttt{checkInDate < :checkOut AND checkOutDate > :checkIn},
        che copre tutti i casi di sovrapposizione parziale e totale.

  \item \textbf{Flyway migration V4}: aggiunge la colonna \texttt{version BIGINT NOT NULL
        DEFAULT 0} alla tabella \texttt{reservations}.
\end{itemize}

\begin{securebox}[title={Codice Sicuro --- Reservation.java (feature/secure-coding-hardening)}]
\begin{lstlisting}[style=java]
@Entity
@Table(name = "reservations")
public class Reservation {

    @Id
    @GeneratedValue(strategy = GenerationType.UUID)
    private UUID id;

    // Optimistic locking: incrementato ad ogni UPDATE.
    // Conflitto concorrente -> ObjectOptimisticLockingFailureException
    // -> HTTP 409 Conflict.
    @Version
    @Column(nullable = false)
    private Long version;

    // ...
}
\end{lstlisting}
\end{securebox}

\begin{securebox}[title={Codice Sicuro --- ReservationRepository.java (feature/secure-coding-hardening)}]
\begin{lstlisting}[style=java]
// SELECT ... FOR UPDATE: serializza le richieste concorrenti
// sulla stessa camera. La seconda transazione attende il commit
// della prima, eliminando la race condition TOCTOU.
@Lock(LockModeType.PESSIMISTIC_WRITE)
@Query("""
    SELECT r FROM Reservation r
    JOIN r.lineItems li
    WHERE r.active = true
    AND r.status NOT IN (CANCELLED, NO_SHOW)
    AND li.roomId IN :roomIds
    AND li.active = true
    AND r.checkInDate < :checkOut
    AND r.checkOutDate > :checkIn
    """)
List<Reservation> findOverlappingReservationsForNew(
    @Param("roomIds") List<UUID> roomIds,
    @Param("checkIn") LocalDate checkIn,
    @Param("checkOut") LocalDate checkOut
);
\end{lstlisting}
\end{securebox}

\begin{securebox}[title={V4\_\_add\_version\_to\_reservations.sql --- Flyway migration}]
\begin{lstlisting}[style=yaml]
ALTER TABLE reservations
    ADD COLUMN IF NOT EXISTS version BIGINT NOT NULL DEFAULT 0;
\end{lstlisting}
\end{securebox}

\begin{securebox}[title={GlobalExceptionHandler.java --- Handler HTTP 409}]
\begin{lstlisting}[style=java]
@ExceptionHandler(ObjectOptimisticLockingFailureException.class)
public ProblemDetail handleOptimisticLockingFailure(
        ObjectOptimisticLockingFailureException ex) {
    LOG.warn("[OptimisticLock] Concurrent modification: {}", ex.getMessage());
    ProblemDetail pd = ProblemDetail.forStatusAndDetail(
        HttpStatus.CONFLICT, "RESERVATION_CONCURRENT_MODIFICATION");
    pd.setTitle("Conflict");
    pd.setType(URI.create("https://hotelpms.com/errors/conflict"));
    return pd;
}
\end{lstlisting}
\end{securebox}

\paragraph{Test} Aggiunti 3 test in \texttt{ReservationServiceImplTest}:
\texttt{\seqsplit{testCreateReservationOverlapThrowsBadRequest}} (verifica che il check di overlap
rifiuti la prenotazione con HTTP 400),
\texttt{\seqsplit{testUpdateReservationOverlapThrowsBadRequest}} (stesso controllo sull'update),
\texttt{\seqsplit{testSaveThrowsOptimisticLockingFailurePropagates}} (verifica che
\texttt{\seqsplit{ObjectOptimisticLockingFailureException}} propaghi correttamente
al GlobalExceptionHandler per la traduzione in HTTP 409).

\paragraph{Riferimento SEI CERT} Regola \textbf{VNA02-J}: \emph{``Ensure that compound operations on shared variables are atomic''}. Il controllo di sovrapposizione e l'inserimento della prenotazione sono operazioni separate; senza locking il check-then-act non è atomico (race condition TOCTOU). Il \texttt{SELECT ... FOR UPDATE} serializza le transazioni concorrenti sulla stessa risorsa.

\paragraph{Nota OWASP A10:2025} Il \texttt{GlobalExceptionHandler} centralizzato copre anche \textbf{A10:2025 --- Mishandling of Exceptional Conditions}, categoria nuova dell'edizione 2025: la \texttt{ObjectOptimisticLockingFailureException} viene intercettata e tradotta in HTTP~409 con body RFC~7807 invece di propagarsi come 500 generico. Questo è il pattern opposto al \emph{fail-open} citato nelle slide del corso (es.\ e-commerce che conferma un ordine dopo timeout del gateway di pagamento).

\paragraph{Riferimento} Commit:\\
\commitref{ca9bf92}{fix(security): T-RES-01 --- optimistic lock + pessimistic write on overlap check}

% ─────────────────────────────────────────────────────────────────────────────
\section{Gestione Sicura delle Sessioni}
% ─────────────────────────────────────────────────────────────────────────────

\subsection{Cookie flags e JWT hardening \threat{T-FE-02} / \threat{T-AUTH-04}}

\paragraph{Descrizione}
Il trasporto del token JWT in un cookie con flag di sicurezza appropriati
costituisce la prima barriera contro le due categorie di furto di token più
comuni: XSS (lettura del token da JavaScript) e intercettazione passiva su canali
non cifrati.
Parallelamente, la gestione del ciclo di vita del JWT — scadenza breve,
rotazione del refresh token, revoca immediata al logout e al cambio password —
determina la finestra di esposizione residua in caso di furto riuscito.

\paragraph{Cookie httpOnly (T-FE-02 — baseline)}
Fin dalla baseline il server imposta il JWT in un cookie con i tre flag
fondamentali:

\begin{securebox}[title={Cookie JWT --- AuthController.java (branch \texttt{main})}]
\begin{lstlisting}[style=java]
// httpOnly:  JavaScript non puo' leggere il token (difesa XSS)
// Secure:    il cookie e' trasmesso solo su HTTPS
// SameSite=Strict: il browser non invia il cookie in richieste cross-site
private ResponseCookie createCookie(String value, int maxAge) {
    return ResponseCookie.from("jwt", value)
        .httpOnly(true)
        .secure(true)
        .sameSite("Strict")
        .path("/")
        .maxAge(maxAge)
        .build();
}
\end{lstlisting}
\end{securebox}

Il flag \texttt{httpOnly} elimina la classe di attacchi in cui uno script
XSS (anche iniettato via dipendenza npm compromessa) cerca di estrarre il
token con \texttt{document.cookie} o \texttt{localStorage.getItem}.
Il flag \texttt{Secure} protegge da intercettazione passiva su reti non
cifrate. \texttt{SameSite=Strict} impedisce l'invio automatico del cookie su
richieste che originano da siti terzi, costituendo la prima linea di difesa
CSRF prima ancora del Double-Submit-Cookie di T-GW-05.

\paragraph{Ciclo di vita del token (T-AUTH-04)}
Il ridisegno completo del ciclo di vita JWT --- access token da 15 minuti,
refresh token da 7 giorni con rotazione e blacklist Redis, revoca immediata
al logout, e token versioning che invalida tutte le sessioni al cambio password
--- è documentato in dettaglio nelle sottosezioni precedenti:

\begin{itemize}
  \item \textbf{\S~Refresh Token Rotation e Blacklist} (T-AUTH-04): la separazione
        access/refresh token riduce la finestra di esposizione da 24 ore a
        \textbf{15 minuti}; il logout revoca il refresh token server-side tramite
        blacklist Redis, chiudendo la sessione indipendentemente dalla scadenza
        naturale.
  \item \textbf{\S~Cambio password non revoca le sessioni attive} (T-AUTH-04 residuo):
        il meccanismo di token versioning (\texttt{tv} claim + Redis key
        \texttt{user:tv:<username>}) garantisce che un token emesso prima del
        cambio password venga rifiutato in fase di refresh, riducendo la finestra
        di esposizione residua ai soli access token in volo (max 15 minuti).
\end{itemize}

\paragraph{Effetto complessivo}
La combinazione dei tre livelli di protezione --- cookie httpOnly, access token
a breve scadenza, revoca server-side --- rende il furto di token una minaccia
gestibile: anche intercettando un access token, un attaccante dispone di una
finestra massima di 15 minuti prima che il token scada. Un refresh token
intercettato è inutile se il legittimo titolare ne fa uso prima (rotation) o se
viene revocato esplicitamente (logout, cambio password).

\subsection{Protezione CSRF \threat{T-GW-05}}

\paragraph{Descrizione}
Il sistema utilizza cookie \texttt{httpOnly} con flag \texttt{SameSite=Strict} per trasportare
il JWT di sessione. Il flag \texttt{SameSite=Strict} è una prima linea di difesa efficace,
ma non sufficiente: browser datati, proxy aziendali e bug specifici di implementazione possono
violarne il rispetto. In assenza di un meccanismo esplicito di verifica dell'origine della
richiesta, un sito malevolo può indurre il browser autenticato dell'utente a inviare richieste
mutanti (POST, PUT, PATCH, DELETE) verso il sistema --- l'attacco noto come
\textbf{Cross-Site Request Forgery (CSRF)}.

Nella baseline non era presente alcuna protezione CSRF esplicita: ogni richiesta con un cookie
JWT valido veniva accettata dall'API Gateway indipendentemente dall'origine.
Classificato OWASP A01 --- Broken Access Control.

\begin{insecurebox}[title={Pattern Vulnerabile --- assenza CSRF token (baseline)}]
\begin{lstlisting}[style=java]
// INSICURO: il gateway accetta qualsiasi richiesta con cookie JWT valido,
// indipendentemente dall'origine. Un sito malevolo puo' scatenare:
//   POST https://hotel.example.com/api/v1/guests
//   Cookie: jwt=<token valido> (inviato automaticamente dal browser)
// senza che il sito possa leggere il token -- ma la richiesta viene eseguita.

// api-gateway: nessun filtro CSRF nella baseline
// -> tutte le richieste POST/PUT/PATCH/DELETE passano se il cookie JWT e' valido
\end{lstlisting}
\end{insecurebox}

\paragraph{Mitigazione applicata --- Double Submit Cookie}
L'intervento implementa il pattern \textbf{Double Submit Cookie} (OWASP CSRF Prevention Cheat
Sheet), che è stateless e si integra naturalmente con l'architettura SPA + JWT:

\begin{enumerate}
  \item \textbf{Emissione del token CSRF} (\texttt{auth-service}): all'esito positivo di
        \texttt{/login}, \texttt{/register} e \texttt{/refresh}, il controller emette un
        cookie \texttt{csrf\_token} contenente un UUID casuale, con flag
        \texttt{Secure; SameSite=Strict} ma senza \texttt{httpOnly}, in modo che il codice
        JavaScript della SPA possa leggerlo tramite \texttt{document.cookie}.
  \item \textbf{Invio dell'header} (frontend): un Axios request interceptor legge il cookie
        \texttt{csrf\_token} e, per ogni richiesta mutante, lo inietta nell'header
        \texttt{X-CSRF-Token}.
  \item \textbf{Validazione} (\texttt{api-gateway/CsrfFilter}): un \texttt{GlobalFilter}
        con ordine \texttt{HIGHEST\_PRECEDENCE + 2} confronta il valore dell'header con
        quello del cookie. Se non corrispondono --- o se uno dei due è assente --- risponde
        con HTTP 403. I percorsi pre-autenticazione (\texttt{/auth/login},
        \texttt{/auth/register}) sono esclusi perché il token non è ancora stato emesso.
\end{enumerate}

\textbf{Perché è sicuro}: un sito malevolo può indurre il browser a inviare il cookie
(questo è il CSRF stesso), ma la \emph{same-origin policy} del browser impedisce al
JavaScript esterno di leggere il valore del cookie \texttt{csrf\_token}. Senza poter leggere
il cookie, il sito malevolo non può costruire l'header \texttt{X-CSRF-Token} corrispondente,
e la richiesta viene rifiutata dal gateway.

\begin{securebox}[title={Codice Sicuro --- CsrfFilter.java (api-gateway)}]
\begin{lstlisting}[style=java]
@Component
public class CsrfFilter implements GlobalFilter, Ordered {

    static final String CSRF_COOKIE_NAME = "csrf_token";
    static final String CSRF_HEADER_NAME = "X-CSRF-Token";

    private static final Set<HttpMethod> SAFE_METHODS = Set.of(
            HttpMethod.GET, HttpMethod.HEAD, HttpMethod.OPTIONS, HttpMethod.TRACE);

    private static final Set<String> EXCLUDED_PATHS = Set.of(
            "/api/v1/auth/login", "/api/v1/auth/register");

    @Override
    public Mono<Void> filter(ServerWebExchange exchange, GatewayFilterChain chain) {
        HttpMethod method = exchange.getRequest().getMethod();
        if (method == null || SAFE_METHODS.contains(method)) {
            return chain.filter(exchange); // metodi sicuri: nessun controllo
        }
        String path = exchange.getRequest().getPath().value();
        if (EXCLUDED_PATHS.contains(path)) {
            return chain.filter(exchange); // percorsi pre-auth: esclusi
        }

        HttpCookie csrfCookie = exchange.getRequest()
                .getCookies().getFirst(CSRF_COOKIE_NAME);
        String headerValue = exchange.getRequest()
                .getHeaders().getFirst(CSRF_HEADER_NAME);

        if (csrfCookie == null || csrfCookie.getValue().isBlank()
                || headerValue == null || headerValue.isBlank()
                || !csrfCookie.getValue().equals(headerValue)) {
            // Cookie assente/blank o mismatch -> CSRF rilevato -> 403
            exchange.getResponse().setStatusCode(HttpStatus.FORBIDDEN);
            return exchange.getResponse().setComplete();
        }
        return chain.filter(exchange); // token valido: passa
    }

    @Override
    public int getOrder() {
        return Ordered.HIGHEST_PRECEDENCE + 2; // dopo SecurityHeadersFilter (+1)
    }
}
\end{lstlisting}
\end{securebox}

\begin{securebox}[title={Codice Sicuro --- AuthController.java (auth-service): emissione csrf\_token}]
\begin{lstlisting}[style=java]
// Emesso su login/register/refresh, cancellato al logout
private static ResponseCookie createCsrfCookie(String value, int maxAge) {
    return ResponseCookie.from("csrf_token", value)
            .httpOnly(false)         // JS deve poterlo leggere
            .secure(true)
            .path("/")
            .maxAge(maxAge)
            .sameSite("Strict")
            .build();
}

// Nel metodo login():
return ResponseEntity.ok()
    .header(SET_COOKIE, createCookie("jwt", ...).toString())
    .header(SET_COOKIE, createCookie("refresh_token", ...).toString())
    .header(SET_COOKIE,
        createCsrfCookie(UUID.randomUUID().toString(), ACCESS_MAX_AGE).toString())
    .build();
\end{lstlisting}
\end{securebox}

\begin{securebox}[title={Codice Sicuro --- api.ts (frontend): Axios request interceptor}]
\begin{lstlisting}[style=typescript]
// Legge il cookie csrf_token (non-httpOnly) e lo inietta come header
function getCsrfToken(): string | null {
  const match = document.cookie.match(/(?:^|;\s*)csrf_token=([^;]*)/);
  return match ? decodeURIComponent(match[1]) : null;
}

api.interceptors.request.use((config) => {
  const method = config.method?.toUpperCase();
  if (method && ['POST','PUT','PATCH','DELETE'].includes(method)) {
    const token = getCsrfToken();
    if (token) config.headers['X-CSRF-Token'] = token;
  }
  return config;
});
\end{lstlisting}
\end{securebox}

\paragraph{Test} Aggiunti 10 test in \texttt{CsrfFilterTest}:
metodi sicuri (GET, HEAD, OPTIONS) passano senza token; percorsi esclusi (login, register) passano senza token;
POST/PUT/PATCH/DELETE senza cookie → 403; senza header → 403; cookie blank → 403; header blank → 403;
token mismatch → 403; token valido → 200 (POST, PUT, DELETE, PATCH).

\paragraph{Riferimento} Commit: \commitref{e146b17}{fix(security): T-GW-05 --- CSRF protection via Double Submit Cookie}

% ─────────────────────────────────────────────────────────────────────────────
\section{Configurazioni di Sicurezza (Difesa in Profondità)}
% ─────────────────────────────────────────────────────────────────────────────

\subsection{Security Headers HTTP \threat{T-GW-03}}

\paragraph{Descrizione}
Nella baseline, l'API Gateway non iniettava alcun header HTTP di sicurezza nelle risposte.
Un browser che riceva una risposta priva di questi header è esposto a una serie di vettori
di attacco ben documentati:
\begin{itemize}
  \item Assenza di \texttt{Strict-Transport-Security} (HSTS): un attaccante on-path può
        forzare il downgrade della connessione da HTTPS a HTTP.
  \item Assenza di \texttt{X-Content-Type-Options}: il browser può effettuare
        \emph{MIME sniffing}, eseguendo come script una risposta dichiarata come
        \texttt{application/json}.
  \item Assenza di \texttt{X-Frame-Options}: la pagina può essere incorporata in un
        \texttt{<iframe>} per attacchi di clickjacking.
  \item Assenza di \texttt{Referrer-Policy}: l'URL completo viene inviato come header
        \texttt{Referer} a terze parti, potenzialmente esponendo path con token o ID.
  \item Assenza di \texttt{Permissions-Policy}: le API del browser (fotocamera, microfono,
        geolocalizzazione) rimangono accessibili a script di terze parti iniettati.
\end{itemize}
Classificato OWASP A02:2025 --- Security Misconfiguration.

\begin{insecurebox}[title={Baseline vulnerabile --- nessun header di sicurezza nelle risposte}]
\begin{lstlisting}[style=java]
// INSICURO: l'API Gateway non aggiunge alcun header di sicurezza.
// HTTP Response example (baseline):
//   HTTP/1.1 200 OK
//   Content-Type: application/json
//   Transfer-Encoding: chunked
//   < nessun HSTS, X-Frame-Options, CSP, X-Content-Type-Options >
\end{lstlisting}
\end{insecurebox}

\paragraph{Mitigazione applicata}
È stato aggiunto il componente \texttt{SecurityHeadersFilter}, un \texttt{GlobalFilter}
Spring Cloud Gateway che si inserisce all'inizio della catena
(\texttt{Ordered.HIGHEST\_PRECEDENCE + 1}) e registra una callback \texttt{beforeCommit}
sulla risposta reattiva. La callback viene eseguita immediatamente prima che gli header HTTP
vengano inviati sulla rete, garantendo la presenza degli header su ogni risposta del gateway,
indipendentemente dalla rotta o dal microservizio downstream.

\begin{securebox}[title={SecurityHeadersFilter.java --- GlobalFilter (api-gateway)}]
\begin{lstlisting}[style=java]
@Component
public class SecurityHeadersFilter implements GlobalFilter, Ordered {

    static final String VALUE_HSTS    = "max-age=31536000; includeSubDomains";
    static final String VALUE_XCTO    = "nosniff";
    static final String VALUE_XFO     = "DENY";
    static final String VALUE_REFERRER = "no-referrer";
    static final String VALUE_PERMISSIONS =
        "camera=(), microphone=(), geolocation=()";
    // API responses sono JSON: nessun contenuto da renderizzare
    static final String VALUE_CSP =
        "default-src 'none'; frame-ancestors 'none'";

    @Override
    public Mono<Void> filter(ServerWebExchange exchange,
                             GatewayFilterChain chain) {
        exchange.getResponse().beforeCommit(() -> {
            HttpHeaders h = exchange.getResponse().getHeaders();
            h.set(HEADER_HSTS,        VALUE_HSTS);
            h.set(HEADER_XCTO,        VALUE_XCTO);
            h.set(HEADER_XFO,         VALUE_XFO);
            h.set(HEADER_REFERRER,    VALUE_REFERRER);
            h.set(HEADER_PERMISSIONS, VALUE_PERMISSIONS);
            h.set(HEADER_CSP,         VALUE_CSP);
            return Mono.empty();
        });
        return chain.filter(exchange);
    }

    @Override
    public int getOrder() { return Ordered.HIGHEST_PRECEDENCE + 1; }
}
\end{lstlisting}
\end{securebox}

\paragraph{Test} Aggiunti 7 test in \texttt{SecurityHeadersFilterTest}:
uno per ciascun header (verifica del valore esatto) e uno per la priorità del filtro.
Il test utilizza \texttt{MockServerWebExchange} per invocare la callback \texttt{beforeCommit}
tramite \texttt{response.setComplete()}, senza avviare alcun contesto Spring.

\paragraph{Riferimento} Commit: \commitref{b901a9f}{fix(security): T-GW-03 + T-FE-04 --- security headers and CSP}

% ─────────────────────────────────────────────────────────────────────────────

\subsection{Content Security Policy \threat{T-FE-04}}

\paragraph{Descrizione}
L'assenza di un header \texttt{Content-Security-Policy} nella risposta che serve il frontend
React lascia il browser privo di istruzioni su quali sorgenti di contenuto siano legittime.
In mancanza di una CSP, uno script iniettato via XSS può caricare risorse da origini
arbitrarie, esfiltrare dati o eseguire operazioni autenticate per conto della vittima.
Classificato OWASP A02:2025 --- Security Misconfiguration.

Problema secondario: la configurazione nginx baseline era definita come stringa \emph{inline}
nel \texttt{Dockerfile}, rendendo impossibile aggiungere header di risposta senza ricostruire
l'immagine con un'istruzione RUN lunga e non manutenibile. Il frontend effettuava inoltre
chiamate API all'indirizzo assoluto \texttt{http://localhost:8080}, rendendo \texttt{connect-src
'self'} inapplicabile (origini diverse per porta).

\begin{insecurebox}[title={Dockerfile baseline --- nginx inline senza header di sicurezza}]
\begin{lstlisting}[style=yaml]
FROM nginx:alpine
COPY --from=build /app/dist /usr/share/nginx/html
RUN rm /etc/nginx/conf.d/default.conf
# Config inline: nessun header di sicurezza, nessun proxy API
RUN echo "server { listen 80; location / { root /usr/share/nginx/html; \
    try_files $uri $uri/ /index.html; } }" \
    > /etc/nginx/conf.d/default.conf
\end{lstlisting}
\end{insecurebox}

\begin{insecurebox}[title={api.ts baseline --- URL assoluto, incompatibile con CSP strict}]
\begin{lstlisting}[style=typescript]
// INSICURO per CSP: URL assoluto con porta diversa
// -> connect-src 'self' non copre http://localhost:8080
const api = axios.create({
  baseURL: import.meta.env.VITE_API_URL || 'http://localhost:8080',
  withCredentials: true,
});
\end{lstlisting}
\end{insecurebox}

\paragraph{Mitigazione applicata}
L'intervento si articola in tre modifiche coordinate:

\begin{enumerate}
  \item \textbf{nginx.conf dedicato} con la CSP completa e tutti gli header di sicurezza,
        copiato nel container via \texttt{COPY nginx.conf /etc/nginx/conf.d/default.conf}.
  \item \textbf{Proxy nginx per \texttt{/api/}}: nginx inoltra tutte le richieste
        \texttt{/api/*} all'API Gateway interno (\texttt{http://api-gateway:8080}),
        rendendo le chiamate API same-origin dal punto di vista del browser e
        consentendo \texttt{connect-src 'self'}.
  \item \textbf{api.ts}: il \texttt{baseURL} è cambiato da \texttt{'http://localhost:8080'}
        a stringa vuota (\texttt{''}), così Axios usa URL relativi che il browser risolve
        sulla stessa origine --- compatibili sia con il proxy Vite in sviluppo sia con
        il proxy nginx in produzione.
\end{enumerate}

\begin{securebox}[title={nginx.conf --- CSP + proxy API + security headers}]
\begin{lstlisting}[style=yaml]
server {
    listen 80;
    root /usr/share/nginx/html;

    add_header Content-Security-Policy
      "default-src 'self'; script-src 'self';
       style-src 'self' 'unsafe-inline'; img-src 'self' data:;
       connect-src 'self'; frame-ancestors 'none';
       form-action 'self'; base-uri 'self'" always;
    add_header X-Content-Type-Options "nosniff"           always;
    add_header X-Frame-Options        "DENY"              always;
    add_header Referrer-Policy        "no-referrer"       always;

    # Proxy API: same-origin per connect-src 'self'
    location /api/ {
        proxy_pass http://api-gateway:8080;
        proxy_set_header X-Real-IP       $remote_addr;
        proxy_set_header X-Forwarded-For $proxy_add_x_forwarded_for;
    }

    location / { try_files $uri $uri/ /index.html; }
}
\end{lstlisting}
\end{securebox}

\begin{securebox}[title={api.ts --- baseURL relativo, same-origin e CSP-compatibile}]
\begin{lstlisting}[style=typescript]
// SICURO: URL relativo -> stesso host/porta del frontend
// Dev:  proxy Vite  /api/* -> localhost:8080
// Prod: proxy nginx /api/* -> api-gateway:8080
const api = axios.create({
  baseURL: import.meta.env.VITE_API_URL || '',
  withCredentials: true,
});
\end{lstlisting}
\end{securebox}

\paragraph{Nota su \texttt{style-src 'unsafe-inline'}}
TailwindCSS inietta classi di utilità come stili inline al runtime; rimuovere
\texttt{'unsafe-inline'} richiederebbe la generazione di nonce o hash per ogni stile,
incompatibile con il meccanismo di purging di Tailwind. Il rischio residuo è accettabile
poiché la CSP blocca comunque \texttt{script-src} all'origine stessa, neutralizzando
il vettore XSS principale.

\paragraph{Riferimento} Commit: \commitref{b901a9f}{fix(security): T-GW-03 + T-FE-04 --- security headers and CSP}

% ─────────────────────────────────────────────────────────────────────────────

\subsection{Output Encoding e Protezione XSS \threat{T-FE-01}}

\paragraph{Descrizione}
Il frontend React espone dati recuperati dalle API direttamente nel DOM (nomi ospiti,
email, descrizioni). Se un componente usasse \texttt{dangerouslySetInnerHTML} o
manipolazioni DOM dirette (\texttt{innerHTML}/\texttt{outerHTML}), un payload XSS
memorizzato nel database verrebbe interpretato come HTML dal browser, consentendo
l'esecuzione di codice arbitrario nel contesto dell'applicazione.
Classificato OWASP A05:2025 --- Injection.

L'audit del codebase ha confermato che nessun componente utilizza
\texttt{dangerously\-Set\-Inner\-HTML}: React JSX escapa automaticamente ogni espressione
renderizzata come contenuto. Tuttavia, in assenza di un guard statico, non vi era
nulla che impedisse l'introduzione futura del pattern vulnerabile.

\begin{insecurebox}[title={Vettore XSS: dangerouslySetInnerHTML con dato da API}]
\begin{lstlisting}[style=typescript]
// INSICURO: interpreta HTML dal dato API come markup reale
// Un payload "firstName": "<img src=x onerror=alert(document.cookie)>"
// verrebbe eseguito dal browser invece di mostrato come testo

function GuestName({ guest }: { guest: GuestResponseDTO }) {
  return (
    <td
      dangerouslySetInnerHTML={{
        __html: `${guest.firstName} ${guest.lastName}`,
      }}
    />
  );
}
\end{lstlisting}
\end{insecurebox}

\begin{insecurebox}[title={Assenza di guard statico --- nessuna regola ESLint impediva il pattern}]
\begin{lstlisting}[style=typescript]
// eslint.config.js baseline: nessuna regola su dangerouslySetInnerHTML
rules: {
  ...reactPerf.configs.recommended.rules,
  // (nessun vincolo XSS)
},
\end{lstlisting}
\end{insecurebox}

\paragraph{Mitigazione applicata}
L'intervento ha operato su due livelli complementari:

\begin{enumerate}
  \item \textbf{Enforcement statico via ESLint (\texttt{no-restricted-syntax})}:
        tre selettori AST vietano i sink XSS primari in tutta la codebase;
        la violazione causa un errore di build (non un warning) rendendo impossibile
        introdurre il pattern senza interrompere la pipeline CI.
  \item \textbf{Test di regressione XSS}: il componente \texttt{Guests} viene
        renderizzato con un payload \texttt{<img src=x onerror=alert(1)>} nel campo
        \texttt{firstName}. Il test verifica che (a) nessun elemento \texttt{<img>}
        con attributo \texttt{onerror} sia presente nel DOM e (b) il payload compaia
        come testo letterale --- comportamento garantito dal JSX escaping di React.
\end{enumerate}

\begin{securebox}[title={eslint.config.js --- no-restricted-syntax vieta i sink XSS}]
\begin{lstlisting}[style=typescript]
// T-FE-01: statically ban XSS sinks
'no-restricted-syntax': [
  'error',
  {
    selector:
      'JSXAttribute[name.name="dangerouslySetInnerHTML"]',
    message:
      '[T-FE-01] dangerouslySetInnerHTML is forbidden ' +
      '-- use React JSX expressions to prevent XSS.',
  },
  {
    selector:
      'AssignmentExpression[left.property.name="innerHTML"]',
    message:
      '[T-FE-01] Direct innerHTML assignment is forbidden ' +
      '-- use React JSX expressions to prevent XSS.',
  },
  {
    selector:
      'AssignmentExpression[left.property.name="outerHTML"]',
    message:
      '[T-FE-01] Direct outerHTML assignment is forbidden ' +
      '-- use React JSX expressions to prevent XSS.',
  },
],
\end{lstlisting}
\end{securebox}

\begin{securebox}[title={Guests.test.tsx --- test di regressione XSS (T-FE-01)}]
\begin{lstlisting}[style=typescript]
it('should escape XSS payload in guest fields (T-FE-01)', async () => {
  const xssFirst = '<img src=x onerror=alert(1)>';
  const xssEmail = '"><svg onload=alert(2)>';
  vi.mocked(guestService.getAllGuests).mockResolvedValueOnce([
    { id: '1', firstName: xssFirst, lastName: 'XSS',
      email: xssEmail, phone: '', city: '', country: '' },
  ] as never);

  const { container } = render(<Guests />);

  await waitFor(() => {
    // Nessun elemento iniettato: React non ha interpretato il payload come HTML
    expect(container.querySelector('img[onerror]')).toBeNull();
    expect(container.querySelector('svg[onload]')).toBeNull();
    // Payload presente come testo letterale (HTML-escaped da JSX)
    const cells = container.querySelectorAll('td');
    expect(cells[0]?.textContent).toContain(xssFirst);
  });
});
\end{lstlisting}
\end{securebox}

\paragraph{Rendering JSX e output encoding}
React, per scelta architetturale, tratta ogni espressione JSX \texttt{\{valore\}}
come testo puro: i caratteri \texttt{<}, \texttt{>} e \texttt{\&} vengono convertiti
nelle corrispondenti entità HTML prima dell'inserimento nel DOM virtuale.
L'unica eccezione esplicita è \texttt{dangerouslySetInnerHTML}, ora vietata
staticamente. Di conseguenza, qualunque dato restituito dalle API --- anche un payload
costruito ad arte --- viene visualizzato come stringa letterale senza possibilità di
esecuzione.

\paragraph{Riferimento} Commit: \commitref{3655906}{fix(security): T-FE-01 --- ESLint no-restricted-syntax bans dangerouslySetInnerHTML/innerHTML; XSS test on Guests}

\subsection{Route Guard basato sul Ruolo \threat{T-FE-03}}

\paragraph{Descrizione}
In React, nascondere un link di navigazione nel layout non impedisce l'accesso alla route
corrispondente: un utente malintenzionato può semplicemente digitare l'URL nel browser e
raggiungere la pagina. Nella baseline, la route \texttt{/owner-dashboard} era protetta
esclusivamente dalla condizione \texttt{isOwnerOrAdmin} nel \texttt{MainLayout} (UI hiding),
che rendeva invisibile la voce di menu per il ruolo \texttt{RECEPTIONIST}.
Il componente \texttt{ProtectedRoute} verificava unicamente lo stato di autenticazione,
senza alcun controllo sul ruolo. Classificato OWASP A01 --- Broken Access Control.

\begin{insecurebox}[title={ProtectedRoute.tsx --- baseline: solo verifica autenticazione}]
\begin{lstlisting}[style=typescript]
// INSICURO: nessun controllo sul ruolo utente
export const ProtectedRoute = () => {
  const isAuthenticated = useAuthStore((state) => state.isAuthenticated);

  if (!isAuthenticated) {
    return <Navigate to="/login" replace />;
  }

  return <Outlet />;  // qualsiasi utente autenticato accede
};
\end{lstlisting}
\end{insecurebox}

\begin{insecurebox}[title={App.tsx --- baseline: /owner-dashboard accessibile a tutti gli autenticati}]
\begin{lstlisting}[style=typescript]
// INSICURO: ProtectedRoute verifica solo l'autenticazione
<Route element={<ProtectedRoute />}>
  <Route element={<MainLayout />}>
    {/* ... altre route ... */}
    {/* Chiunque sia autenticato puo' digitare /owner-dashboard */}
    <Route path="/owner-dashboard" element={<OwnerDashboard />} />
  </Route>
</Route>
\end{lstlisting}
\end{insecurebox}

\paragraph{Intervento}
\texttt{ProtectedRoute} è stato esteso con una prop opzionale \texttt{allowedRoles?: Role[]}.
Se specificata, il componente verifica che \texttt{user.role} sia incluso nella allowlist;
in caso contrario, redirige verso \texttt{/} (home). La route \texttt{/owner-dashboard}
è ora annidata dentro un secondo \texttt{ProtectedRoute} con
\texttt{allowedRoles=\{['OWNER', 'ADMIN']\}}, garantendo l'enforcement a livello di routing
e non solo a livello di UI. L'array viene dichiarato come costante di modulo
(\texttt{OWNER\_ADMIN\_ROLES}) per rispettare la regola ESLint
\texttt{react-perf/jsx-no-new-array-as-prop}.

\begin{securebox}[title={ProtectedRoute.tsx --- con prop allowedRoles (T-FE-03)}]
\begin{lstlisting}[style=typescript]
import type { Role } from '../types/auth.types';

interface ProtectedRouteProps {
  allowedRoles?: Role[];
}

export const ProtectedRoute = (
  { allowedRoles }: ProtectedRouteProps = {}
) => {
  const isAuthenticated = useAuthStore((state) => state.isAuthenticated);
  const user = useAuthStore((state) => state.user);

  if (!isAuthenticated) {
    return <Navigate to="/login" replace />;
  }

  if (allowedRoles && (!user || !allowedRoles.includes(user.role))) {
    return <Navigate to="/" replace />;  // ruolo non autorizzato
  }

  return <Outlet />;
};
\end{lstlisting}
\end{securebox}

\begin{securebox}[title={App.tsx --- /owner-dashboard con doppio guard (T-FE-03)}]
\begin{lstlisting}[style=typescript]
// Costante di modulo: evita new array ad ogni render
const OWNER_ADMIN_ROLES = ['OWNER', 'ADMIN'] as const;

// ...

<Route element={<ProtectedRoute />}>        {/* guard: autenticazione */}
  <Route element={<MainLayout />}>
    {/* ... altre route (accessibili a tutti gli autenticati) ... */}
    <Route element={
      <ProtectedRoute allowedRoles={OWNER_ADMIN_ROLES} />  {/* guard: ruolo */}
    }>
      <Route path="/owner-dashboard" element={<OwnerDashboard />} />
    </Route>
  </Route>
</Route>
\end{lstlisting}
\end{securebox}

\paragraph{Test di verifica}
Quattro nuovi casi in \texttt{ProtectedRoute.test.tsx}: ADMIN ammesso, OWNER ammesso,
RECEPTIONIST rediretto a \texttt{/}, utente \texttt{null} con \texttt{allowedRoles} impostato
rediretto a \texttt{/}. Tutti i 144 test rimangono verdi; lint zero warnings; TypeScript clean.

\paragraph{Riferimento} Commit: \commitref{ccfa431}{fix(security): T-FE-03 --- role-based route guard on ProtectedRoute}

\paragraph{Estensione --- Pattern ``role-gated service call'' (T-FE-03b)}
Una variante del medesimo pattern è stata individuata nel \texttt{dashboardService}:
la funzione \texttt{getDashboardStats()} invocava incondizionatamente l'endpoint
\texttt{/api/v1/reports/owner} (protetto da \texttt{@PreAuthorize("hasAnyRole('OWNER','ADMIN')")}),
causando HTTP~403 per i ruoli \texttt{RECEPTIONIST} e \texttt{MANAGER}.
La risposta di errore propagava all'intero Zustand store, rendendo la dashboard
inutilizzabile per tutti i ruoli non privilegiati.
A differenza della vulnerabilità T-FE-03 (accesso diretto alla route), qui la chiamata
API non autorizzata avveniva in modo silenzioso, senza che l'utente potesse accorgersene.
Il backend era correttamente protetto e non restituiva dati sensibili, ma il design
configurava comunque una violazione OWASP~A01.
Fix: il parametro \texttt{isOwnerOrAdmin: boolean} è propagato da \texttt{dashboardStore.fetchStats()}
a \texttt{Dashboard.tsx}; la chiamata al billing report è omessa se \texttt{false},
\texttt{pendingRevenue} ritorna \texttt{null}.
Commit: \commitref{14a43bb}{fix(dashboard): gate owner report behind role + replace stats with daily KPIs}

\subsection{Protezione Actuator Endpoints \threat{T-CFG-01}}

\paragraph{Descrizione}
Spring Boot Actuator espone endpoint di diagnostica e monitoraggio (\texttt{/actuator/health},
\texttt{/actuator/prometheus}, \texttt{/actuator/info}) sulla stessa porta dell'API applicativa.
Nella baseline, tutti e 9 i microservizi pubblicavano la propria porta applicativa
(da \texttt{8080} a \texttt{8087} e \texttt{8888}) con questi endpoint liberamente accessibili.
In particolare:

\begin{itemize}
  \item \texttt{/actuator/prometheus} espone metriche interne del processo JVM (thread pool,
        connection pool del database, latenze HTTP, circuit breaker state) — informazioni
        preziose per un attaccante in fase di reconnaissance.
  \item \texttt{/actuator/health} con \texttt{show-details: always} rivela lo stato dei
        sotto-sistemi (database UP/DOWN, Redis ping, Zipkin) e i dettagli delle loro
        configurazioni di connessione.
\end{itemize}

Classificato OWASP A02:2025 --- Security Misconfiguration.

\begin{insecurebox}[title={Configurazione Vulnerabile --- auth-service.yml (baseline)}]
\begin{lstlisting}[style=yaml]
# INSICURO: actuator sulla stessa porta dell'API (8087), accessibile
# a chiunque raggiunga la porta pubblicata nel docker-compose
management:
  endpoints:
    web:
      exposure:
        include: "prometheus,health,info"
  endpoint:
    health:
      show-details: always   # dettagli DB/Redis esposti senza autenticazione
\end{lstlisting}
\end{insecurebox}

\begin{insecurebox}[title={docker-compose.yml (baseline) --- porta applicativa pubblicata}]
\begin{lstlisting}[style=yaml]
auth-service:
  ports:
    - "8087:8087"   # /actuator/prometheus raggiungibile da qualsiasi client
\end{lstlisting}
\end{insecurebox}

\paragraph{Mitigazione applicata}
La mitigazione separa il piano di gestione dal piano applicativo tramite la proprietà
\texttt{management.server.port: 8090}. Tutti e 9 i servizi (incluso il config-service)
legano ora il management server sulla porta \texttt{8090}, che \textbf{non viene mai
pubblicata} nel \texttt{docker-compose.yml}.

Il risultato è il seguente modello di accesso:

\begin{itemize}
  \item \textbf{Internet / client esterno} → \texttt{:808x} (API) — nessuna rotta actuator.
  \item \textbf{Prometheus / healthcheck Docker} → \texttt{:8090} (management) — raggiungibile
        solo dall'interno della rete Docker \texttt{hotel\_network}.
  \item \texttt{show-details} impostato a \texttt{when-authorized}: gli endpoint health
        restituiscono solo \texttt{"status": "UP"} a chiamanti anonimi.
\end{itemize}

\begin{securebox}[title={Configurazione Sicura --- auth-service.yml (feature/secure-coding-hardening)}]
\begin{lstlisting}[style=yaml]
# SICURO: management server separato su porta interna non pubblicata
management:
  server:
    port: 8090             # porta dedicata, non nel blocco "ports" del compose
  endpoints:
    web:
      exposure:
        include: "prometheus,health,info"
  endpoint:
    health:
      probes:
        enabled: true
      show-details: when-authorized  # dettagli visibili solo ad autenticati
\end{lstlisting}
\end{securebox}

\begin{securebox}[title={docker-compose.yml (post-hardening) --- porta management non pubblicata}]
\begin{lstlisting}[style=yaml]
auth-service:
  ports:
    - "8087:8087"   # solo API esposta; porta 8090 non in "ports"
  healthcheck:
    test: ["CMD-SHELL",
      "wget -qO- http://localhost:8090/actuator/health/liveness || exit 1"]
\end{lstlisting}
\end{securebox}

\begin{securebox}[title={docker/prometheus/prometheus.yml --- scraping su porta interna}]
\begin{lstlisting}[style=yaml]
scrape_configs:
  - job_name: "auth-service"
    metrics_path: /actuator/prometheus
    static_configs:
      - targets: ["auth-service:8090"]  # rete Docker interna, non pubblica
\end{lstlisting}
\end{securebox}

La stessa configurazione è applicata simmetricamente a tutti e 9 i servizi
(\texttt{api-gateway}, \texttt{auth-service}, \texttt{guest-service},\\
\texttt{inventory-service}, \texttt{reservation-service}, \texttt{stay-service},\\
\texttt{billing-service}, \texttt{fb-service}, \texttt{config-service}).

\paragraph{Riferimento} Commit:\\
\commitref{b0cf898}{fix(security): T-CFG-01 --- move actuator to internal-only management port}

% ─────────────────────────────────────────────────────────────────────────────
\subsection{Rate Limiting — Estensione a Tutte le Route e Key Resolver Proxy-Aware \threat{T-GW-02}}
% ─────────────────────────────────────────────────────────────────────────────

\paragraph{Descrizione}
Il rate limiting baseline presentava due vulnerabilità distinte che lo rendevano aggirabile
in produzione:

\begin{enumerate}
  \item \textbf{Copertura parziale}: il filtro \texttt{RequestRateLimiter} era configurato
        \emph{unicamente} sulla route \texttt{auth-service} (\texttt{/api/v1/auth/**}).
        Le sei route business (guest, inventory, reservation, stay, billing, fb) erano
        completamente prive di rate limiting.
        Un attaccante autenticato — o uno script che sfrutti un account compromesso —
        poteva inondare qualsiasi endpoint business senza alcun limite, causando un
        denial-of-service verso gli altri tenant del sistema multi-hotel.

  \item \textbf{Key resolver non proxy-aware}: il \texttt{remoteAddrKeyResolver} baseline
        leggeva il campo TCP \texttt{RemoteAddress} dell'exchange.
        In un deployment con reverse proxy o load-balancer (nginx, AWS ALB), questo campo
        contiene l'IP del proxy e \emph{non} quello del client originale: tutti i client
        condividono quindi un unico bucket comune, rendendo il rate limiting di fatto
        inutile (un singolo attaccante esaurisce il budget dell'intero sistema).
\end{enumerate}

Classificato OWASP A02:2025 --- Security Misconfiguration.

\begin{insecurebox}[title={Baseline vulnerabile — rate limiting solo su auth, key resolver non proxy-aware}]
\begin{lstlisting}[style=yaml]
# INSICURO 1: le route business non hanno RequestRateLimiter
- id: guest-service
  uri: "http://guest-service:8083"
  predicates:
    - Path=/api/v1/guests/**
  filters:
    - AuthenticationFilter   # nessun rate limiter -> flooding illimitato

# INSICURO 2: key resolver usa TCP RemoteAddress; dietro un proxy
# tutti i client condividono il medesimo bucket (IP del proxy)
@Bean
public KeyResolver remoteAddrKeyResolver() {
    return exchange -> Mono.just(
        Objects.requireNonNull(
            exchange.getRequest().getRemoteAddress(),  // IP del proxy, non del client!
            "Remote address must not be null")
            .getAddress().getHostAddress());
}
\end{lstlisting}
\end{insecurebox}

\paragraph{Mitigazione applicata}

\subparagraph{Key resolver proxy-aware}
Il \texttt{remoteAddrKeyResolver} è stato aggiornato per leggere il header
\texttt{X-Forwarded-For} quando presente, estraendo il primo IP della lista
(il client originale secondo la specifica RFC 7239).  Solo in assenza del header
si ricade sull'\texttt{RemoteAddress} TCP.

È stato inoltre introdotto un secondo bean, \texttt{userKeyResolver}, dedicato alle
route autenticate.  Poiché il filtro \texttt{AuthenticationFilter} è dichiarato
\emph{prima} di \texttt{RequestRateLimiter} nella catena di ogni route, quando il rate
limiter viene invocato l'header \texttt{X-Auth-User} è già stato iniettato nel request
mutato.  Il resolver legge tale header e assegna a ciascun utente autenticato un bucket
indipendente (chiave \texttt{"user:<username>"}), cadendo indietro sull'IP se l'header
mancasse.  Il risultato è che:
\begin{itemize}
  \item un account compromesso non può esaurire il quota di altri account (isolamento
        per-utente);
  \item le richieste pre-autenticazione (es. endpoint non protetti o fallback) continuano
        ad essere limitate per IP.
\end{itemize}

\subparagraph{Estensione a tutte le route}
Il filtro \texttt{RequestRateLimiter} è stato aggiunto a tutte e sei le route business,
con soglie calibrate sul carico atteso di un operatore alberghiero:
\begin{itemize}
  \item Route autenticate: \texttt{replenishRate: 20}, \texttt{burstCapacity: 50}
        (per-user via \texttt{userKeyResolver}).
  \item Route auth: \texttt{replenishRate: 5}, \texttt{burstCapacity: 10}\\
        (per-IP via \texttt{remote\-Addr\-Key\-Resolver}).
\end{itemize}

\begin{securebox}[title={RateLimiterConfig.java — key resolver proxy-aware e per-utente (api-gateway)}]
\begin{lstlisting}[style=java]
@Bean
public KeyResolver remoteAddrKeyResolver() {
    return exchange -> {
        final String forwarded =
            exchange.getRequest().getHeaders().getFirst("X-Forwarded-For");
        if (forwarded != null && !forwarded.isBlank()) {
            return Mono.just(forwarded.split(",")[0].trim()); // client originale
        }
        return Mono.just(
            Objects.requireNonNull(
                exchange.getRequest().getRemoteAddress(),
                "Remote address must not be null").getAddress().getHostAddress());
    };
}

@Bean
public KeyResolver userKeyResolver() {
    return exchange -> {
        final String user =
            exchange.getRequest().getHeaders().getFirst("X-Auth-User");
        if (user != null && !user.isBlank()) {
            return Mono.just("user:" + user);   // bucket per-utente
        }
        final String forwarded =
            exchange.getRequest().getHeaders().getFirst("X-Forwarded-For");
        if (forwarded != null && !forwarded.isBlank()) {
            return Mono.just("ip:" + forwarded.split(",")[0].trim());
        }
        return Mono.just("ip:" + Objects.requireNonNull(
            exchange.getRequest().getRemoteAddress(),
            "Remote address must not be null").getAddress().getHostAddress());
    };
}
\end{lstlisting}
\end{securebox}

\begin{securebox}[title={api-gateway.yml — rate limiting su route business autenticate}]
\begin{lstlisting}[style=yaml]
- id: guest-service
  uri: "${GW_GUEST_SERVICE_URI:http://guest-service:8083}"
  predicates:
    - Path=/api/v1/guests/**
  filters:
    - AuthenticationFilter         # 1. valida JWT, inietta X-Auth-User
    - name: RequestRateLimiter     # 2. usa X-Auth-User per bucket per-utente
      args:
        redis-rate-limiter.replenishRate: 20
        redis-rate-limiter.burstCapacity: 50
        redis-rate-limiter.requestedTokens: 1
        key-resolver: "#{@userKeyResolver}"
# Identica configurazione applicata a: inventory, reservation,
# stay, billing, fb-service
\end{lstlisting}
\end{securebox}

\paragraph{Test} Aggiunti 6 test in \texttt{RateLimiterConfigTest}:
\texttt{remoteAddrKeyResolver} — XFF presente $\Rightarrow$ primo IP; XFF assente $\Rightarrow$ RemoteAddress; XFF blank $\Rightarrow$ RemoteAddress.
\texttt{userKeyResolver} — \texttt{X-Auth-User} presente $\Rightarrow$ \texttt{"user:<username>"}; assente + XFF presente $\Rightarrow$ \texttt{"ip:<xff>"}; entrambi assenti $\Rightarrow$ \texttt{"ip:<remoteAddr>"}.

\paragraph{Riferimento} Commit: \commitref{25c7b57}{fix(security): T-GW-02 --- rate limiting on all routes with proxy-aware key resolvers}

% ─────────────────────────────────────────────────────────────────────────────
\subsection{CORS Misconfiguration --- Whitelist Header Esplicita \threat{T-GW-04}}
% ─────────────────────────────────────────────────────────────────────────────

\paragraph{Descrizione}
La politica CORS (Cross-Origin Resource Sharing) dell'API Gateway consente al browser
di inviare richieste autenticate dall'applicazione React (porta 5173 in sviluppo) verso
il gateway (porta 8080).  Poiché \texttt{allowCredentials: true} è obbligatorio (le
richieste portano il cookie httpOnly JWT), qualsiasi permissività sugli header diventa
rischiosa: un header non necessario accettato in preflight potrebbe essere sfruttato
da uno script iniettato per inviare informazioni fuori banda verso l'API.

Nella baseline la configurazione YAML dichiarava:

\begin{lstlisting}[style=yaml]
allowedHeaders: "*"
\end{lstlisting}

Il valore wildcard istruisce il gateway a riflettere \emph{qualsiasi} header nella
risposta \texttt{Access-Control-Allow-Headers} della preflight OPTIONS, indipendentemente
da ciò che il server abbia effettivamente bisogno di leggere.
In abbinamento a \texttt{allowCredentials: true}, questo amplia la superficie di
attacco: script di terze parti iniettati (XSS) possono includere header personalizzati
arbitrari — ad esempio \texttt{X-Custom-Attack} — che il browser invierà con le credenziali.
Classificato OWASP A02:2025 --- Security Misconfiguration.

Un secondo problema minore riguardava l'assenza di \texttt{maxAge}: senza questa
proprietà, il browser non può mettere in cache le risposte alle preflight OPTIONS,
generando una richiesta OPTIONS aggiuntiva prima di \emph{ogni} richiesta mutante.
Nella pratica d'uso del gestionale (centinaia di operazioni per sessione) questo
moltiplica inutilmente le richieste al gateway.

\begin{insecurebox}[title={Configurazione Vulnerabile --- api-gateway.yml (baseline, cors globalcors)}]
\begin{lstlisting}[style=yaml]
spring:
  cloud:
    gateway:
      globalcors:
        cors-configurations:
          '[/**]':
            allowedOrigins:
              "${GW_CORS_ALLOWED_ORIGINS:http://localhost:5173}"
            allowedMethods: [GET, POST, PUT, PATCH, DELETE, OPTIONS]
            allowedHeaders: "*"   # wildcard: qualsiasi header e' riflesso
            allowCredentials: true
            # maxAge assente: nessuna cache delle preflight
\end{lstlisting}
\end{insecurebox}

\paragraph{Mitigazione applicata}
La whitelist degli header accettati è stata sostituita con una lista esplicita e minimale:

\begin{itemize}
  \item \textbf{\texttt{Content-Type}}: necessario per il body JSON delle richieste
        POST/PUT/PATCH.
  \item \textbf{\texttt{X-CSRF-Token}}: necessario per il meccanismo di protezione
        Double-Submit-Cookie introdotto in T-GW-05; senza questa voce, il browser
        non otterrebbe l'approvazione della preflight per l'header e le richieste
        mutanti fallirebbero in produzione.
\end{itemize}

L'header \texttt{Authorization} non è necessario: il token JWT transita nel cookie
httpOnly, mai in un header esplicito.

È stato aggiunto \texttt{maxAge: 3600} per istruire il browser a memorizzare
in cache il risultato della preflight OPTIONS per un'ora, riducendo il numero
di richieste ridondanti senza abbassare la sicurezza.

L'origine consentita era già correttamente vincolata tramite la variabile
d'ambiente \texttt{GW\_CORS\_ALLOWED\_ORIGINS} (default \texttt{http://localhost:5173}),
senza l'uso di wildcard né pattern globbing. In produzione, nginx serve il frontend
sulla stessa origine del proxy API, rendendo il CORS non necessario in quel contesto.

\begin{securebox}[title={Configurazione Sicura --- api-gateway.yml (feature/secure-coding-hardening)}]
\begin{lstlisting}[style=yaml]
spring:
  cloud:
    gateway:
      globalcors:
        cors-configurations:
          '[/**]':
            allowedOrigins:
              "${GW_CORS_ALLOWED_ORIGINS:http://localhost:5173}"
            allowedMethods: [GET, POST, PUT, PATCH, DELETE, OPTIONS]
            allowedHeaders:      # whitelist minimale e dichiarativa
              - Content-Type     # body JSON richieste mutanti
              - X-CSRF-Token     # Double-Submit-Cookie (T-GW-05)
            allowCredentials: true
            maxAge: 3600         # cache preflight: 1 ora
\end{lstlisting}
\end{securebox}

\paragraph{Effetto}
Un client cross-origin può includere nei propri header \emph{solo} \texttt{Content-Type}
e \texttt{X-CSRF-Token}; qualsiasi altro header personalizzato viene rifiutato già in
fase di preflight dal browser, prima che la richiesta raggiunga il gateway.
La combinazione con \texttt{allowedOrigins} esplicito e \texttt{allowCredentials: true}
forma un contratto CORS minimale e corretto: accetta solo ciò che è strettamente
necessario al funzionamento dell'applicazione.

\paragraph{Nota sull'assenza di test unitari}
La configurazione CORS è gestita interamente da Spring Cloud Gateway tramite il
blocco YAML \texttt{globalcors}, senza alcun bean Java da testare in isolamento.
Il comportamento corretto è verificato dall'integrazione end-to-end (Playwright E2E):
il browser esegue le preflight verso \texttt{http://localhost:5173} e, se la risposta
OPTIONS è errata, l'intera suite E2E fallisce a livello di \texttt{NetworkError}.

\paragraph{Riferimento} Commit: \commitref{73e2788}{fix(security): T-GW-04 --- restrict CORS allowedHeaders whitelist and add preflight maxAge}

% ═════════════════════════════════════════════════════════════════════════════
\section{Broken Access Control --- Isolamento Multi-Tenant}
% ═════════════════════════════════════════════════════════════════════════════

% ─────────────────────────────────────────────────────────────────────────────
\subsection{IDOR e Data Leak sul Guest Service \threat{T-GST-01} \threat{T-GST-03}}
% ─────────────────────────────────────────────────────────────────────────────

\paragraph{Descrizione}
Il sistema gestisce strutture ricettive multiple (multi-hotel). Ogni utente autenticato
appartiene a un hotel specifico, identificato dal campo \texttt{hotel\_id} estratto dal JWT
e iniettato dall'API Gateway come header \texttt{X-Auth-Hotel}.

\textbf{T-GST-01 (IDOR):} il \texttt{GuestController} esponeva endpoint
\path{GET /api/v1/guests/{id}} e \path{PUT /api/v1/guests/{id}} senza verificare
l'ownership dell'hotel. Un receptionist dell'Hotel A, conoscendo o indovinando un UUID,
poteva leggere e modificare profili ospite dell'Hotel B. Classificato OWASP A01 ---
Broken Access Control.

\textbf{T-GST-03 (Multi-tenant leak):} le query \texttt{findAll()} e \texttt{searchByKeyword()}
non filtravano per \texttt{hotel\_id}, restituendo ospiti di tutti gli hotel alla prima
chiamata paginata. I dati PII (nome, documento, data di nascita) risultavano accessibili
cross-hotel. Classificato OWASP A01 --- Broken Access Control.

\begin{insecurebox}[title={Codice Vulnerabile --- GuestServiceImpl.java (baseline)}]
\begin{lstlisting}[style=java]
// INSICURO: lookup per solo UUID, nessun check ownership
private Guest resolveGuest(UUID id) {
    return guestRepository.findById(id)   // trova chiunque, qualsiasi hotel
        .orElseThrow(() -> new NotFoundException("GUEST_NOT_FOUND"));
}

// INSICURO: lista tutti gli ospiti di tutti gli hotel
public Page<GuestResponse> getAllGuests(Pageable pageable) {
    return guestRepository.findAll(pageable)  // nessun filtro hotel_id
        .map(guestMapper::toResponse);
}
\end{lstlisting}
\end{insecurebox}

\paragraph{Mitigazione applicata}
L'intervento si articola in quattro livelli:

\begin{enumerate}
    \item \textbf{Schema DB} --- Flyway V2 aggiunge la colonna \texttt{hotel\_id UUID NOT NULL}
          alla tabella \texttt{guests}, con indice parziale su \texttt{active = TRUE}.
    \item \textbf{Extraction dal SecurityContext} --- \texttt{InternalAuthFilter} legge
          \texttt{X-Auth-Hotel} e lo memorizza in \path{Authentication.details}. Il service
          lo estrae con \texttt{extractHotelId()} senza dipendere dal controller.
    \item \textbf{Query scoped} --- tutte le query del repository usano varianti
          hotel-scoped (\texttt{findByIdAndHotelId}, \texttt{findAllByHotelId},
          \texttt{searchByKeywordAndHotelId}, \texttt{findAllByIdInAndHotelId}).
    \item \textbf{IDOR-safe 404} --- \texttt{resolveGuest(id, hotelId)} chiama
          \texttt{findByIdAndHotelId}: se il guest esiste ma appartiene a un altro hotel
          il metodo restituisce \texttt{Optional.empty()} e il service lancia
          \texttt{NotFoundException} (HTTP 404), senza rivelare l'esistenza della risorsa.
\end{enumerate}

\begin{securebox}[title={Codice Sicuro --- GuestServiceImpl.java (fix)}]
\begin{lstlisting}[style=java]
// SICURO: lookup combinato id + hotel_id, IDOR-safe con 404
private Guest resolveGuest(UUID id, UUID hotelId) {
    return guestRepository.findByIdAndHotelId(id, hotelId)
        .orElseThrow(() -> new NotFoundException("GUEST_NOT_FOUND"));
    // 404 anche se il guest esiste ma appartiene ad altro hotel:
    // l'attaccante non puo' distinguere "non esiste" da "non tuo"
}

// SICURO: elenco filtrato per hotel del caller
public Page<GuestResponse> getAllGuests(Pageable pageable) {
    UUID hotelId = extractHotelId();          // da Authentication.details
    return guestRepository
        .findAllByHotelId(hotelId, pageable)  // solo il proprio hotel
        .map(guestMapper::toResponse);
}

// SICURO: hotelId estratto dal SecurityContext (non dal client)
private UUID extractHotelId() {
    Authentication auth = SecurityContextHolder.getContext().getAuthentication();
    Object details = auth.getDetails();       // iniettato da InternalAuthFilter
    return UUID.fromString((String) details); // mai controllabile dal client
}
\end{lstlisting}
\end{securebox}

\begin{securebox}[title={InternalAuthFilter.java --- estrazione X-Auth-Hotel}]
\begin{lstlisting}[style=java]
// Il gateway inietta X-Auth-Hotel dal JWT; il filter lo memorizza nei details
String hotelId = request.getHeader("X-Auth-Hotel");
// Richiesta rifiutata (401) se l'header e' assente
if (!StringUtils.hasText(hotelId)) {
    rejectRequest(response, "Missing required gateway authentication headers");
    return;
}
UsernamePasswordAuthenticationToken auth = new UsernamePasswordAuthenticationToken(
    username, "", List.of(new SimpleGrantedAuthority("ROLE_" + role)));
auth.setDetails(hotelId);   // tenant context propagato al service layer
SecurityContextHolder.getContext().setAuthentication(auth);
\end{lstlisting}
\end{securebox}

\begin{securebox}[title={V2\_\_add\_hotel\_id\_to\_guests.sql --- Flyway migration}]
\begin{lstlisting}[style=sql]
ALTER TABLE guests
    ADD COLUMN hotel_id UUID NOT NULL
    DEFAULT '00000000-0000-0000-0000-000000000000';
ALTER TABLE guests ALTER COLUMN hotel_id DROP DEFAULT;

CREATE INDEX IF NOT EXISTS idx_guests_hotel_id
    ON guests (hotel_id) WHERE active = TRUE;
\end{lstlisting}
\end{securebox}

\paragraph{Riferimento} Commit:\\
\commitref{b483eac}{fix(security): T-GST-01 + T-GST-03 --- hotel\_id isolation on guest-service}

% ─────────────────────────────────────────────────────────────────────────────
\subsection{IDOR sul Reservation Service \threat{T-RES-02}}
% ─────────────────────────────────────────────────────────────────────────────

\paragraph{Descrizione}
Il \texttt{reservation-service} gestisce le prenotazioni di camera. Prima della
mitigazione, l'endpoint \texttt{GET /api/v1/reservations/\{id\}} e le operazioni
correlate (PUT, DELETE, PATCH) non verificavano che la prenotazione richiesta
appartenesse all'hotel del caller. Un receptionist dell'Hotel A, conoscendo o
indovinando un UUID valido, poteva leggere, modificare o cancellare prenotazioni
dell'Hotel B. Classificato OWASP A01 --- Broken Access Control.

Analogamente, \texttt{getAllReservations()} eseguiva \texttt{findAll()} senza
alcun filtro su \texttt{hotel\_id}, esponendo l'intero dataset di prenotazioni a
qualsiasi utente autenticato.

\begin{insecurebox}[title={Codice Vulnerabile --- ReservationServiceImpl.java (baseline)}]
\begin{lstlisting}[style=java]
// INSICURO: lookup per solo UUID, nessun controllo ownership hotel
private Reservation findReservationByIdOrThrow(UUID id) {
    return reservationRepository.findById(id)  // restituisce qualsiasi hotel
        .orElseThrow(() -> new NotFoundException("RESERVATION_NOT_FOUND"));
}

// INSICURO: lista tutte le prenotazioni di tutti gli hotel
public Page<ReservationResponse> getAllReservations(Pageable pageable) {
    return reservationRepository.findAll(pageable)  // nessun filtro hotel_id
        .map(...);
}
\end{lstlisting}
\end{insecurebox}

\paragraph{Mitigazione applicata}
L'intervento replica il pattern gia' applicato al \texttt{guest-service} (T-GST-01/T-GST-03)
e si articola in tre livelli:

\begin{enumerate}
    \item \textbf{InternalAuthFilter aggiornato} --- il filter legge l'header
          \texttt{X-Auth-Hotel} iniettato dal gateway e lo salva in
          \texttt{Authentication.details}. Richieste prive dell'header sono
          rifiutate con HTTP 401.
    \item \textbf{Query hotel-scoped} --- il repository espone due nuovi metodi:\\
          \texttt{\seqsplit{findByIdAndHotelId(id, hotelId)}} e
          \texttt{\seqsplit{findAllByHotelId(hotelId, pageable)}}. Il service usa
          \texttt{findByIdAndHotelId} per tutte le operazioni su singola
          risorsa (GET, PUT, DELETE, PATCH).
    \item \textbf{IDOR-safe 404} --- se la prenotazione esiste ma appartiene a
          un altro hotel, \texttt{findByIdAndHotelId} restituisce
          \texttt{Optional.empty()} e il service lancia \texttt{NotFoundException}
          (HTTP 404), senza rivelare l'esistenza della risorsa all'attaccante.
    \item \textbf{hotelId impostato alla creazione} --- \texttt{createReservation}
          recupera il \texttt{hotelId} dal SecurityContext tramite
          \texttt{resolveHotelId()} e lo persiste sull'entità, garantendo che
          ogni nuova prenotazione sia legata all'hotel del caller.
\end{enumerate}

\begin{securebox}[title={Codice Sicuro --- ReservationServiceImpl.java (fix)}]
\begin{lstlisting}[style=java]
// SICURO: lookup combinato id + hotel_id (IDOR-safe con 404)
private Reservation findReservationByIdAndHotelOrThrow(UUID id, UUID hotelId) {
    return reservationRepository.findByIdAndHotelId(id, hotelId)
        .orElseThrow(() -> new NotFoundException("RESERVATION_NOT_FOUND"));
    // 404 anche se esiste ma appartiene ad altro hotel
}

// SICURO: lista filtrata per hotel del caller
public Page<ReservationResponse> getAllReservations(Pageable pageable) {
    UUID hotelId = resolveHotelId();
    return reservationRepository
        .findAllByHotelId(hotelId, safePageable)
        .map(...);
}

// SICURO: hotelId estratto dal SecurityContext, mai dal client
private UUID resolveHotelId() {
    Authentication auth = SecurityContextHolder.getContext().getAuthentication();
    if (auth == null || !(auth.getDetails() instanceof String hotelIdStr)) {
        throw new IllegalStateException("Hotel ID cannot be null");
    }
    return UUID.fromString(hotelIdStr);
}

// SICURO: hotelId persistito sul record alla creazione
public ReservationResponse createReservation(ReservationRequest request) {
    UUID hotelId = resolveHotelId();
    Reservation reservation = reservationMapper.toEntity(request);
    reservation.setHotelId(hotelId);  // tenant-bound, non alterabile dal client
    ...
}
\end{lstlisting}
\end{securebox}

\begin{securebox}[title={ReservationRepository.java --- metodi hotel-scoped}]
\begin{lstlisting}[style=java]
// Lookup singola risorsa con check ownership (Spring Data query derivata)
Optional<Reservation> findByIdAndHotelId(UUID id, UUID hotelId);

// Lista paginata filtrata per hotel
Page<Reservation> findAllByHotelId(UUID hotelId, Pageable pageable);
\end{lstlisting}
\end{securebox}

\paragraph{Riferimento} Commit: \commitref{3e93f49}{fix(security): T-RES-02 --- hotel\_id isolation on reservation-service}

% ─────────────────────────────────────────────────────────────────────────────
\subsection{IDOR sul Billing Service \threat{T-BILL-01}}
% ─────────────────────────────────────────────────────────────────────────────

\paragraph{Descrizione}
Il \texttt{billing-service} gestisce fatture e pagamenti. Prima della mitigazione,
\texttt{InvoiceServiceImpl.getInvoice(id)} e \texttt{PaymentServiceImpl.addPayment(invoiceId, ...)}
usavano \texttt{findById(id)}, che restituisce qualsiasi fattura indipendentemente dall'hotel
di appartenenza. Un receptionist dell'Hotel A, indovinando un UUID valido, poteva:
\begin{itemize}
  \item leggere le fatture dell'Hotel B (\texttt{GET /api/v1/invoices/\{id\}});
  \item registrare pagamenti su fatture di altri hotel (\texttt{POST}\linebreak\texttt{/api/v1/invoices/\{invoiceId\}/payments});
  \item ottenere la lista completa di tutte le fatture di tutti gli hotel via \texttt{getAllInvoices()}.
\end{itemize}
Classificato OWASP A01 --- Broken Access Control.

\begin{insecurebox}[title={Codice Vulnerabile --- InvoiceServiceImpl.java (baseline)}]
\begin{lstlisting}[style=java]
// INSICURO: lookup per solo UUID, nessun controllo ownership hotel
public InvoiceResponse getInvoice(UUID id) {
    Invoice invoice = invoiceRepository.findById(id)  // qualsiasi hotel
        .orElseThrow(() -> new NotFoundException("INVOICE_NOT_FOUND"));
    return invoiceMapper.toResponse(invoice);
}

// INSICURO: lista tutte le fatture di tutti gli hotel
public Page<InvoiceResponse> getAllInvoices(Pageable pageable) {
    return invoiceRepository.findAll(pageable)  // nessun filtro hotel_id
        .map(invoiceMapper::toResponse);
}
\end{lstlisting}
\end{insecurebox}

\begin{insecurebox}[title={Codice Vulnerabile --- PaymentServiceImpl.java (baseline)}]
\begin{lstlisting}[style=java]
// INSICURO: addPayment carica la fattura senza verifica hotel
public PaymentResponse addPayment(UUID invoiceId, PaymentRequest request) {
    Invoice invoice = invoiceRepository.findById(invoiceId)
        .orElseThrow(() -> new NotFoundException("INVOICE_NOT_FOUND"));
    // ... aggiunge il pagamento indipendentemente dall'hotel
}
\end{lstlisting}
\end{insecurebox}

\paragraph{Mitigazione applicata}
L'intervento replica il pattern consolidato di T-GST-01 e T-RES-02:

\begin{enumerate}
    \item \textbf{InternalAuthFilter aggiornato} --- il filter legge l'header
          \texttt{X-Auth-Hotel} iniettato dal gateway e lo salva in
          \texttt{Authentication.details}. Richieste prive dell'header sono
          rifiutate con HTTP 401.
    \item \textbf{Query hotel-scoped nel repository} --- aggiunti tre metodi:\\
          \texttt{findByIdAndHotelId(id, hotelId)},
          \texttt{findByHotelId(hotelId, pageable)} (paginata) e\\
          \texttt{\seqsplit{findFirstByReservationIdAndHotelIdOrderByIssueDateDesc(reservationId,hotelId)}}.
    \item \textbf{IDOR-safe 404} --- \texttt{findByIdAndHotelId} restituisce
          \texttt{Optional.empty()} se la fattura esiste ma appartiene ad
          un altro hotel; il service lancia \texttt{NotFoundException} (HTTP 404),
          senza rivelare l'esistenza della risorsa all'attaccante.
    \item \textbf{hotelId impostato alla creazione} --- \texttt{createInvoice} imposta\\
          \texttt{\seqsplit{invoice.setHotelId(resolveHotelId())}}, ignorando qualsiasi valore
          passato nel corpo della richiesta.
    \item \textbf{FeignHeaderConfig aggiornato} --- propaga \texttt{X-Auth-Hotel} su
          tutte le chiamate Feign verso \texttt{guest-service} e
          \texttt{reservation-service}.
\end{enumerate}

\begin{securebox}[title={Codice Sicuro --- InvoiceServiceImpl.java (fix)}]
\begin{lstlisting}[style=java]
// SICURO: lookup combinato id + hotel_id (IDOR-safe con 404)
public InvoiceResponse getInvoice(UUID id) {
    UUID hotelId = resolveHotelId();
    Invoice invoice = invoiceRepository.findByIdAndHotelId(id, hotelId)
        .orElseThrow(() -> new NotFoundException("INVOICE_NOT_FOUND"));
    return invoiceMapper.toResponse(invoice);
}

// SICURO: lista paginata filtrata per hotel del caller
public Page<InvoiceResponse> getAllInvoices(Pageable pageable) {
    UUID hotelId = resolveHotelId();
    return invoiceRepository.findByHotelId(hotelId, safePageable)
        .map(invoiceMapper::toResponse);
}

// SICURO: hotelId estratto dal SecurityContext, mai dal client
public InvoiceResponse createInvoice(InvoiceRequest request) {
    // ... validazioni guest e reservation ...
    Invoice invoice = invoiceMapper.toEntity(request);
    invoice.setHotelId(resolveHotelId());  // tenant-bound
    return invoiceMapper.toResponse(invoiceRepository.save(invoice));
}

private UUID resolveHotelId() {
    Authentication auth = SecurityContextHolder.getContext().getAuthentication();
    if (auth == null || !(auth.getDetails() instanceof String hotelIdStr)) {
        throw new IllegalStateException("MISSING_HOTEL_CONTEXT");
    }
    return UUID.fromString(hotelIdStr);
}
\end{lstlisting}
\end{securebox}

\begin{securebox}[title={Codice Sicuro --- PaymentServiceImpl.java (fix)}]
\begin{lstlisting}[style=java]
// SICURO: lookup combinato id + hotel_id (IDOR-safe con 404)
public PaymentResponse addPayment(UUID invoiceId, PaymentRequest request) {
    UUID hotelId = resolveHotelId();
    Invoice invoice = invoiceRepository.findByIdAndHotelId(invoiceId, hotelId)
        .orElseThrow(() -> new NotFoundException("INVOICE_NOT_FOUND"));
    // ... aggiunge il pagamento solo se la fattura appartiene al caller
}
\end{lstlisting}
\end{securebox}

\begin{securebox}[title={InvoiceRepository.java --- metodi hotel-scoped}]
\begin{lstlisting}[style=java]
// Lookup singola fattura con check ownership (Spring Data query derivata)
Optional<Invoice> findByIdAndHotelId(UUID id, UUID hotelId);

// Lista paginata filtrata per hotel
Page<Invoice> findByHotelId(UUID hotelId, Pageable pageable);

// Ultima fattura per prenotazione, scoped per hotel
Optional<Invoice> findFirstByReservationIdAndHotelIdOrderByIssueDateDesc(
    UUID reservationId, UUID hotelId);
\end{lstlisting}
\end{securebox}

\paragraph{Test}
Aggiunti/aggiornati test in \texttt{InvoiceServiceImplTest} e \texttt{PaymentServiceImplTest}:
il \texttt{SecurityContextHolder} viene inizializzato in \texttt{setUp()} con un
\texttt{\seqsplit{UsernamePasswordAuthenticationToken}} il cui \texttt{details} è il \texttt{hotelId} del
caller, e ripulito in \texttt{tearDown()}. Aggiunto il caso
\texttt{\seqsplit{shouldThrowWhenInvoiceBelongsToDifferentHotel}} che verifica il comportamento IDOR-safe
(la query hotel-scoped restituisce \texttt{Optional.empty()} → HTTP 404).

\paragraph{Riferimento} Commit: \commitref{4a44eea}{fix(security): T-BILL-01 --- IDOR fix on billing-service (hotel-scoped invoice/payment lookup)}

% ─────────────────────────────────────────────────────────────────────────────
\subsection{Data Leak Cross-Tenant sul Report Finanziario Owner \threat{T-BILL-04}}
% ─────────────────────────────────────────────────────────────────────────────

\paragraph{Descrizione}
Trovato durante il design di una regola ArchUnit per imporre lo scoping \texttt{hotel\_id}
su tutte le query multi-tenant (non da un pentest mirato) --- l'endpoint
\texttt{GET /api/v1/reports/owner}, protetto da \texttt{@PreAuthorize("hasAnyRole('OWNER','ADMIN')")},
non estraeva mai l'\texttt{hotelId} dal contesto di sicurezza. Il controller passava le
sole date al service, che interrogava \texttt{InvoiceRepository.findByIssueDateBetween(start, end)}
--- una query priva di qualunque parametro \texttt{hotel\_id}. Il report aggregava fatturato,
conteggio fatture e lista fatture (inclusi dati ospite) di \emph{tutti} gli hotel della
piattaforma. Non serviva alcuna enumerazione di UUID: era il comportamento di default
dell'endpoint per qualunque Owner/Admin autenticato, su un hotel qualsiasi. Il bug era
invisibile in produzione con un solo hotel attivo, ma immediatamente sfruttabile non
appena la piattaforma passa multi-tenant. Classificato OWASP A01 --- Broken Access Control.

\begin{insecurebox}[title={Codice Vulnerabile --- OwnerReportController/ServiceImpl.java (baseline)}]
\begin{lstlisting}[style=java]
// INSICURO: nessun hotelId estratto dal contesto autenticato
@GetMapping("/owner")
@PreAuthorize("hasAnyRole('OWNER', 'ADMIN')")
public ResponseEntity<OwnerFinancialReportDto> getOwnerFinancialReport(
        LocalDate startDate, LocalDate endDate) {
    OwnerFinancialReportDto report =
        ownerReportService.getFinancialReport(startDate, endDate);  // no hotel scope
    return ResponseEntity.ok(report);
}

// INSICURO: aggrega fatture di TUTTI gli hotel
public OwnerFinancialReportDto getFinancialReport(LocalDate startDate, LocalDate endDate) {
    List<Invoice> invoices = invoiceRepository.findByIssueDateBetween(start, end);
    // ... nessun filtro hotel_id ...
}
\end{lstlisting}
\end{insecurebox}

\paragraph{Mitigazione applicata}
Stesso pattern già consolidato per T-BILL-01/T-GST-01/T-RES-02, applicato a un endpoint
mai coperto in precedenza:

\begin{enumerate}
    \item \textbf{extractHotelId() nel controller} --- stesso helper già presente in
          \texttt{InvoiceController}, legge \texttt{Authentication.details} (popolato
          dall'header \texttt{X-Auth-Hotel} via \texttt{InternalAuthFilter}).
    \item \textbf{Nuovo metodo repository hotel-scoped} ---
          \texttt{findByHotelIdAndIssueDateBetween(hotelId, start, end)} sostituisce
          l'unscoped \texttt{findByIssueDateBetween}.
    \item \textbf{Firma del service estesa} --- \texttt{getFinancialReport(hotelId,
          startDate, endDate)}; l'\texttt{hotelId} arriva sempre dal contesto
          autenticato, mai dal client.
    \item \textbf{Pulizia correlata} --- rimosso anche \texttt{InvoiceRepository.findByGuestId}
          (dead code, zero chiamanti, stesso rischio strutturale se mai richiamato in futuro).
\end{enumerate}

\begin{securebox}[title={Codice Sicuro --- OwnerReportController/ServiceImpl.java (fix)}]
\begin{lstlisting}[style=java]
// SICURO: hotelId risolto dal SecurityContext, mai dal client
@GetMapping("/owner")
@PreAuthorize("hasAnyRole('OWNER', 'ADMIN')")
public ResponseEntity<OwnerFinancialReportDto> getOwnerFinancialReport(
        LocalDate startDate, LocalDate endDate) {
    UUID hotelId = extractHotelId();
    OwnerFinancialReportDto report =
        ownerReportService.getFinancialReport(hotelId, startDate, endDate);
    return ResponseEntity.ok(report);
}

// SICURO: aggrega solo le fatture dell'hotel del caller
public OwnerFinancialReportDto getFinancialReport(
        UUID hotelId, LocalDate startDate, LocalDate endDate) {
    List<Invoice> invoices =
        invoiceRepository.findByHotelIdAndIssueDateBetween(hotelId, start, end);
    // ...
}
\end{lstlisting}
\end{securebox}

\paragraph{Test}
Nuovo \texttt{shouldScopeQueryToTheAuthenticatedHotelOnly} in \texttt{OwnerReportServiceImplTest}
verifica che il repository sia sempre invocato con l'\texttt{hotelId} del caller.
\texttt{OwnerReportControllerTest} aggiornato con \texttt{SecurityContextHolder} popolato
(stesso pattern di T-BILL-01), asserzioni \texttt{verify(...)} sull'\texttt{hotelId} propagato.

\paragraph{Riferimento} Commit: \commitref{c983946}{fix(security): T-BILL-04 --- cross-tenant data leak on Owner financial report}

% ─────────────────────────────────────────────────────────────────────────────
\subsection{Data Leak Cross-Tenant sul Pre-fill Check-in \threat{T-STAY-06}}
% ─────────────────────────────────────────────────────────────────────────────

\paragraph{Descrizione}
Trovato nella stessa esplorazione di T-BILL-04. \texttt{StayServiceImpl.getLastCompletedStayForGuest(guestId)}
--- usato per pre-compilare il form di check-in con i dati dell'ultimo soggiorno di un
ospite di ritorno --- interrogava \texttt{StayRepository.findTopByGuestIdAndStatusOrderByActualCheckInTimeDesc(guestId,
CHECKED\_OUT)} senza alcun parametro \texttt{hotel\_id}. L'unico controllo presente era che
il profilo ospite fosse ancora attivo su guest-service, non che il soggiorno restituito
appartenesse all'hotel del chiamante. Un receptionist dell'Hotel A, per un
\texttt{guestId} noto, poteva ottenere numero camera, date e stato fattura/Alloggiati
dell'ultimo soggiorno di quell'ospite anche se avvenuto presso l'Hotel B. Classificato
OWASP A01 --- Broken Access Control.

\begin{insecurebox}[title={Codice Vulnerabile --- StayServiceImpl.java (baseline)}]
\begin{lstlisting}[style=java]
// INSICURO: nessun hotelId, solo verifica che il guest sia attivo
public Optional<StayResponse> getLastCompletedStayForGuest(UUID guestId) {
    GuestResponse guest = guestClient.getGuestById(guestId);  // solo attivo?
    return stayRepository
        .findTopByGuestIdAndStatusOrderByActualCheckInTimeDesc(guestId, CHECKED_OUT)
        .map(stayMapper::toDto);  // soggiorno di QUALSIASI hotel
}
\end{lstlisting}
\end{insecurebox}

\paragraph{Mitigazione applicata}
Stesso pattern di T-STAY-04/T-STAY-05: nuova query hotel-scoped
\texttt{findTopByGuestIdAndHotelIdAndStatusOrderByActualCheckInTimeDesc}; il service e il
controller richiedono ora \texttt{hotelId}, risolto da \texttt{StayController.extractHotelId()}
(mai dal client).

\begin{securebox}[title={Codice Sicuro --- StayServiceImpl.java (fix)}]
\begin{lstlisting}[style=java]
// SICURO: hotelId propagato dal controller, mai dal client
public Optional<StayResponse> getLastCompletedStayForGuest(UUID guestId, UUID hotelId) {
    GuestResponse guest = guestClient.getGuestById(guestId);
    return stayRepository
        .findTopByGuestIdAndHotelIdAndStatusOrderByActualCheckInTimeDesc(
            guestId, hotelId, CHECKED_OUT)
        .map(stayMapper::toDto);
}
\end{lstlisting}
\end{securebox}

\paragraph{Test}
Nuovi \texttt{shouldReturnLastCompletedStayScopedToTheAuthenticatedHotel} e
\texttt{shouldReturnEmptyWhenGuestServiceUnavailableForPreFill} in
\texttt{StayServiceImplTest} --- l'endpoint non aveva alcuna copertura di regressione
in precedenza, coerente con come il gap sia passato inosservato.

\paragraph{Riferimento} Commit: [pending --- da aggiornare al push]

% ─────────────────────────────────────────────────────────────────────────────
\section{Gestione Segreti}
% ─────────────────────────────────────────────────────────────────────────────

\subsection{Rimozione segreti in chiaro \threat{T-CFG-02}}

\paragraph{Descrizione}
Nella baseline il sistema presentava due categorie di segreti in chiaro nei file versionati:

\begin{enumerate}
  \item \textbf{JWT\_SECRET con fallback hardcoded} in
        \path{config-service/src/main/resources/config/auth-service.yml}: il valore di default
        \texttt{\seqsplit{bXktMzItYnl0ZS1zZWNyZXQta2V5LWZvci10ZXN0LWxvY2FsLWRldi0xMjM0NQ==}}
        (base64 della stringa nota \texttt{my-32-byte-secret-key-for-test-local-dev-12345})
        era committato nel repository. Se la variabile d'ambiente \texttt{JWT\_SECRET} non veniva
        impostata, il servizio avviava silenziosamente con questo segreto debole e noto.
        Un attaccante con accesso al repository --- o che ne conosce la struttura --- poteva
        forgiare JWT validi per qualsiasi utente.

  \item \textbf{POSTGRES\_PASSWORD hardcoded} a \texttt{password} nel \texttt{docker-compose.yml}
        e in tutti e 7 i blocchi \texttt{SPRING\_DATASOURCE\_PASSWORD}: la password del database
        era in chiaro in un file versionato, identica per tutti gli ambienti, senza possibilità
        di rotazione senza modifica del codice sorgente.
\end{enumerate}

Classificato OWASP A04:2025 --- Cryptographic Failures.

\begin{insecurebox}[title={Configurazione Vulnerabile --- auth-service.yml (baseline)}]
\begin{lstlisting}[style=yaml]
# INSICURO: fallback hardcoded — se JWT_SECRET non e' impostata,
# viene usato un segreto debole e noto committato nel repository
jwt:
  secret: "${JWT_SECRET:bXktMzItYnl0ZS1zZWNyZXQta2V5LWZvci10ZXN0LWxvY2FsLWRldi0xMjM0NQ==}"
\end{lstlisting}
\end{insecurebox}

\begin{insecurebox}[title={docker-compose.yml (baseline) --- segreti DB in chiaro}]
\begin{lstlisting}[style=yaml]
postgres:
  environment:
    POSTGRES_PASSWORD: password   # hardcoded, identico in tutti gli ambienti

auth-service:
  environment:
    - SPRING_DATASOURCE_PASSWORD=password   # idem su tutti e 7 i servizi
\end{lstlisting}
\end{insecurebox}

\paragraph{Mitigazione applicata}
L'intervento si articola in tre modifiche coordinate:

\begin{enumerate}
  \item \textbf{Rimozione del fallback} da \texttt{auth-service.yml}: il valore di default
        è stato eliminato. La proprietà diventa \texttt{"\$\{JWT\_SECRET\}"} senza alternativa:
        se la variabile d'ambiente non è presente, Spring Boot fallisce all'avvio con un errore
        esplicito (\emph{fail-fast}), impedendo che il servizio parta con un segreto insicuro.

  \item \textbf{Parametrizzazione di \texttt{POSTGRES\_PASSWORD}} nel \texttt{docker-compose.yml}:
        la password del database è ora \texttt{"\$\{POSTGRES\_PASSWORD\}"} sia per il container
        PostgreSQL sia per tutti e 7 i servizi applicativi. Il valore effettivo viene caricato
        esclusivamente dal file \texttt{.env} (gitignored).

  \item \textbf{Aggiornamento degli script di setup}
        (\texttt{setup\-hmac\-secret.sh} e \texttt{setup\-hmac\-secret.ps1}):
        entrambi gli script generano ora, in modo idempotente,
        anche \texttt{JWT\_SECRET} (48 byte random in Base64)
        e \texttt{POSTGRES\_PASSWORD} (16 byte random in esadecimale)
        nel file \texttt{.env}, oltre al già gestito
        \texttt{INTERNAL\_HMAC\_SECRET}.
\end{enumerate}

\begin{securebox}[title={Configurazione Sicura --- auth-service.yml (feature/secure-coding-hardening)}]
\begin{lstlisting}[style=yaml]
# SICURO: nessun fallback — il servizio fallisce all'avvio
# se JWT_SECRET non e' impostata nell'ambiente (fail-fast)
jwt:
  secret: "${JWT_SECRET}"
\end{lstlisting}
\end{securebox}

\begin{securebox}[title={docker-compose.yml (post-hardening) --- segreti da variabili d'ambiente}]
\begin{lstlisting}[style=yaml]
postgres:
  environment:
    POSTGRES_PASSWORD: ${POSTGRES_PASSWORD}   # da .env (gitignored)

auth-service:
  environment:
    - SPRING_DATASOURCE_PASSWORD=${POSTGRES_PASSWORD}
    - JWT_SECRET=${JWT_SECRET}                # iniettato nel container, non nel repo
\end{lstlisting}
\end{securebox}

\begin{securebox}[title={setup-hmac-secret.sh (estratto) --- generazione idempotente di tutti i segreti}]
\begin{lstlisting}[style=bash]
# Genera JWT_SECRET se non gia' presente in .env
if ! grep -q "^JWT_SECRET=" "$ENV_FILE"; then
    JWT_SECRET="$(openssl rand -base64 48 | tr -d '\n')"
    printf "JWT_SECRET=%s\n" "$JWT_SECRET" >> "$ENV_FILE"
fi

# Genera POSTGRES_PASSWORD se non gia' presente in .env
if ! grep -q "^POSTGRES_PASSWORD=" "$ENV_FILE"; then
    POSTGRES_PASSWORD="$(openssl rand -hex 16)"
    printf "POSTGRES_PASSWORD=%s\n" "$POSTGRES_PASSWORD" >> "$ENV_FILE"
fi
\end{lstlisting}
\end{securebox}

Il file \texttt{.env} rimane gitignored e viene generato localmente da ogni sviluppatore
tramite lo script di setup, oppure iniettato dal sistema di CI/CD come secret di ambiente.
Nessun segreto transita più attraverso il repository.

\paragraph{Riferimento} Commit: \commitref{91d9484}{fix(security): T-CFG-02 --- remove JWT\_SECRET fallback, parameterize POSTGRES\_PASSWORD in docker-compose}

\paragraph{Gap residuo — fallback in JwtUtil dell'API Gateway}
L'intervento descritto sopra rimuoveva il fallback soltanto da \texttt{auth-service.yml}.
Un secondo segreto debole era rimasto nascosto in
\path{api-gateway/src/main/java/com/hotelpms/gateway/util/JwtUtil.java}: l'annotazione
\texttt{@Value} conteneva lo stesso default Base64 hardcoded, rendendo di fatto inefficace
la protezione se l'API Gateway veniva avviato senza la variabile d'ambiente.

\begin{insecurebox}[title={JwtUtil.java (api-gateway) --- fallback hardcoded residuo}]
\begin{lstlisting}[style=java]
// INSICURO: il gateway accetta JWT firmati con il segreto debole
// se JWT_SECRET non e' presente nell'ambiente
@Value("${jwt.secret:bXktMzItYnl0ZS1zZWNyZXQta2V5LWZvci10ZXN0LWxvY2FsLWRldi0xMjM0NQ==}")
private String secret;
\end{lstlisting}
\end{insecurebox}

\begin{insecurebox}[title={api-gateway.yml (baseline) --- proprieta' jwt.secret assente}]
\begin{lstlisting}[style=yaml]
# jwt.secret non configurato: il gateway legge il fallback dalla
# annotazione @Value nel codice sorgente (segreto debole e noto)
internal:
  hmac:
    secret: ${INTERNAL_HMAC_SECRET}
\end{lstlisting}
\end{insecurebox}

\paragraph{Seconda mitigazione applicata}
\begin{enumerate}
  \item \textbf{Rimozione del default dall'annotazione} \texttt{@Value}: il parametro
        \texttt{jwt.secret} non ha più valore alternativo. L'assenza della variabile d'ambiente
        causa un fallimento esplicito all'avvio del gateway (\emph{fail-fast}), allineando il
        comportamento a quello già in vigore per \texttt{auth-service}.

  \item \textbf{Aggiunta di \texttt{jwt.secret} in \texttt{api-gateway.yml}}: la chiave viene
        ora dichiarata esplicitamente nel file di configurazione centralizzata, con il binding
        obbligatorio \texttt{\$\{JWT\_SECRET\}}. In questo modo il valore viene risolto
        dall'ambiente (Docker secret o file \texttt{.env}) in modo coerente con tutti gli altri
        segreti del sistema.
\end{enumerate}

\begin{securebox}[title={JwtUtil.java (post-hardening) --- fail-fast senza fallback}]
\begin{lstlisting}[style=java]
// SICURO: nessun default — il container non si avvia se JWT_SECRET
// non e' definita; nessun segreto debole puo' essere usato silenziosamente
@Value("${jwt.secret}")
private String secret;
\end{lstlisting}
\end{securebox}

\begin{securebox}[title={api-gateway.yml (post-hardening) --- jwt.secret esplicito da env}]
\begin{lstlisting}[style=yaml]
jwt:
  secret: "${JWT_SECRET}"    # obbligatorio: fail-fast se assente

internal:
  hmac:
    secret: ${INTERNAL_HMAC_SECRET}
\end{lstlisting}
\end{securebox}

\paragraph{Riferimento} Commit: \commitref{6dc6ea4}{fix(security): T-CFG-02 --- rimuovi fallback JWT\_SECRET hardcoded da api-gateway JwtUtil}

% ─────────────────────────────────────────────────────────────────────────────
% ═════════════════════════════════════════════════════════════════════════════
\section{Analisi SAST con CodeQL e Supply Chain Security \textnormal{\small\textcolor{warnyellow}{[Appendice A]}}}
% ═════════════════════════════════════════════════════════════════════════════

\begin{notebox}[title={Nota --- A03:2025 Software Supply Chain Failures + OWASP SAMM Verification (AI-assisted)}]
Questa sezione copre \textbf{A03:2025 --- Software Supply Chain Failures}, categoria nuova
dell'edizione 2025 (nel 2021 era solo ``Vulnerable and Outdated Components''): l'analisi
SAST con GitHub CodeQL, la gestione delle CVE con Trivy e l'hardening CI/CD rientrano
esattamente in questa categoria. Corrisponde alla fase \textbf{Verification} di OWASP SAMM.
Implementata con supporto AI; la critica metodologica è nell'Appendice~B.
\end{notebox}

\section{GAP-2 --- Test E2E Attack-Path (Zero Coverage di Sicurezza)}
% ─────────────────────────────────────────────────────────────────────────────

\paragraph{Descrizione}
I test E2E \textbf{Playwright} esistenti coprono esclusivamente gli
\emph{happy path} funzionali (login con successo, creazione prenotazione,
check-in). Non esiste alcun test che verifichi il comportamento del sistema
quando riceve risposte \emph{di sicurezza} dal backend: 401 Unauthorized, 403
Forbidden, 404 Not Found per IDOR, 429 Too Many Requests per brute-force.
Questa assenza lascia scoperti i meccanismi di hardening implementati nelle
sezioni precedenti: un bug di regressione nel frontend (ad esempio, un modal
che si chiude silenziosamente su 403) non sarebbe rilevato da alcuna suite
automatizzata.

\paragraph{Categoria OWASP}
\begin{itemize}
  \item \textbf{A01:2025} — Broken Access Control (IDOR, RBAC, accesso non autenticato)
  \item \textbf{A07:2025} — Authentication Failures (JWT scaduto, brute-force 429)
  \item \textbf{A01:2025} — CSRF (403 su mutation)
\end{itemize}

\begin{insecurebox}[title={Baseline --- Suite E2E con zero copertura di percorsi di attacco}]
\begin{lstlisting}[style=typescript]
// e2e/auth.spec.ts  (baseline)
// Solo happy path: login OK, redirect a dashboard, form visibile.
// Nessun test per:
//   - JWT assente/scaduto (401) -> redirect a /login
//   - Refresh fallito -> performLogout()
//   - 429 Too Many Requests -> errore generico (no detail account)
//   - 403 CSRF -> toast error, modal aperto (no silent success)

// e2e/booking-checkin.spec.ts  (baseline)
// Solo happy path: prenotazione -> check-in.
// Nessun test per:
//   - IDOR cross-hotel: PUT /guests/{id} -> 404
//   - URL manipulation: /reservations/edit/{cross-hotel-uuid} -> 404
//   - RBAC: RECEPTIONIST -> /owner-dashboard -> redirect a /
\end{lstlisting}
\end{insecurebox}

\begin{securebox}[title={Codice Sicuro --- 11 test E2E attack-path (3 suite)}]
\begin{lstlisting}[style=typescript]
// ── security-auth.spec.ts ─────────────────────────────────────────────────

// T-AUTH-SEC-02: sessione scaduta durante la navigazione
test('expired session (401 on data + failed refresh) -> /login', async ({page}) => {
  // /me: call #1 -> 200, call #2+ -> 401
  await page.route('**/api/v1/auth/me', async (route) => { /* counter */ });
  // /refresh -> 401 (refresh token revocato o scaduto)
  await page.route('**/api/v1/auth/refresh', route =>
    route.fulfill({ status: 401, body: '{"error":"Unauthorized"}' }));
  // GET /guests -> 401: scatena l'interceptor -> refresh -> logout
  await page.route('**/api/v1/guests', route =>
    route.fulfill({ status: 401, body: '{"error":"Unauthorized"}' }));
  await page.goto('/guests');
  // performLogout() -> window.location.href = '/login'
  await expect(page).toHaveURL(/\/login/, { timeout: 10_000 });
});

// T-AUTH-SEC-03: brute-force rate limit (429) -> errore generico, no detail
test('429 -> generic error, no account detail exposed', async ({ page }) => {
  await page.route('**/api/v1/auth/login', route =>
    route.fulfill({ status: 429, body: '{"detail":"Too Many Requests"}' }));
  await page.goto('/login');
  await page.getByRole('textbox', { name: /username/i }).fill('admin');
  await page.getByLabel(/password/i).fill('wrongpassword');
  await page.locator('[data-testid="login-submit"]').click();
  const errorParagraph =
    page.locator('[data-testid="login-form"] .bg-error-container p');
  await expect(errorParagraph).toBeVisible({ timeout: 5_000 });
  const msg = await errorParagraph.textContent();
  // Nessun dettaglio su lock, ban, tentativi rimasti
  expect(msg).not.toMatch(/lock|ban|remain|attempt/i);
  // Submit ri-abilitato: no UI-level lockout permanente
  await expect(page.locator('[data-testid="login-submit"]')).toBeEnabled();
});

// T-AUTH-SEC-04: CSRF rejection (403) -> toast error, no silent success
test('403 CSRF rejection -> toast error, modal stays open', async ({ page }) => {
  await page.route('**/api/v1/guests', async (route) => {
    if (route.request().method() === 'GET')
      await route.fulfill({ status: 200, body: '{"content":[],...}' });
    else // POST senza X-CSRF-Token valido
      await route.fulfill({ status: 403, body: '{"detail":"Forbidden"}' });
  });
  await page.goto('/guests');
  await page.getByRole('button', { name: /add guest/i }).click();
  await expect(page.locator('[role="dialog"]')).toBeVisible();
  await page.locator('input[name="firstName"]').fill('Test');
  await page.locator('input[name="lastName"]').fill('Guest');
  await page.locator('input[name="email"]').fill('test@example.com');
  await page.locator('button[form="guest-form"]').click();
  // Toast error (role=alert da ToastContainer)
  await expect(page.getByRole('alert')).toBeVisible({ timeout: 5_000 });
  // Modal rimane aperto: nessun successo silenzioso
  await expect(page.locator('[role="dialog"]')).toBeVisible();
});

// ── security-rbac.spec.ts ─────────────────────────────────────────────────

// T-RBAC-SEC-01: RECEPTIONIST -> /owner-dashboard -> ProtectedRoute -> /
test('RECEPTIONIST -> /owner-dashboard -> redirected to /', async ({ page }) => {
  await page.route('**/api/v1/auth/me', route =>
    route.fulfill({ status: 200,
      body: JSON.stringify({ username: 'r', role: 'RECEPTIONIST' }) }));
  await page.goto('/owner-dashboard');
  await expect(page).toHaveURL(/\/$/, { timeout: 10_000 });
  await expect(
    page.getByRole('heading', { name: /owner dashboard/i })
  ).not.toBeVisible();
});

// ── security-idor.spec.ts ─────────────────────────────────────────────────

// T-IDOR-SEC-02: URL manipulation su prenotazione cross-hotel -> 404 -> error
test('cross-hotel reservation edit via URL -> error state', async ({ page }) => {
  const id = 'res-hotel-b-99999';
  await page.route(`**/api/v1/reservations/${id}`, route =>
    route.fulfill({ status: 404,
      body: '{"detail":"RESERVATION_NOT_FOUND"}' }));
  await page.goto(`/reservations/edit/${id}`);
  // loadInitialData() cattura il 404 e chiama setError()
  await expect(
    page.locator('.bg-error-container').first()
  ).toBeVisible({ timeout: 10_000 });
  // Nessun dato di prenotazione popolato
  const dateInputs = page.locator('input[type="date"]');
  if (await dateInputs.count() > 0)
    await expect(dateInputs.first()).toHaveValue('');
});
\end{lstlisting}
\end{securebox}

\paragraph{Copertura ottenuta}
Le tre suite aggiungono \textbf{11 test E2E} che coprono i seguenti vettori:

\begin{center}
\small
\begin{tabularx}{\textwidth}{l l X l}
\toprule
\textbf{Suite} & \textbf{Test ID} & \textbf{Vettore di attacco} & \textbf{Assert chiave} \\
\midrule
security-auth  & T-AUTH-SEC-01 & Accesso non autenticato          & URL $=$ \texttt{/login} \\
security-auth  & T-AUTH-SEC-02 & Sessione scaduta + refresh fail  & performLogout $\to$ \texttt{/login} \\
security-auth  & T-AUTH-SEC-03 & Brute-force (429)                & Errore generico, no detail \\
security-auth  & T-AUTH-SEC-04 & CSRF (403 su POST)               & Toast error, modal aperto \\
security-rbac  & T-RBAC-SEC-01 & RECEPTIONIST $\to$ route ADMIN   & Redirect a \texttt{/} \\
security-rbac  & T-RBAC-SEC-02 & ADMIN $\to$ route ADMIN/OWNER    & Accesso concesso \\
security-rbac  & T-RBAC-SEC-03 & OWNER $\to$ route ADMIN/OWNER    & Accesso concesso \\
security-rbac  & T-RBAC-SEC-04 & Non autenticato $\to$ route prot.& Redirect a \texttt{/login} \\
security-idor  & T-IDOR-SEC-01 & Edit ospite cross-hotel (404)    & Toast error, modal aperto \\
security-idor  & T-IDOR-SEC-02 & URL manipulation risorsa (404)   & Error state, no dati \\
security-idor  & T-IDOR-SEC-03 & Lista ospiti hotel-scoped        & Solo ospiti hotel corrente \\
\bottomrule
\end{tabularx}
\end{center}

\paragraph{Riferimento} Commit:\\
\commitref{f2cbc7d}{fix: GAP-2 — 11 E2E security attack-path tests (auth/RBAC/IDOR)}

% ─────────────────────────────────────────────────────────────────────────────
\section{GAP-3 --- fb-service InternalAuthFilter Struttura Inconsistente}
% ─────────────────────────────────────────────────────────────────────────────

\paragraph{Descrizione}
Il filtro \texttt{InternalAuthFilter} del \textbf{fb-service} presentava una
struttura del \emph{presence check} divergente rispetto agli altri cinque
microservizi (guest, billing, reservation, stay, inventory).
Mentre il pattern canonico adottato da tutti gli altri servizi verifica la
presenza dei quattro header obbligatori in un \textbf{singolo} \texttt{if}
combinato, il fb-service utilizzava \textbf{due} blocchi \texttt{if} separati:
il primo controllava \texttt{X-Auth-User}, \texttt{X-Auth-Role} e
\texttt{X-Internal-Signature} (escludendo \texttt{X-Auth-Hotel}); il secondo
controllava \texttt{X-Auth-Hotel} in modo isolato, con un \texttt{LOG.warn}
dedicato che rivelava il \texttt{username} già estratto prima della verifica
d'integrità HMAC.

\noindent
L'impatto diretto sulla sicurezza è contenuto — la firma HMAC include
\texttt{hotelId} nel payload firmato (\texttt{username:role:hotelId}),
quindi un header mancante avrebbe comunque causato il fallimento del controllo
HMAC. Tuttavia la struttura inconsistente:
\begin{itemize}
  \item Aumentava la \emph{superficie di manutenzione}: una futura modifica al
        pattern di un solo servizio sarebbe passata inosservata.
  \item Violava il principio di \textbf{uniformità difensiva}: filtri con
        struttura diversa rendono il code review meno affidabile
        (difficile identificare deviazioni dal pattern atteso).
  \item Il \texttt{LOG.warn} nel secondo \texttt{if} emetteva il \texttt{username}
        in chiaro anche prima della validazione della firma, benché in un contesto
        interno non pubblico.
\end{itemize}

\paragraph{Classificazione OWASP}
\textbf{A01 --- Broken Access Control} (pattern di controllo accessi
non uniforme tra i servizi interni).

\vspace{0.5em}
\begin{insecurebox}[title={Codice vulnerabile --- fb-service InternalAuthFilter (presenza separata)}]
\begin{lstlisting}[language=Java]
// fb-service (PRIMA): due if separati -- hotelId escluso dal primo check
if (!StringUtils.hasText(username) || !StringUtils.hasText(role)
        || !StringUtils.hasText(signature)) {           // hotelId assente!
    rejectRequest(response, "Missing required gateway authentication headers");
    return;
}

if (!StringUtils.hasText(hotelId)) {
    // LOG emette username PRIMA della verifica HMAC
    LOG.warn("[InternalAuthFilter] Missing X-Auth-Hotel header for user={}", username);
    rejectRequest(response, "Missing X-Auth-Hotel header");
    return;
}
\end{lstlisting}
\end{insecurebox}

\vspace{0.5em}
\begin{securebox}[title={Codice sicuro --- fb-service InternalAuthFilter (singolo if canonico)}]
\begin{lstlisting}[language=Java]
// fb-service (DOPO): singolo if combinato -- identico agli altri 5 servizi
if (!StringUtils.hasText(username) || !StringUtils.hasText(role)
        || !StringUtils.hasText(hotelId) || !StringUtils.hasText(signature)) {
    rejectRequest(response, "Missing required gateway authentication headers");
    return;
}

// LOG.warn rimane solo per mismatch HMAC (dopo che tutti gli header sono presenti)
if (!isSignatureValid(username, role, hotelId, signature)) {
    LOG.warn("[InternalAuthFilter] HMAC signature mismatch for user={}", username);
    rejectRequest(response, "Invalid internal request signature");
    return;
}
\end{lstlisting}
\end{securebox}

\paragraph{Test di copertura}\mbox{}\\
Sono stati aggiunti 8 test JUnit~5 in\\
\path{fb-service/.../security/InternalAuthFilterTest.java}
che coprono tutti i path del filtro:

\begin{center}
\small
\begin{tabular}{l p{5.5cm} p{5.5cm}}
\toprule
\textbf{Nested class} & \textbf{Test} & \textbf{Assert chiave} \\
\midrule
PresenceCheck  & \seqsplit{shouldRejectWhenUsernameMissing}    & HTTP 401 \\
PresenceCheck  & shouldRejectWhenRoleMissing                    & HTTP 401 \\
PresenceCheck  & \seqsplit{shouldRejectWhenHotelIdMissing}      & HTTP 401 \\
PresenceCheck  & \seqsplit{shouldRejectWhenSignatureMissing}    & HTTP 401 \\
HmacCheck      & \seqsplit{shouldRejectWhenSignatureInvalid}    & HTTP 401 \\
HmacCheck      & \seqsplit{shouldRejectWhenSignatureComputedForDifferentHotelId} & HTTP 401 (tenant isolation) \\
HmacCheck      & \seqsplit{shouldAuthenticateWithValidHeaders}  & HTTP 200 + SecurityContext (name/details/authority) \\
ActuatorBypass & \seqsplit{shouldBypassFilterForActuatorEndpoint} & HTTP 200 senza headers \\
\bottomrule
\end{tabular}
\end{center}

\paragraph{Riferimento} Commit:
\commitref{9c299b1}{fix(security): GAP-3 --- fb-service InternalAuthFilter single combined presence check}

% ─────────────────────────────────────────────────────────────────────────────
\section{GAP-4 --- Audit Log Non Aggregati (SIEM Assente)}
% ─────────────────────────────────────────────────────────────────────────────

\paragraph{Descrizione}
I log strutturati emessi dai microservizi con prefissi \texttt{[AUTH]},
\texttt{[STAY]} e \texttt{[BILLING]} venivano scritti esclusivamente su
\texttt{stdout} del container Docker, senza alcun meccanismo di aggregazione,
ricerca o alerting centralizzato.
In assenza di un sistema di raccolta, risultava impossibile rispondere a
domande operative di sicurezza quali:
\begin{itemize}
  \item \emph{«Quanti \texttt{LOGIN\_FAILED} ha prodotto l'utente X nelle
        ultime 24 ore?»}
  \item \emph{«Ci sono stati picchi anomali di \texttt{ACCOUNT\_LOCKED} in
        una finestra temporale ristretta?»}
  \item \emph{«Quali richieste hanno fallito la verifica HMAC nelle ultime
        ore?»}
\end{itemize}

\noindent
La mancanza di aggregazione log costituisce una violazione diretta di
\textbf{OWASP A09 --- Security Logging and Monitoring Failures}: anche se i
singoli eventi vengono emessi correttamente, la loro utilità è nulla se non
possono essere consultati, correlati e alertati in modo centralizzato.

\paragraph{Classificazione OWASP}
\textbf{A09 --- Security Logging and Monitoring Failures} (log prodotti ma
non aggregati, nessuna capacità di rilevamento anomalie o risposta a incidenti).

\paragraph{Soluzione implementata}
La soluzione è composta da tre livelli cooperanti:

\begin{enumerate}
  \item \textbf{Infrastruttura}: aggiunta di \textbf{Loki}
        (\texttt{grafana/loki:2.9.0}) come aggregatore di log e
        \textbf{Grafana} (\texttt{grafana/grafana:10.2.0}) come layer di
        visualizzazione e alerting, entrambi nella rete Docker isolata
        \texttt{observability-network} con volumi persistenti.
  \item \textbf{Appender Logback}: integrazione della libreria
        \texttt{loki-logback-appender:1.5.2} nei tre servizi con audit log
        strutturati (auth-service, stay-service, billing-service).
        Ogni servizio riceve un file \texttt{logback-spring.xml} dedicato.
  \item \textbf{Dashboard Grafana}: provisioning automatico di una dashboard
        \emph{Security Audit Log} con sei pannelli operativi: contatori
        \texttt{LOGIN\_FAILED}, \texttt{ACCOUNT\_LOCKED} e \texttt{HMAC
        Failures} nell'ultima ora, grafico di tendenza degli eventi
        di sicurezza al minuto, e visualizzazione raw dei log per
        ciascun prefisso di audit.
\end{enumerate}

\vspace{0.5em}
\begin{insecurebox}[title={Situazione precedente --- log su stdout, nessuna aggregazione}]
\begin{lstlisting}[language=Java, caption={auth-service: audit event emesso su stdout senza destinazione SIEM}]
// AuthServiceImpl.java (PRIMA) -- @Slf4j, log solo su stdout container
LOG.warn("[AUTH] LOGIN_FAILED user={} reason=BAD_PASSWORD attempt={}/{}",
         username, failedAttempts, lockoutThreshold);
// Il log scompare nel buffer del container; nessun sistema lo raccoglie.
// Impossibile aggregare, cercare, correlare o alertare su questi eventi.
\end{lstlisting}
\end{insecurebox}

\vspace{0.5em}
\begin{securebox}[title={Livello 1 --- Loki e Grafana in docker-compose (observability-network)}]
\begin{lstlisting}[language=bash, caption={docker-compose.yml: servizi Loki e Grafana aggiunti}]
loki:
  image: grafana/loki:2.9.0
  command: -config.file=/etc/loki/config.yaml
  volumes:
    - ./docker/loki/loki-config.yaml:/etc/loki/config.yaml
    - loki_data:/loki
  networks: [observability-network]

grafana:
  image: grafana/grafana:10.2.0
  environment:
    - GF_SECURITY_ADMIN_PASSWORD=admin
    - GF_USERS_ALLOW_SIGN_UP=false
  volumes:
    - ./docker/grafana/provisioning:/etc/grafana/provisioning
    - grafana_data:/var/lib/grafana
  depends_on:
    loki: { condition: service_healthy }
  networks: [observability-network]
\end{lstlisting}
\end{securebox}

\vspace{0.5em}
\begin{securebox}[title={Livello 2 --- logback-spring.xml con LokiAppender (auth-service)}]
\begin{lstlisting}[language=XML, caption={logback-spring.xml: Loki attivo solo nel profilo Docker; test usano solo console}]
<configuration>
  <include resource=
    "org/springframework/boot/logging/logback/defaults.xml"/>
  <springProperty scope="context" name="appName"
    source="spring.application.name" defaultValue="auth-service"/>

  <appender name="CONSOLE"
            class="ch.qos.logback.core.ConsoleAppender">
    <encoder>
      <pattern>${CONSOLE_LOG_PATTERN}</pattern>
    </encoder>
  </appender>

  <!-- Attivo solo con SPRING_PROFILES_ACTIVE=auth-service (Docker) -->
  <springProfile name="auth-service">
    <appender name="LOKI"
              class="com.github.loki4j.logback.Loki4jAppender">
      <http>
        <url>${LOKI_URL:-http://loki:3100}/loki/api/v1/push</url>
      </http>
      <format>
        <label>
          <pattern>service=${appName},level=%level</pattern>
        </label>
        <message>
          <pattern>[%level] [%thread] %logger{36} | %msg%n%ex</pattern>
        </message>
        <sortByTime>true</sortByTime>
      </format>
    </appender>

    <appender name="LOKI_ASYNC"
              class="ch.qos.logback.classic.AsyncAppender">
      <appender-ref ref="LOKI"/>
      <queueSize>512</queueSize>
      <discardingThreshold>0</discardingThreshold>
      <!-- neverBlock: Loki irraggiungibile non blocca mai i thread app -->
      <neverBlock>true</neverBlock>
    </appender>

    <root level="INFO">
      <appender-ref ref="CONSOLE"/>
      <appender-ref ref="LOKI_ASYNC"/>
    </root>
  </springProfile>

  <!-- Profilo default (test / sviluppo locale): solo console -->
  <springProfile name="!auth-service">
    <root level="INFO">
      <appender-ref ref="CONSOLE"/>
    </root>
  </springProfile>
</configuration>
\end{lstlisting}
\end{securebox}

\vspace{0.5em}
\begin{securebox}[title={Livello 3 --- Query LogQL nella dashboard Grafana}]
\begin{lstlisting}[language=bash, caption={LogQL: pannelli dashboard "Security Audit Log"}]
# LOGIN_FAILED nell'ultima ora (pannello stat con soglia rossa a 20)
sum(count_over_time(
  {service="auth-service"} |= "LOGIN_FAILED" [1h]
))

# ACCOUNT_LOCKED nell'ultima ora (soglia arancio a 1)
sum(count_over_time(
  {service="auth-service"} |= "ACCOUNT_LOCKED" [1h]
))

# HMAC Failures su qualsiasi servizio (soglia rossa a 1)
sum(count_over_time(
  {service=~".+", level="WARN"} |= "HMAC" [1h]
))

# Trend eventi/min — grafico time-series
sum by (service) (rate(
  {service=~"auth-service|stay-service|billing-service"}
  |= "LOGIN_FAILED|CHECK_IN_FAILED|CHECK_OUT_FAILED" [5m]
))
\end{lstlisting}
\end{securebox}

\paragraph{Scelte progettuali rilevanti}

\begin{itemize}
  \item \textbf{Profilo Spring come guard}: il \texttt{<springProfile>} in
        \texttt{logback-spring.xml} assicura che l'appender Loki sia attivo
        esclusivamente quando il servizio viene avviato in Docker
        (\texttt{SPRING\_PROFILES\_ACTIVE=auth-service}). In ambiente di test
        JUnit non è presente alcun profilo attivo: solo il \texttt{ConsoleAppender}
        è attivo, nessuna connessione a Loki viene tentata.
  \item \textbf{AsyncAppender con \texttt{neverBlock=true}}: se Loki è
        irraggiungibile (es. container non ancora avviato), i log vengono
        scartati silenziosamente dalla coda senza mai bloccare i thread
        dell'applicazione. La resilienza funzionale dell'applicazione non
        dipende dalla disponibilità del sistema di logging.
  \item \textbf{Label Loki \texttt{service/level}}: il pattern di label\\
        \texttt{service=\$\{appName\},level=\%level} permette query LogQL
        efficienti con filtro per servizio e livello di log, abilitando
        l'aggregazione cross-service dei log di audit.
  \item \textbf{Provisioning automatico Grafana}: datasource Loki e dashboard
        vengono caricati automaticamente all'avvio tramite i file YAML in
        \path{docker/grafana/provisioning/}, eliminando ogni configurazione
        manuale post-deploy.
\end{itemize}

\paragraph{Riferimento} Commit:
\commitref{fe99e2b}{feat(security): GAP-4 --- Loki log aggregation SIEM for audit events}

% ═════════════════════════════════════════════════════════════════════════════
\section{GAP-5 \textbar{} T-GST-05 --- GDPR Data Retention: Assenza di Policy sul Ciclo di Vita dei PII Ospite}
% ─────────────────────────────────────────────────────────────────────────────

\paragraph{Descrizione}
Il \texttt{guest-service} persiste in modo permanente i dati personali degli
ospiti — nome, cognome, data di nascita, indirizzo, email, numero di telefono
e, attraverso l'entità \texttt{IdentityDocument}, tipo documento, numero,
date di emissione e scadenza, paese di rilascio.
Queste informazioni vengono create al momento della prima prenotazione e
rimangono nel database a tempo indeterminato: non esiste alcun meccanismo
automatico di cancellazione, anonimizzazione o scadenza.

L'unica operazione di rimozione attualmente disponibile è il \emph{soft-delete}
(\texttt{active = false} via \texttt{@SQLDelete}), che nasconde il record dalle
query applicative ma lascia i dati fisicamente presenti nel database.
Questo comportamento è incompatibile con l'articolo~17 del GDPR
(\emph{diritto alla cancellazione}): un interessato che eserciti tale diritto
non ottiene la cancellazione effettiva dei propri dati, bensì la loro semplice
invisibilità a livello applicativo.

Inoltre non esiste alcun campo di consenso sull'entità \texttt{Guest}
(\texttt{gdprConsentDate}, \texttt{consentPurpose}), rendendo impossibile
dimostrare la base giuridica del trattamento in caso di ispezione da parte
del Garante Privacy.

\paragraph{Classificazione OWASP}
\textbf{A04 --- Insecure Design}: la mancanza di una policy di ritenzione è
un difetto di progettazione del modello dei dati, non una vulnerabilità
puntuale. Il rischio non dipende da un bug nel codice ma dall'assenza di un
controllo progettato fin dall'inizio.

\paragraph{Contesto normativo italiano}
La situazione è resa più complessa dalla coesistenza di obblighi di legge
che \emph{impongono} la conservazione per periodi minimi e del diritto GDPR
che ne chiede la cancellazione su richiesta:

\begin{center}
\small
\begin{tabularx}{\textwidth}{X l l l}
\toprule
\textbf{Obbligo} & \textbf{Dato} & \textbf{Durata minima} & \textbf{Prevalenza} \\
\midrule
Circ.\ PS / D.Lgs.\ 196/2003 & Dati alloggiati & 5 anni & Su richiesta GDPR \\
Codice Civile art.\ 2220 & Fatture / pagamenti & 10 anni & Su richiesta GDPR \\
GDPR art.\ 17 par.\ 3 lett.\ b & — & Deroga per obblighi legali & Prevale sulla richiesta \\
\bottomrule
\end{tabularx}
\end{center}

\noindent
In pratica: la richiesta di cancellazione di un ospite che abbia soggiornato
negli ultimi 5 anni deve essere \emph{rifiutata motivandola per legge}, non
eseguita. Trascorso il periodo obbligatorio, i dati devono invece essere
cancellati o anonimizzati.

\paragraph{Confronto con PMS professionali}
I sistemi PMS commerciali (Opera, Mews, Cloudbeds, Protel) adottano
uniformemente il seguente modello:

\begin{enumerate}
  \item \textbf{Retention period configurabile per hotel}: l'operatore imposta
        il periodo di conservazione (tipicamente 5--7 anni dall'ultimo
        soggiorno). Un job notturno identifica i profili scaduti e li
        \emph{anonimizza}: il nome viene sostituito con un identificativo
        pseudonimo, tutti i campi PII vengono azzerati e i documenti di
        identità vengono cancellati fisicamente. Lo storico contabile e gli
        stay associati vengono mantenuti per la rendicontazione fiscale.
  \item \textbf{Hard-delete con guardia legale}: l'endpoint di cancellazione
        su richiesta dell'interessato verifica preventivamente l'esistenza di
        obblighi legali prevalenti (fatture nel periodo fiscale, dati alloggiati
        nel periodo PS). Se il blocco esiste, la richiesta viene rifiutata con
        un messaggio che indica la data a partire dalla quale la cancellazione
        sarà possibile. Altrimenti, PII e documenti vengono cancellati
        fisicamente e il profilo viene anonimizzato.
  \item \textbf{Separazione fisica PII / dati operativi}: i dati di identità
        (entità \texttt{IdentityDocument}) sono mantenuti in una tabella
        separata dal profilo operativo (preferenze, storico soggiorni),
        consentendo retention differenziata: i dati di identità scadono prima,
        il profilo operativo (senza PII) può essere conservato più a lungo per
        analisi storiche e CRM.
\end{enumerate}

\noindent
L'architettura attuale del sistema presenta già una base favorevole:
\texttt{IdentityDocument} è un'entità separata da \texttt{Guest} con
\texttt{CascadeType.ALL}, il che rende tecnicamente semplice implementare una
retention differenziata.

\paragraph{Soluzione implementata}

I quattro step del piano sono stati realizzati integralmente nel
commit \texttt{56109eb}.

\begin{enumerate}
  \item \textbf{Campo consenso --- Flyway V4}: la migration
        \texttt{V4\_\_add\_gdpr\_consent\_date.sql} aggiunge la colonna
        \texttt{gdpr\_consent\_date DATE NOT NULL} alla tabella \texttt{guests},
        con backfill automatico da \texttt{created\_at} per i record esistenti
        (proxy legalmente difendibile del momento di registrazione in hotel).
        \texttt{GuestServiceImpl.createGuest()} stampa \texttt{LocalDate.now()}
        su ogni nuovo profilo.

  \item \textbf{Retention job notturno}:
        \texttt{GuestRetentionJobServiceImpl} è annotato
        \texttt{@Scheduled(cron = "0 0 2 * * *")} e gira ogni notte alle 02:00.
        Usa il pre-filtro conservativo \texttt{FISCAL\_MIN\_YEARS = 10} anni
        per identificare i candidati senza Feign call: solo i profili la cui
        \texttt{gdprConsentDate} risale a più di 10 anni fa entrano nella
        pipeline. Per ciascun candidato verifica le due hold indipendenti tramite
        \texttt{StayServiceClient} e \texttt{BillingServiceClient}; in caso di
        indisponibilità del servizio il circuit-breaker ritorna
        \texttt{hasStays = true} (fail-closed: nessun ospite viene mai
        anonimizzato per errore). L'anonimizzazione imposta
        \texttt{firstName = "GDPR"}, \texttt{lastName = "ERASED\_<id-prefix>"},
        azzera tutti i campi PII, svuota \texttt{identityDocuments}
        (cascade hard-delete) e pone \texttt{active = false}.

  \item \textbf{Hard-delete con guardia legale duale}:
        \texttt{GuestServiceImpl.deleteGuest()} esegue due check in sequenza:
        \begin{itemize}
          \item \textbf{TULPS} --- chiama \texttt{GET /api/v1/stays/guest/\{id\}/last-stay-date}
                (stay-service) e verifica che \texttt{lastStayDate + max(retentionYears, 5)} sia
                già trascorso.
          \item \textbf{Fiscale} --- chiama \texttt{GET /api/v1/invoices/guest/\{id\}/last-invoice-date}
                (billing-service) e verifica che \texttt{lastInvoiceDate + 10 anni} sia
                già trascorso.
        \end{itemize}
        Se almeno un vincolo è attivo, l'operazione è rifiutata con
        \texttt{HTTP 451 Unavailable For Legal Reasons} (RFC 7725) e un corpo
        \texttt{GdprLegalHoldResponse} contenente il campo \texttt{unlocksAt}
        (la data di sblocco più lontana) e \texttt{legalBasis}
        (\texttt{TULPS}, \texttt{FISCAL} o \texttt{TULPS\_AND\_FISCAL}).
        Il \texttt{GlobalExceptionHandler} intercetta
        \texttt{GdprLegalHoldException} e produce la risposta 451.

  \item \textbf{Configurabilità per hotel --- Flyway V5}:
        la migration \path{V5__add_guest_privacy_settings.sql} crea la
        tabella \texttt{\seqsplit{guest\_privacy\_settings}} (una riga per hotel).\\
        \texttt{\seqsplit{GuestPrivacySettingsController}} espone
        \texttt{GET/PUT /api/v1/guests/settings} con Jakarta Bean Validation
        (\texttt{@Min(5)}) che impedisce di scendere sotto il minimo TULPS.
        Il periodo di retention configurato viene rispettato sia dalla guardia
        legale hard-delete sia dal retention job.
\end{enumerate}

\paragraph{Fix collaterale: FeignHeaderConfig mancante}
La \texttt{FeignHeaderConfig} del guest-service non calcolava la firma
HMAC-SHA256 (\texttt{X-Internal-Signature}), a differenza di tutti gli altri
servizi del progetto. Di conseguenza ogni chiamata Feign in uscita dal
guest-service veniva rifiutata con HTTP 401 dagli \texttt{InternalAuthFilter}
dei servizi downstream, e il circuit-breaker si attivava silenziosamente.
Il file è stato allineato al pattern canonico (identico a stay-service)
con l'aggiunta della dipendenza da \texttt{\$\{internal.hmac.secret\}} e del
calcolo HMAC \texttt{username:role:hotelId} (T-GW-07).

\paragraph{Test}
Sono stati scritti 17 test suddivisi in tre classi, con copertura
dei boundary value imposti dalla legge italiana:

\begin{itemize}
  \item \texttt{GuestServiceImplDeleteGuardTest} (7 test): verifica i casi limite
        temporali — soggiorno a 4 anni e 364 giorni (bloccato), esattamente
        5 anni (bloccato), 5 anni e 1 giorno (liberato); fattura a esattamente
        10 anni (bloccata), 10 anni e 1 giorno (liberata); entrambi i vincoli
        attivi con \texttt{unlocksAt} corretto; retention configurabile
        per hotel a 7 anni.
  \item \texttt{GuestPrivacySettingsServiceImplTest} (5 test): get-or-create,
        default creation, update, update-when-none-exists, no-save-on-existing.
  \item \texttt{GuestRetentionJobServiceImplTest} (5 test): anonimizzazione
        completa, skip per TULPS hold, skip per fiscal hold, ospite senza
        storico, nessun candidato.
\end{itemize}

\paragraph{Stato}
\textbf{\textcolor{fixgreen}{RISOLTO}} --- commit \texttt{56109eb},
fix warning \texttt{066f611}.
Tutti e quattro gli step del piano di mitigazione sono implementati;
la copertura del test include i boundary value temporali critici richiesti
dalla normativa italiana.

% ═════════════════════════════════════════════════════════════════════════════
\chapter{Gestione Autenticazione e Autorizzazione}
% ═════════════════════════════════════════════════════════════════════════════

Questo capitolo presenta una visione sintetica e trasversale del sistema
di identità e controllo degli accessi risultante dall'intero ciclo di hardening.
I dettagli implementativi di ciascun intervento sono documentati nel
Capitolo~3; qui l'obiettivo è offrire una lettura architetturale complessiva
che mostri come i singoli interventi si raccordino in un modello coerente
di difesa in profondità.

% ─────────────────────────────────────────────────────────────────────────────
\section{Ridisegno del flusso di autenticazione}
% ─────────────────────────────────────────────────────────────────────────────

Il flusso di autenticazione nella baseline era un unico JWT con scadenza di
24 ore, privo di revoca e senza binding al tenant dell'utente.
Il ridisegno, completato in cinque iterazioni successive (commit
\commitref{a90e462}{T-AUTH-02}, \commitref{4e1f1e9}{T-AUTH-03},
\commitref{446ad8d}{T-AUTH-03 upgrade Argon2id},
\commitref{5ce3010}{T-AUTH-04}, \commitref{b469e21}{T-GW-06},
\commitref{b473265}{T-AUTH-04 residuo}), produce il seguente flusso
di autenticazione:

\begin{enumerate}
  \item \textbf{Login} (\texttt{POST /api/v1/auth/login}):
    verifica lockout → verifica password (Argon2id via DelegatingPasswordEncoder, memory-hard) → genera
    \emph{access token} (TTL 15 min, claim \texttt{role}, \texttt{hotelId},
    \texttt{tv}) + \emph{refresh token} (TTL 7 giorni, \texttt{jti} UUID,
    \texttt{typ=refresh}) → emette tre cookie:
    \texttt{jwt}, \texttt{refresh\_token} (entrambi httpOnly/Secure/SameSite=Strict)
    e \texttt{csrf\_token} (non-httpOnly, per il Double-Submit-Cookie).
    Scrive la chiave Redis \texttt{user:tv:<username>} con il token version corrente.

  \item \textbf{Richiesta autenticata}:
    il browser invia automaticamente il cookie JWT.
    L'\texttt{AuthenticationFilter} del gateway valida la firma HMAC-SHA256
    del JWT, estrae \texttt{username}, \texttt{role}, \texttt{hotelId},
    inietta \texttt{X-Auth-User}, \texttt{X-Auth-Role}, \texttt{X-Auth-Hotel}
    e \texttt{X-Internal-Signature} sulla richiesta verso il microservizio.
    Il \texttt{CsrfFilter} (ordine \texttt{HIGHEST\_PRECEDENCE+2}) verifica
    che l'header \texttt{X-CSRF-Token} corrisponda al cookie \texttt{csrf\_token}
    per ogni metodo mutante.

  \item \textbf{Refresh} (\texttt{POST /api/v1/auth/refresh}):
    valida il refresh token (firma, scadenza, tipo, JTI non blacklistato,
    token version Redis) → blacklista il JTI consumato → emette nuovo
    token pair con token version aggiornata.

  \item \textbf{Cambio password} (\texttt{POST /api/v1/auth/change-password}):
    verifica la password corrente → incrementa \texttt{tokenVersion} in DB
    → sovrascrive Redis \texttt{user:tv:<username>} → emette nuovo token pair
    (sessione corrente mantenuta) → ogni refresh token precedente viene rifiutato
    alla prossima rotazione (TV mismatch → 401).

  \item \textbf{Logout} (\texttt{POST /api/v1/auth/logout}):
    blacklista il JTI del refresh token → azzera i tre cookie
    (\texttt{jwt}, \texttt{refresh\_token}, \texttt{csrf\_token}).
    Il log \texttt{[AUTH] LOGOUT | user=...} garantisce tracciabilità.

  \item \textbf{Reset password admin} (\texttt{PATCH /api/v1/auth/users/\{id\}/reset-password}):
    ADMIN/OWNER invia nuova password → Argon2id encode →
    \texttt{mustChangePassword=true} → incrementa \texttt{tokenVersion}
    → sovrascrive Redis \texttt{user:tv:<username>} → 204 No Content.
    Tutte le sessioni attive dell'utente vengono invalidate (TV mismatch
    al prossimo refresh). Al login successivo il sistema reindirizza a
    \texttt{/profile} per impostare una nuova password personale.
    Commit: \commitref{4e6a7f4}{G3 --- reset password con invalidazione sessioni (ADMIN/OWNER)}.
\end{enumerate}

\paragraph{Fix inizializzazione hash admin — Flyway V6.}
L'adozione di \texttt{DelegatingPasswordEncoder} richiede che ogni hash in database
sia prefissato con il nome dell'algoritmo (es.\ \texttt{\{bcrypt\}}) per permettere
il dispatch corretto al verificatore.
Il seed iniziale dell'utente \texttt{admin} in Flyway era privo di tale prefisso,
causando un crash \texttt{IllegalArgumentException} al primo login.
La migrazione \texttt{V6\_\_fix\_admin\_password\_bcrypt\_prefix.sql} corregge
l'hash con una guardia idempotente
(\texttt{WHERE NOT password LIKE '\textbackslash\{bcrypt\textbackslash\}\%'})
che risulta no-op su installazioni già corrette.
Commit: \commitref{953be80}{fix(security): V6 Flyway --- prefisso bcrypt mancante sull'hash admin iniziale}.

\paragraph{Limitazione operativa — recovery accesso amministrativo.}
In assenza di un secondo account ADMIN attivo, il ripristino dell'accesso
amministrativo richiede accesso diretto al database PostgreSQL del container
\texttt{auth-service} (es.\ \texttt{UPDATE user\_account SET active = true WHERE ...}
e reset dell'hash). Non esiste un endpoint applicativo di recovery:
questa è una limitazione operativa nota, da documentare nel runbook di deployment.

% ─────────────────────────────────────────────────────────────────────────────
\section{Modello RBAC}
% ─────────────────────────────────────────────────────────────────────────────

Il sistema adotta un modello \textbf{Role-Based Access Control} a tre ruoli,
con enforcement a tre livelli:

\begin{center}
\small
\begin{tabular}{llp{7cm}}
\toprule
\textbf{Ruolo} & \textbf{Accesso} & \textbf{Perimetro operativo} \\
\midrule
\texttt{ADMIN}       & Totale & Tutte le funzionalità, tutte le risorse dell'hotel \\
\texttt{OWNER}       & Totale & Identico ad ADMIN + dashboard reportistica \\
\texttt{RECEPTIONIST}& Limitato & Ospiti, prenotazioni, check-in/out, F\&B;
                                  nessun accesso a fatture e dashboard finanziaria \\
\bottomrule
\end{tabular}
\end{center}

\paragraph{Enforcement a tre livelli}

\begin{enumerate}
  \item \textbf{Gateway (stateless)}: l'\texttt{AuthenticationFilter} valida il JWT
        e inietta \texttt{X-Auth-Role}; le route del gateway non filtrano per ruolo
        (delegato al layer applicativo), ma rifiutano qualsiasi richiesta priva di JWT
        valido con HTTP 401.

  \item \textbf{Microservizi (Spring Security)}: ciascun microservizio
        popola il \texttt{SecurityContext} tramite \texttt{InternalAuthFilter}
        con un \texttt{\seqsplit{SimpleGrantedAuthority("ROLE\_" + role)}} derivato dall'header
        \texttt{X-Auth-Role}. I metodi di servizio e i controller usano
        \texttt{@PreAuthorize} o \texttt{hasRole()} di Spring Security per
        verificare il ruolo sul thread corrente. Un audit post-sviluppo ha
        identificato due endpoint (\texttt{POST /stays/reports/alloggiati/submit}
        e \texttt{PUT /stays/settings}) privi di \texttt{@PreAuthorize}
        lato server — l'accesso era limitato solo via UI.
        Entrambi sono stati corretti aggiungendo
        \texttt{@PreAuthorize("hasAnyRole('ADMIN','OWNER')")}
        con \texttt{@EnableMethodSecurity} abilitato su \texttt{SecurityConfig}.
        \texttt{AccessDeniedException} è gestito dal \texttt{GlobalExceptionHandler}
        che restituisce HTTP 403 anziché 500; il comportamento è verificato da
        \texttt{StayControllerSecurityTest} con test \texttt{@WebMvcTest} e HMAC reale.
        Commit: \commitref{52e869c}{fix(security): RBAC --- @PreAuthorize su StayController (submit/download alloggiati)}
        e \commitref{44edf35}{fix(security): RBAC --- @PreAuthorize su HotelSettingsController.updateSettings}.

  \item \textbf{Frontend (React)}: il componente \texttt{ProtectedRoute}
        con prop \texttt{allowedRoles} redirige a \texttt{/} gli utenti
        non autorizzati prima di renderizzare la pagina (T-FE-03). Questo
        terzo livello è una misura \emph{UX}; la garanzia di sicurezza
        effettiva è fornita dai livelli 1 e 2.
\end{enumerate}

% ─────────────────────────────────────────────────────────────────────────────
\section{Gestione del ciclo di vita del JWT}
% ─────────────────────────────────────────────────────────────────────────────

Il JWT in produzione porta le seguenti claim obbligatorie:

\begin{center}
\small
\begin{tabular}{llp{6.5cm}}
\toprule
\textbf{Claim} & \textbf{Tipo} & \textbf{Scopo} \\
\midrule
\texttt{sub}     & \texttt{String} & Username dell'utente autenticato \\
\texttt{role}    & \texttt{String} & Ruolo (\texttt{ADMIN}, \texttt{OWNER}, \texttt{RECEPTIONIST}) \\
\texttt{hotelId} & \texttt{String} (UUID) & Tenant binding: iniettato come \texttt{X-Auth-Hotel} \\
\texttt{tv}      & \texttt{int}    & Token version: invalidazione bulk al cambio password \\
\texttt{iat}     & \texttt{long}   & Issued at (standard JWT) \\
\texttt{exp}     & \texttt{long}   & Expiration: 15 min (access), 7 giorni (refresh) \\
\texttt{jti}     & \texttt{String} (UUID) & Solo refresh token: identificatore per blacklist Redis \\
\texttt{typ}     & \texttt{String} & Solo refresh token: valore fisso \texttt{"refresh"} \\
\bottomrule
\end{tabular}
\end{center}

\paragraph{Stato della revoca}
Il sistema gestisce due categorie di revoca:

\begin{itemize}
  \item \textbf{Individuale} (logout, furto rilevato): il JTI del refresh token
        viene scritto in Redis con chiave \texttt{rt:blacklist:<jti>} e TTL
        pari alla vita residua del token. Scadenza automatica, nessun job di
        manutenzione.
  \item \textbf{Bulk} (cambio password): la chiave Redis
        \texttt{user:tv:<username>} viene sovrascritta con il nuovo valore
        di \texttt{tokenVersion}. Qualsiasi refresh token con \texttt{tv} inferiore
        viene rifiutato in fase di rotation, senza enumerare i JTI.
\end{itemize}

\paragraph{Firma HMAC inter-service}
Oltre alla firma JWT (HS256), ogni richiesta interna è autenticata da una firma
HMAC-SHA256 calcolata su \texttt{"username:role:hotelId"} e trasmessa nell'header
\texttt{X-Internal-Signature}. L'\texttt{InternalAuthFilter} di ogni microservizio
verifica la firma con confronto \texttt{MessageDigest.isEqual} (constant-time)
prima di popolare il \texttt{SecurityContext}, rendendo impossibile sia
l'impersonation sia il tampering del tenant context (T-GW-01, T-GW-07).

% ─────────────────────────────────────────────────────────────────────────────
\section{Protezioni contro attacchi comuni all'autenticazione}
% ─────────────────────────────────────────────────────────────────────────────

La tabella seguente sintetizza le protezioni attive contro i principali vettori
di attacco sull'autenticazione, con riferimento alla categoria OWASP e al commit
di implementazione.

\begin{center}
\small
\begin{tabular}{p{3.2cm} p{4.5cm} p{3.0cm} p{2.8cm}}
\toprule
\textbf{Vettore} & \textbf{Protezione} & \textbf{OWASP} & \textbf{Commit} \\
\midrule
Brute force credenziali
  & Lockout 5 tentativi / 15 min; rate limit gateway 5~req/s per IP
  & A07:2025 & \texttt{a90e462} \\
User enumeration
  & Messaggio uniforme \texttt{INVALID\_CREDENTIALS}
  & A07:2025 & baseline \\
Password debole in DB
  & Argon2id memory-hard (DelegatingPasswordEncoder); BCrypt legacy migrati on login
  & A04:2025 & \texttt{4e1f1e9}, \texttt{446ad8d} \\
Token theft (access)
  & TTL 15 min; cookie httpOnly/Secure/SameSite
  & A04:2025/\linebreak[2]A07:2025 & \texttt{5ce3010} \\
Token theft (refresh)
  & Rotation + blacklist JTI Redis; bulk revoca TV
  & A07:2025 & \texttt{5ce3010}, \texttt{b473265} \\
Session hijacking CSRF
  & Double-Submit Cookie (header \texttt{X-CSRF-Token} == cookie)
  & A01:2025 & \texttt{e146b17} \\
JWT forging
  & HS256 con segreto 48 byte random; fail-fast senza \texttt{JWT\_SECRET}
  & A04:2025 & \texttt{91d9484}, \texttt{6dc6ea4} \\
Header injection inter-service
  & HMAC-SHA256 su \texttt{username:role:hotelId} (constant-time compare)
  & A01:2025 & baseline, \texttt{5bdc594} \\
Tenant hopping (IDOR)
  & \texttt{hotelId} nel JWT; X-Auth-Hotel iniettato; query hotel-scoped; 404 IDOR-safe
  & A01 & \texttt{b469e21} + \path{T-GST/RES/BILL/FB} \\
XSS → token exfiltration
  & httpOnly cookie; ESLint ban \texttt{dangerously\-Set\-Inner\-HTML}; CSP nginx
  & A05:2025 & \texttt{3655906}, \texttt{b901a9f} \\
\bottomrule
\end{tabular}
\end{center}

\paragraph{Postura complessiva}
Al termine del ciclo di hardening, il sistema implementa tutti i controlli
minimi raccomandati dall'OWASP Authentication Cheat Sheet e dalla NIST SP 800-63B
per applicazioni aziendali: password hashing con funzione lenta, lockout
progressivo, revoca dei token, binding crittografico del tenant context, e
audit logging di ogni evento di autenticazione verso un aggregatore centralizzato
(Loki + Grafana). La superficie di attacco residua nella \emph{nostra codebase}
è ridotta a un singolo alert aperto nel Security tab di GitHub:
\textbf{CVE-2026-42577} (Netty native epoll, rischio pratico nullo perché
Spring Boot non attiva il trasporto native epoll in configurazione JDK NIO).
Gli aggiornamenti a \textbf{Spring Boot 3.5.x} e \textbf{Spring Cloud 2025.0.x}
hanno risolto CVE-2026-22739 (config-server), CVE-2025-41253 (gateway EL
injection) e le CVE Netty della serie 4.1.13x. I 108~alert Trivy rilevati
sulle immagini di terze parti (Grafana, postgres) sono stati analizzati,
triaggiati e dismissi con motivazione tecnica documentata (si veda DEP-CVE-04).

% ═════════════════════════════════════════════════════════════════════════════
% APPENDICI
% ═════════════════════════════════════════════════════════════════════════════
\appendix

% ═════════════════════════════════════════════════════════════════════════════
\chapter{Appendice A --- Implementazioni Avanzate (AI-assisted)}
\label{app:enterprise}
% ═════════════════════════════════════════════════════════════════════════════

Questa appendice raccoglie le implementazioni di sicurezza sviluppate con supporto
diretto di \textbf{Claude Code} (Anthropic) che vanno oltre i requisiti minimi
della traccia d'esame. Il codice è funzionante e integrato nel sistema; è documentato
qui per completezza architetturale, non rivendicato come competenza autonoma esclusiva.

% ─────────────────────────────────────────────────────────────────────────────
\section{Autenticazione avanzata}
% ─────────────────────────────────────────────────────────────────────────────

\paragraph{T-AUTH-04 residual --- Token versioning e revoca bulk (Redis)}
Al cambio password, il sistema incrementa \texttt{tokenVersion} in PostgreSQL e
aggiorna la chiave Redis \texttt{user:tv:<username>}. Qualsiasi refresh token con
\texttt{tv} inferiore al valore corrente è rifiutato alla prossima rotazione,
invalidando tutte le sessioni precedenti. La finestra di esposizione residua è
limitata ai soli access token in volo (massimo 15 minuti).
Commit: \commitref{b473265}{fix(security): GAP-1 --- password change revokes sessions via tokenVersion}

\paragraph{T-AUTH-05 --- Audit log strutturati (OWASP A09)}
Ogni evento di autenticazione (login/fallimento/lockout/logout) produce una riga
strutturata con prefisso \texttt{[AUTH]}, livello INFO/WARN, username e motivazione.
Corrispondenza con OWASP A09:2025 --- Security Logging \& Alerting Failures.
Commit: \commitref{9a3d7f2}{fix(security): T-AUTH-05 --- structured auth audit log}

% ─────────────────────────────────────────────────────────────────────────────
\section{Sicurezza inter-service avanzata}
% ─────────────────────────────────────────────────────────────────────────────

\paragraph{T-GW-07 --- HMAC-SHA256 inter-service con binding sul tenant}
Header \texttt{X-Internal-Signature} = HMAC-SHA256 di \texttt{"user:role:hotelId"}.
Verificato con \texttt{MessageDigest.isEqual} (constant-time) da ogni microservizio.
Impedisce impersonation e tampering del tenant context.
Commit: \commitref{b469e21}{fix(security): T-GW-07 --- HMAC inter-service tenant binding}

\paragraph{T-GW-08 --- Anti-replay (timestamp + nonce) su HMAC inter-service}
Il payload firmato T-GW-07 (\texttt{"user:role:hotelId"}) non includeva alcun elemento
temporale: una firma catturata restava una credenziale valida indefinitamente, senza
possibilità di revoca se non ruotando il segreto condiviso. Esteso a
\texttt{"user:role:hotelId:timestamp:nonce"}: l'API Gateway genera timestamp
(epoch millis) e nonce (UUID) casuali ad ogni richiesta, propagati come
\texttt{X-Auth-Timestamp}/\texttt{X-Auth-Nonce}. Ogni \texttt{InternalAuthFilter} verifica
una finestra di tolleranza di 60s sul timestamp e reclama il nonce su un
\texttt{NonceStore} backed da Redis (\texttt{SETNX}, TTL 120s) — un secondo reclamo dello
stesso nonce entro la finestra viene rifiutato come replay.
Commit: \commitref{5dd8ed8}{feat(security): T-GW-08 --- timestamp+nonce anti-replay on internal HMAC auth}

\paragraph{T-GW-06 --- Propagazione X-Auth-Hotel end-to-end}
Aggiunto campo \texttt{hotelId} a \texttt{UserAccount} e claim JWT corrispondente.
L'API Gateway inietta \texttt{X-Auth-Hotel} su ogni richiesta a valle.
Senza questo intervento tutti i servizi multi-tenant restituivano HTTP~401.
Commit: \commitref{b469e21}{fix(security): T-GW-06 --- X-Auth-Hotel propagation end-to-end}

\paragraph{T-GW-02 --- Rate limiting Redis token bucket proxy-aware}
Key resolver aggiornato per leggere \texttt{X-Forwarded-For} anziché l'IP TCP
(dietro un proxy tutti i client condividerebbero un unico bucket).
Rate limiting esteso a tutte le route business, non solo \texttt{/auth/**}.
Commit: \commitref{3f8a1c2}{fix(security): T-GW-02 --- proxy-aware rate limiting all routes}

% ─────────────────────────────────────────────────────────────────────────────
\section{Configurazione e SDLC avanzati}
% ─────────────────────────────────────────────────────────────────────────────

\paragraph{T-CFG-03 --- Autenticazione HTTP Basic sul config-service}
Il config-service (porta 8888) espone JWT\_SECRET, POSTGRES\_PASSWORD e
INTERNAL\_HMAC\_SECRET. HTTP Basic Auth aggiunta per impedire lettura anonima.
Commit: \commitref{c934f6e}{fix(security): T-CFG-03 --- HTTP Basic Auth on config-service}

\paragraph{SAST/CodeQL e supply chain CI/CD (OWASP SAMM Verification)}
10 alert CodeQL risolti: permessi minimi \texttt{contents: read} su ogni job CI
(OWASP A02:2025 --- Security Misconfiguration), 9 falsi positivi CSRF su microservizi interni chiusi con motivazione.
Corrisponde alla fase \textbf{Verification} di OWASP SAMM e al principio shift-left.
Commit: \commitref{92c0050}{fix(ci): add permissions: contents: read to all CI jobs}

\paragraph{CVE scanning con Trivy (DEP-CVE-01/02/04)}
Pipeline CI esegue Trivy su ogni immagine Docker.
CVE risolte: \texttt{CVE-2026-33870/33871} (Netty request smuggling + HTTP/2 DoS),
\texttt{CVE-2025-41235/41253} (Spring Cloud Gateway header forwarding + EL injection).
108 alert su immagini terze parti triaggiati con motivazione documentata.

\paragraph{Aggiornamento dipendenze frontend (DEP-CVE-05)}
L'analisi del security tab GitHub ha rilevato 7~CVE sulle dipendenze npm del frontend.
Risolte con upgrade mirati senza breaking change (OWASP~A06:2025):
\begin{itemize}
  \item \texttt{react-router-dom} + \texttt{react-router} \texttt{7.13.1}~→~\texttt{7.17.0}:
    \texttt{\#36} XSS via RSC redirect, \texttt{\#37} XSS via \texttt{Location} header,
    \texttt{\#38} RCE via turbo-stream, \texttt{\#39} open redirect,
    \texttt{\#40}/\texttt{\#41} DoS via \texttt{\_\_manifest} e single-fetch.
  \item \texttt{esbuild} \texttt{0.27.3}~→~\texttt{0.28.1} (transitivo di Vite,
    forzato via \texttt{npm overrides}): \texttt{CVE-2025-42577} arbitrary file
    read tramite dev server su Windows.
\end{itemize}
Tutti i 344 test rimasti verdi; build e lint zero warning.
Commit: \commitref{b1be3e8}{fix(deps): resolve 7 Dependabot CVEs --- react-router 7.17.0, esbuild 0.28.1}

% ═════════════════════════════════════════════════════════════════════════════
\chapter{Appendice B --- Utilizzo dell'IA e Revisione Critica}
\label{app:ia}
% ═════════════════════════════════════════════════════════════════════════════

\section{Metodologia}

Il progetto è stato sviluppato a partire da una codebase personale preesistente
(branch \texttt{pre-secure-coding}), usata come baseline insicura.
\textbf{Claude Code} (Anthropic, modello \texttt{claude-sonnet-4-6}) è stato impiegato
come strumento di analisi e supporto all'implementazione, \emph{non} come sostituto
del ragionamento critico dello studente.

In questa appendice sono documentati i momenti in cui l'output dell'AI è stato valutato
criticamente, evidenziando i casi in cui il suggerimento era errato, incompleto o
incompatibile con il contesto specifico del sistema.

% ─────────────────────────────────────────────────────────────────────────────
\section{Episodio 1 --- JWT storage: localStorage suggerito dall'AI}
% ─────────────────────────────────────────────────────────────────────────────

\paragraph{Prompt utilizzato}
\begin{quote}
``Implementa un sistema di autenticazione JWT con Spring Boot e React.
Fornisci il codice per generare il token nel backend e per archiviarlo nel frontend.''
\end{quote}

\paragraph{Output dell'AI}
L'AI ha generato un'implementazione funzionante dove il frontend archiviava
il token con \texttt{localStorage.setItem('token', response.data.token)}
e lo inviava come header \texttt{Authorization: Bearer <token>} su ogni richiesta Axios.

\paragraph{Critica allo studente}
\texttt{localStorage} è accessibile a qualsiasi script JavaScript della pagina.
Un attacco XSS --- anche tramite una dipendenza npm compromessa --- può leggere
il token con \texttt{localStorage.getItem('token')} e inviarlo a un server
dell'attaccante. L'AI ha prodotto codice funzionale ma insicuro per OWASP A07
(Identification and Authentication Failures). Il problema non è nell'algoritmo
JWT ma nel vettore di esfiltrazione abilitato dallo storage scelto.

\paragraph{Correzione applicata (T-FE-02 / T-AUTH-04)}
Il token è stato spostato in un cookie \texttt{httpOnly} impostato dal server:

\begin{securebox}[title={Cookie httpOnly al posto di localStorage (T-FE-02)}]
\begin{lstlisting}[style=java]
// Il browser non puo' leggere il cookie con JavaScript (difesa XSS)
private ResponseCookie createCookie(String value, int maxAge) {
    return ResponseCookie.from("jwt", value)
        .httpOnly(true)      // blocca lettura da JavaScript
        .secure(true)        // trasmesso solo su HTTPS
        .sameSite("Strict")  // blocca invio cross-site
        .path("/")
        .maxAge(maxAge)
        .build();
}
\end{lstlisting}
\end{securebox}

% ─────────────────────────────────────────────────────────────────────────────
\section{Episodio 2 --- BCrypt: cost factor obsoleto per il 2025/26}
% ─────────────────────────────────────────────────────────────────────────────

\paragraph{Prompt utilizzato}
\begin{quote}
``Sostituisci la memorizzazione della password in chiaro con un hashing sicuro.
Usa BCrypt con Spring Security.''
\end{quote}

\paragraph{Output dell'AI}
L'AI ha generato \texttt{new BCryptPasswordEncoder()} senza parametri, che usa
il cost factor di default di Spring Security: \textbf{10}.
Il codice era strutturalmente corretto ma usava il parametro di default.

\paragraph{Critica allo studente}
Il cost factor 10 era adeguato per hardware del 2011. Le raccomandazioni OWASP
al 2025/26 indicano cost~$\geq$~12 come minimo su hardware moderno, per garantire
almeno 300 ms di calcolo su singolo core. Con cost~10 su hardware attuale il
calcolo richiede $\sim$50 ms, rendendo il brute force $2^2 = 4\times$ più rapido
del necessario (SEI CERT MSC62-J).

\paragraph{Correzione applicata (T-AUTH-03)}

\begin{securebox}[title={BCrypt cost 12 esplicito (T-AUTH-03, SEI CERT MSC62-J)}]
\begin{lstlisting}[style=java]
// Cost 12: OWASP raccomanda >= 12 per hardware 2025.
// Cost 10 (default Spring Security) era ottimale per hardware 2011.
@Bean
public PasswordEncoder passwordEncoder() {
    return new BCryptPasswordEncoder(12);
}
\end{lstlisting}
\end{securebox}

% ─────────────────────────────────────────────────────────────────────────────
\section{Episodio 3 --- IDOR: l'AI ha ignorato il contesto multi-tenant}
% ─────────────────────────────────────────────────────────────────────────────

\paragraph{Prompt utilizzato}
\begin{quote}
``Nel guest-service, un utente può accedere ai dati di ospiti che non appartengono
alla sua struttura. Correggi il problema IDOR.''
\end{quote}

\paragraph{Output dell'AI}
L'AI ha suggerito di aggiungere un controllo basato sull'\texttt{userId} del
chiamante, verificando \texttt{guest.createdBy == currentUserId}:
\begin{lstlisting}[style=java]
// Suggerimento AI -- ERRATO nel contesto multi-hotel
Optional<Guest> findByIdAndCreatedBy(UUID guestId, UUID userId);
\end{lstlisting}

\paragraph{Critica allo studente}
Il suggerimento era errato nel contesto di un PMS multi-hotel.
Le risorse appartengono a un \emph{hotel}, non a un singolo utente.
Con il modello proposto dall'AI: un ADMIN non potrebbe accedere agli ospiti
creati da un RECEPTIONIST dello stesso hotel; un RECEPTIONIST trasferito
manterrebbe accesso agli ospiti del precedente. La dimensione di isolamento
corretta è \texttt{hotel\_id} estratto dal JWT, non \texttt{userId}.

\paragraph{Correzione applicata (T-GST-01 / T-GST-03)}

\begin{securebox}[title={Isolamento per hotel\_id (T-GST-01, OWASP A01)}]
\begin{lstlisting}[style=java]
// Isolamento per hotel, non per utente
Optional<Guest> findByIdAndActiveTrueAndHotelId(UUID id, UUID hotelId);

// hotelId estratto dal SecurityContext, mai dal client (IDOR-safe)
private UUID resolveHotelId() {
    Authentication auth =
        SecurityContextHolder.getContext().getAuthentication();
    return UUID.fromString((String) auth.getDetails());
}
\end{lstlisting}
\end{securebox}

% ─────────────────────────────────────────────────────────────────────────────
\section{Episodio 4 --- CSRF: pattern incompatibile con SPA React}
% ─────────────────────────────────────────────────────────────────────────────

\paragraph{Prompt utilizzato}
\begin{quote}
``Implementa la protezione CSRF per il gateway API. L'applicazione usa cookie
httpOnly per il JWT e Spring Cloud Gateway (WebFlux).''
\end{quote}

\paragraph{Output dell'AI}
L'AI ha suggerito il filtro CSRF standard di Spring Security con
\texttt{CookieCsrfTokenRepository.withHttpOnlyFalse()}, che attende il token
in un campo \texttt{\_csrf} del form HTML o nell'header \texttt{X-CSRF-TOKEN}.

\paragraph{Critica allo studente}
Il filtro standard Spring Security è progettato per applicazioni MVC che rendono
form HTML. Non è compatibile con una SPA React che usa chiamate Axios JSON perché:
(1) il gateway è Spring Cloud Gateway (WebFlux/reactive), non Spring MVC;
(2) il token CSRF viene esposto solo su endpoint specifici con richiesta asincrona
aggiuntiva, incompatibile con l'architettura SPA senza rendering server-side.

\paragraph{Correzione applicata (T-GW-05)}
Implementato il pattern \textbf{Double Submit Cookie} con un filtro personalizzato:
il cookie \texttt{csrf\_token} (non-httpOnly, leggibile dalla SPA) viene confrontato
con l'header \texttt{X-CSRF-Token} (aggiunto da Axios). Un sito malevolo non può
costruire l'header perché la same-origin policy impedisce di leggere il cookie.

% ─────────────────────────────────────────────────────────────────────────────
\section{Episodio 5 --- User enumeration: messaggi ``utili per il debug''}
% ─────────────────────────────────────────────────────────────────────────────

\paragraph{Prompt utilizzato}
\begin{quote}
``Scrivi il metodo login() per AuthServiceImpl. Distingui i casi
errore per facilitare il debug: utente non trovato e password errata.''
\end{quote}

\paragraph{Output dell'AI}
\begin{lstlisting}[style=java]
// Output AI -- INSICURO: messaggi distinti
.orElseThrow(() ->
    new BadCredentialsException("USER_NOT_FOUND"));  // diverso!

if (!passwordEncoder.matches(...)) {
    throw new BadCredentialsException("INVALID_PASSWORD");  // diverso!
}
\end{lstlisting}

\paragraph{Critica allo studente}
Messaggi distinti consentono la \emph{user enumeration}: un attaccante può
inviare richieste automatizzate e identificare username validi dalla risposta.
OWASP A07 cita esplicitamente questo pattern come vettore per brute force mirati.
\textbf{L'AI ha motivato la scelta come ``utile per il debug''}, dimostrando
che non considera il contesto di sicurezza nella formulazione dei messaggi di errore.

\paragraph{Correzione applicata (T-AUTH-01)}

\begin{securebox}[title={Messaggio uniforme --- nessuna information disclosure (T-AUTH-01)}]
\begin{lstlisting}[style=java]
// SICURO: identico per USER_NOT_FOUND e INVALID_PASSWORD
.orElseThrow(() ->
    new BadCredentialsException("INVALID_CREDENTIALS"));

if (!passwordEncoder.matches(request.password(),
                              user.getPasswordHash())) {
    throw new BadCredentialsException("INVALID_CREDENTIALS");
}
\end{lstlisting}
\end{securebox}

% ─────────────────────────────────────────────────────────────────────────────
\section{Riepilogo della revisione critica}
% ─────────────────────────────────────────────────────────────────────────────

\begin{center}
\small
\begin{tabular}{p{2.8cm} p{4.5cm} p{4.5cm}}
\toprule
\textbf{Episodio} & \textbf{Errore/limite dell'AI} & \textbf{Correzione studente} \\
\midrule
JWT storage
  & \texttt{localStorage} (vulnerabile a XSS)
  & Cookie \texttt{httpOnly} + \texttt{SameSite=Strict} \\
BCrypt cost factor
  & Default 10 (hardware 2011)
  & Cost 12 esplicito (OWASP 2025) \\
IDOR multi-tenant
  & Isolamento per \texttt{userId} (errato)
  & Isolamento per \texttt{hotel\_id} (corretto) \\
CSRF su SPA
  & Filter Spring MVC (incompatibile WebFlux)
  & Double Submit Cookie custom \\
User enumeration
  & Messaggi distinti ``per debug''
  & Messaggio uniforme \texttt{INVALID\_CREDENTIALS} \\
\bottomrule
\end{tabular}
\end{center}

In tutti e cinque i casi l'AI ha prodotto codice \emph{funzionale} ma con difetti
di sicurezza identificabili applicando i principi del corso: OWASP Top~10, SEI CERT
e la conoscenza del contesto architetturale del sistema.

\end{document}
